% =====================================================================
%  Anatomía de un LLM — Fase 0 — Módulo 0.1
%  Lógica y técnicas de demostración
%  Compilar: latexmk -r ../latexmkrc P01-demostraciones.tex
% =====================================================================
\documentclass[11pt,a4paper,twoside,openany]{book}
\usepackage{../estilo-apuntes}

\newcommand{\ApuntesCurso}{Anatomía de un LLM}
\newcommand{\ApuntesModulo}{Módulo 0.1}

% Notación propia del módulo
\newcommand{\N}{\mathbb{N}}
\newcommand{\Z}{\mathbb{Z}}
\newcommand{\Q}{\mathbb{Q}}
\renewcommand{\V}{\mathrm{V}}% valor de verdad «verdadero» (se redefine \V de la plantilla)
\newcommand{\F}{\mathrm{F}}
\newcommand{\xor}{\oplus}
\newcommand{\divides}{\mid}
\newcommand{\nodivide}{\nmid}
\newcommand{\eqlog}{\equiv}
\DeclarePairedDelimiter{\floor}{\lfloor}{\rfloor}
\DeclarePairedDelimiter{\ceil}{\lceil}{\rceil}

\hypersetup{pdftitle={Módulo 0.1. Lógica y técnicas de demostración},pdfauthor={Anatomía de un LLM}}

\begin{document}
\frontmatter
% ------------------------------------------------------------------ Portada
\begin{titlepage}
\centering
\vspace*{2.5cm}
{\large\scshape Anatomía de un LLM}\par
\vspace{0.4cm}
{\normalsize Fundamentos matemáticos de los modelos de lenguaje}\par
\vspace{2.2cm}
\rule{\textwidth}{0.8pt}\par\vspace{0.5cm}
{\small\scshape Fase 0 \textperiodcentered\ Preliminares matemáticos}\par\vspace{0.3cm}
{\Huge\bfseries Módulo 0.1\par}\vspace{0.3cm}
{\LARGE Lógica y técnicas de demostración\par}
\vspace{0.5cm}\rule{\textwidth}{0.4pt}\par
\vspace{1.6cm}
{\large Apuntes del curso}\par
\vfill
{\small Edición de octubre de 2026}\par
\end{titlepage}

% ------------------------------------------------------- Presentación
\chapter*{Presentación del módulo}
\addcontentsline{toc}{chapter}{Presentación del módulo}
\markboth{Presentación del módulo}{}

Los apuntes del curso enuncian definiciones, teoremas y proposiciones, y casi todos los resultados van acompañados de una demostración. Leer esas demostraciones, comprobarlas y escribir otras nuevas exige un lenguaje preciso y un repertorio pequeño de técnicas que se repiten una y otra vez: la demostración directa, el contrarrecíproco, la reducción al absurdo, la división en casos y la inducción. Este módulo, el primero de la fase de preliminares, presenta ese lenguaje y esas técnicas desde el nivel de bachillerato, con la misma organización que los cursos introductorios de demostración de las universidades de referencia, y ofrece una colección amplia de ejercicios cuyas soluciones están redactadas como modelo de escritura.

\section*{Objetivos}
Al finalizar el módulo, el estudiante debe ser capaz de:
\begin{enumerate}[label=\arabic*.]
  \item construir tablas de verdad, reconocer tautologías y equivalencias lógicas y distinguir una implicación de su recíproca y de su contrarrecíproca;
  \item traducir enunciados con cuantificadores al lenguaje simbólico y negarlos correctamente, incluidos los de cuantificadores anidados;
  \item demostrar enunciados condicionales y bicondicionales de forma directa y por contrarrecíproco;
  \item emplear la reducción al absurdo y la demostración por casos, y reproducir las pruebas de la irracionalidad de $\sqrt2$ y de la infinitud de los números primos;
  \item demostrar por inducción simple y fuerte, y relacionar la inducción con el principio del buen orden;
  \item demostrar existencia y unicidad, refutar enunciados mediante contraejemplos y redactar una demostración clara y completa.
\end{enumerate}

\section*{Requisitos previos}
Matemáticas de bachillerato: aritmética de los números enteros y racionales, potencias y raíces, manipulación algebraica de expresiones, ecuaciones de primer y segundo grado y la noción intuitiva de conjunto. No se presupone ningún módulo anterior.

\section*{Organización}
\begin{center}
\begin{tabular}{@{}cll@{}}
\toprule
Tema & Contenido & Fuentes principales \\
\midrule
1 & Proposiciones, conectivas y tablas de verdad & \cite{hammack}, \cite{velleman}, \cite{mcs} \\
2 & Cuantificadores y negación & \cite{hammack}, \cite{velleman}, \cite{cs103} \\
3 & Demostración directa y por contrarrecíproco & \cite{hammack}, \cite{velleman}, \cite{cs103} \\
4 & Reducción al absurdo y demostración por casos & \cite{hammack}, \cite{mcs}, \cite{cs103} \\
5 & Inducción simple y fuerte; buen orden & \cite{hammack}, \cite{mcs}, \cite{velleman} \\
6 & Existencia, unicidad, contraejemplos y redacción & \cite{hammack}, \cite{cs103}, \cite{polya} \\
\bottomrule
\end{tabular}
\end{center}

Cada tema termina con una lista de lecturas y con una colección extensa de ejercicios graduados en tres niveles: $\star$ comprensión, $\star\star$ cálculo o práctica mecánica y $\star\star\star$ demostración. Los ejercicios marcados con \emph{Fuente} están adaptados de cursos y libros de texto universitarios; el enunciado se ha redactado de nuevo y, en algunos casos, se han cambiado los datos. Las soluciones completas, paso a paso, están en el \cref{ap:soluciones}. En los ejercicios de demostración, el apartado «Desarrollo» de la solución contiene la demostración redactada tal como debería entregarse; conviene leerla como modelo de escritura, no solo como comprobación del resultado.

\section*{Dónde se usa en el curso}
\begin{center}
\small
\begin{tabularx}{\textwidth}{@{}>{\raggedright\arraybackslash}p{4.6cm}X@{}}
\toprule
Concepto & Módulos del curso que lo emplean \\
\midrule
Lectura y escritura de demostraciones & Todos los módulos con teoremas, del módulo~1 al~23. \\
Condición necesaria y suficiente, contraejemplo & M04 (criterio de la hessiana), M07 (convexidad), M12 (fusiones de BPE con solapamiento). \\
Cuantificadores anidados ($\forall\varepsilon\,\exists\delta$) & M04 (límites y derivadas), M07 (convergencia de algoritmos), M01 (límites de la softmax con temperatura). \\
Inyectividad por contrarrecíproco & M01 (tokenizador como codificación y decodificación), M12 (UTF-8 y BPE). \\
Reducción al absurdo & M04 (condiciones de mínimo con el teorema de Taylor), M06 (códigos y cadenas). \\
Inducción simple & M01 (regla de la cadena), M02 (aplicaciones lineales), M06 (desigualdad de Jensen), M07 (medias exponenciales de Adam), M10 (grafos computacionales), M12 (BPE), M13 (máscara causal), M17 (\emph{all-reduce} y FlashAttention). \\
Existencia y unicidad & M02 y M08 (soluciones de sistemas y mínimos cuadrados), M07 (mínimo de una función estrictamente convexa). \\
\bottomrule
\end{tabularx}
\end{center}

\section*{Sobre las fuentes}
El texto es de redacción propia. La organización sigue la de los capítulos~2 y 4 a~10 de \emph{Book of Proof}, de Hammack~\cite{hammack}, y la de los capítulos~1 a~3 y~6 de \emph{How to Prove It}, de Velleman~\cite{velleman}, dos libros de texto habituales en los cursos universitarios de introducción a la demostración. Las partes sobre lógica proposicional, buen orden e inducción siguen además los capítulos~1, 2, 3 y~5 del texto \emph{Mathematics for Computer Science}~\cite{mcs}, del curso 6.042J del MIT~\cite{mit6042}. Las recomendaciones de redacción proceden de la lista de comprobación del curso CS103 de Stanford~\cite{cs103}, de la §5.3 de~\cite{hammack} y de Polya~\cite{polya}.

\section*{Notación}
\begin{center}
\begin{tabular}{@{}ll@{}}
\toprule
Símbolo & Significado \\
\midrule
$\neg P$, $P\wedge Q$, $P\vee Q$ & negación, conjunción y disyunción (inclusiva) \\
$P\Rightarrow Q$, $P\Leftrightarrow Q$ & condicional («si $P$, entonces $Q$») y bicondicional \\
$\V$, $\F$ & valores de verdad «verdadero» y «falso» \\
$\Phi\eqlog\Psi$ & equivalencia lógica de dos fórmulas \\
$\forall$, $\exists$, $\exists!$ & «para todo», «existe», «existe un único» \\
$\N=\{1,2,3,\dots\}$, $\N_0=\N\cup\{0\}$ & números naturales, sin y con el cero \\
$\Z$, $\Q$, $\R$ & enteros, racionales y reales \\
$a\divides b$, $a\nodivide b$ & $a$ divide a $b$; $a$ no divide a $b$ \\
$a\equiv b\pmod n$ & $a$ y $b$ son congruentes módulo $n$ \\
$\floor{x}$, $\ceil{x}$ & parte entera por defecto y por exceso \\
$\blacksquare$ & fin de una demostración \\
\bottomrule
\end{tabular}
\end{center}
El símbolo $\equiv$ se usa con dos sentidos que el contexto distingue: entre fórmulas lógicas indica equivalencia lógica, y en la notación $a\equiv b\pmod n$ indica congruencia.

\tableofcontents

\mainmatter
% =====================================================================
\chapter{Proposiciones, conectivas y tablas de verdad}
\label{tema:1}
% =====================================================================

Una demostración es una cadena de afirmaciones enlazadas por reglas de razonamiento. Antes de estudiar las técnicas de demostración conviene fijar qué es una afirmación, cómo se combinan varias mediante conectivas y cuándo dos combinaciones dicen lo mismo. Este tema desarrolla la lógica proposicional con el enfoque del capítulo~2 de~\cite{hammack} (§2.1–2.6 y §2.11) y del capítulo~1 de~\cite{velleman}; la parte algebraica sigue la §3.4 de~\cite{mcs}.

\section{Proposiciones y conectivas}

\begin{definicion}[Proposición]\label{def:proposicion}
Una \emph{proposición} es un enunciado del que tiene sentido afirmar que es verdadero o que es falso, y que es exactamente una de las dos cosas. Su \emph{valor de verdad} es $\V$ (verdadero) o $\F$ (falso).
\end{definicion}

\begin{ejemplo}\label{ej:proposiciones}
«$7$ es un número primo» es una proposición verdadera; «$2+2=5$» es una proposición falsa; «todo número par mayor que~2 es suma de dos primos» (conjetura de Goldbach) es una proposición, aunque hoy no se sepa cuál es su valor de verdad. En cambio, «¿qué hora es?», «cierre la puerta» y «$x+1=3$» no son proposiciones: las dos primeras no afirman nada y la tercera depende del valor de $x$, que no está fijado. Enunciados como el último se llaman \emph{predicados} y se estudian en el \cref{tema:2}.
\end{ejemplo}

A partir de proposiciones $P$ y $Q$ se forman proposiciones compuestas mediante \emph{conectivas lógicas}. El valor de verdad de la proposición compuesta depende solo de los de $P$ y $Q$, y se describe con una \emph{tabla de verdad}.

\begin{definicion}[Negación, conjunción y disyunción]\label{def:conectivas}
Sean $P$ y $Q$ proposiciones.
\begin{enumerate}
  \item La \emph{negación} $\neg P$ («no $P$») es verdadera si y solo si $P$ es falsa.
  \item La \emph{conjunción} $P\wedge Q$ («$P$ y $Q$») es verdadera si y solo si $P$ y $Q$ son ambas verdaderas.
  \item La \emph{disyunción} $P\vee Q$ («$P$ o $Q$») es verdadera si y solo si al menos una de las dos es verdadera.
\end{enumerate}
\end{definicion}

\begin{table}[h]
\centering
\caption{Tablas de verdad de las conectivas básicas.}\label{tab:conectivas}
\begin{tabular}{@{}cc|ccccc@{}}
\toprule
$P$ & $Q$ & $\neg P$ & $P\wedge Q$ & $P\vee Q$ & $P\Rightarrow Q$ & $P\Leftrightarrow Q$ \\
\midrule
$\V$ & $\V$ & $\F$ & $\V$ & $\V$ & $\V$ & $\V$ \\
$\V$ & $\F$ & $\F$ & $\F$ & $\V$ & $\F$ & $\F$ \\
$\F$ & $\V$ & $\V$ & $\F$ & $\V$ & $\V$ & $\F$ \\
$\F$ & $\F$ & $\V$ & $\F$ & $\F$ & $\V$ & $\V$ \\
\bottomrule
\end{tabular}
\end{table}

La disyunción matemática es \emph{inclusiva}: $P\vee Q$ es verdadera también cuando lo son las dos. La disyunción \emph{exclusiva}, «o $P$ o $Q$, pero no ambas», se escribe $P\xor Q$ y se estudia en el \cref{eje:1-xor}.

\section{Condicional y bicondicional}

\begin{definicion}[Condicional]\label{def:condicional}
El \emph{condicional} $P\Rightarrow Q$ («si $P$, entonces $Q$») es falso únicamente cuando $P$ es verdadera y $Q$ es falsa. $P$ se llama \emph{hipótesis} o antecedente, y $Q$, \emph{tesis} o consecuente.
\end{definicion}

La tabla del condicional (\cref{tab:conectivas}) se entiende como una promesa: «si apruebas el examen, te invito a cenar» solo se incumple si se aprueba y no hay cena. Si no se aprueba, la promesa no se ha roto, haya cena o no. Por eso un condicional con hipótesis falsa es verdadero; se dice que es \emph{vacuamente} verdadero.

El condicional $P\Rightarrow Q$ admite varias lecturas equivalentes, que conviene reconocer porque aparecen constantemente en los enunciados de los teoremas~\cite[§2.3]{hammack}:
\begin{center}
\begin{tabular}{@{}l@{}}
«si $P$, entonces $Q$»; \quad «$Q$ si $P$»; \quad «$P$ solo si $Q$»; \\
«$P$ es condición suficiente para $Q$»; \quad «$Q$ es condición necesaria para $P$».
\end{tabular}
\end{center}

\begin{definicion}[Bicondicional]\label{def:bicondicional}
El \emph{bicondicional} $P\Leftrightarrow Q$ («$P$ si y solo si $Q$») es verdadero cuando $P$ y $Q$ tienen el mismo valor de verdad. Se lee también «$P$ es condición necesaria y suficiente para $Q$».
\end{definicion}

\begin{definicion}[Recíproco, contrarrecíproco e inverso]\label{def:reciproco}
Dado el condicional $P\Rightarrow Q$, se llaman \emph{recíproco} a $Q\Rightarrow P$, \emph{contrarrecíproco} a $\neg Q\Rightarrow\neg P$ e \emph{inverso} a $\neg P\Rightarrow\neg Q$.
\end{definicion}

\begin{ejemplo}\label{ej:reciproco}
Sea $n$ un entero y considérese «si $n$ es múltiplo de~4, entonces $n$ es par», que es verdadero. Su recíproco, «si $n$ es par, entonces es múltiplo de~4», es falso: $n=2$ lo refuta. Su contrarrecíproco, «si $n$ no es par, entonces no es múltiplo de~4», es verdadero. La \cref{prop:leyes}(v) muestra que un condicional y su contrarrecíproco siempre tienen el mismo valor de verdad, mientras que el recíproco puede no tenerlo.
\end{ejemplo}

\section{Fórmulas, tautologías y contradicciones}

\begin{definicion}[Fórmula proposicional]
Las \emph{fórmulas proposicionales} sobre las variables $P_1,\dots,P_n$ son las expresiones que se obtienen a partir de las variables aplicando un número finito de veces las conectivas $\neg$, $\wedge$, $\vee$, $\Rightarrow$ y $\Leftrightarrow$, con paréntesis para indicar el orden. Una \emph{asignación} da a cada variable un valor de verdad, y la tabla de verdad de la fórmula lista su valor para cada asignación.
\end{definicion}

\begin{proposicion}\label{prop:filas}
Una fórmula con $n$ variables tiene una tabla de verdad de $2^{n}$ filas.
\end{proposicion}

\begin{proof}
Una asignación es una elección independiente de $\V$ o $\F$ para cada una de las $n$ variables. Por el principio del producto, hay $2\cdot 2\cdots 2=2^{n}$ elecciones.
\end{proof}

\begin{definicion}[Tautología y contradicción]
Una fórmula es una \emph{tautología} si es verdadera para todas las asignaciones, y una \emph{contradicción} si es falsa para todas.
\end{definicion}

\begin{ejemplo}\label{ej:tercero-excluido}
$P\vee\neg P$ (principio del tercero excluido) es una tautología, y $P\wedge\neg P$ es una contradicción: en cada fila exactamente una de $P$ y $\neg P$ es verdadera. Más interesante es la fórmula $\bigl((P\Rightarrow Q)\wedge P\bigr)\Rightarrow Q$. Su única forma de ser falsa sería con $Q$ falsa y $(P\Rightarrow Q)\wedge P$ verdadera; pero entonces $P$ es verdadera y $P\Rightarrow Q$ también, lo que obliga a que $Q$ sea verdadera. Por tanto es una tautología; es la regla de \emph{modus ponens} de la \cref{sec:inferencia}.
\end{ejemplo}

El razonamiento del ejemplo anterior evita escribir la tabla completa: para comprobar que un condicional es tautología basta ver que la fila «hipótesis verdadera y tesis falsa» es imposible. Esta forma de pensar es ya una demostración por reducción al absurdo (\cref{tema:4}).

\section{Equivalencia lógica}

\begin{definicion}[Equivalencia lógica]\label{def:equivalencia}
Dos fórmulas $\Phi$ y $\Psi$ en las mismas variables son \emph{lógicamente equivalentes}, y se escribe $\Phi\eqlog\Psi$, si tienen el mismo valor de verdad en todas las asignaciones; equivalentemente, si $\Phi\Leftrightarrow\Psi$ es una tautología.
\end{definicion}

\begin{proposicion}[Leyes lógicas]\label{prop:leyes}
Para proposiciones cualesquiera $P$, $Q$ y $R$:
\begin{enumerate}
  \item doble negación: $\neg\neg P\eqlog P$;
  \item leyes de De Morgan: $\neg(P\wedge Q)\eqlog\neg P\vee\neg Q$ \ y \ $\neg(P\vee Q)\eqlog\neg P\wedge\neg Q$;
  \item distributivas: $P\wedge(Q\vee R)\eqlog(P\wedge Q)\vee(P\wedge R)$ \ y \ $P\vee(Q\wedge R)\eqlog(P\vee Q)\wedge(P\vee R)$;
  \item condicional como disyunción: $P\Rightarrow Q\eqlog\neg P\vee Q$;
  \item contrarrecíproco: $P\Rightarrow Q\eqlog\neg Q\Rightarrow\neg P$;
  \item negación del condicional: $\neg(P\Rightarrow Q)\eqlog P\wedge\neg Q$;
  \item bicondicional: $P\Leftrightarrow Q\eqlog(P\Rightarrow Q)\wedge(Q\Rightarrow P)$.
\end{enumerate}
Además, $\wedge$ y $\vee$ son conmutativas y asociativas.
\end{proposicion}

\begin{proof}
Cada apartado se comprueba con una tabla de verdad. Se detallan (ii), (iv) y (v); los demás son análogos (\cref{eje:1-leyes}).

(ii) $\neg(P\wedge Q)$ es verdadera exactamente cuando $P\wedge Q$ es falsa, es decir, cuando al menos una de $P$ o $Q$ es falsa, es decir, cuando al menos una de $\neg P$ o $\neg Q$ es verdadera, que es la condición para que $\neg P\vee\neg Q$ sea verdadera. La segunda ley se razona igual intercambiando los papeles de «todas» y «al menos una».

(iv) $\neg P\vee Q$ es falsa exactamente cuando $\neg P$ y $Q$ son falsas, es decir, cuando $P$ es verdadera y $Q$ falsa. Es la misma única fila en la que $P\Rightarrow Q$ es falsa.

(v) Por (iv), $\neg Q\Rightarrow\neg P\eqlog\neg\neg Q\vee\neg P$, que por (i) y la conmutatividad equivale a $\neg P\vee Q$, y de nuevo por (iv), a $P\Rightarrow Q$.
\end{proof}

Las leyes de la \cref{prop:leyes} permiten transformar fórmulas sin construir tablas, igual que las propiedades de la suma y el producto permiten manipular expresiones algebraicas~\cite[§3.4]{mcs}.

\begin{ejemplo}\label{ej:simplificar}
Se simplifica $\neg(P\Rightarrow\neg Q)$:
\[
\neg(P\Rightarrow\neg Q)\ \eqlog\ P\wedge\neg\neg Q\ \eqlog\ P\wedge Q,
\]
por (vi) y (i). Así, negar «si llueve, no salgo» equivale a afirmar «llueve y salgo».
\end{ejemplo}

\section{Razonamientos válidos}\label{sec:inferencia}

\begin{definicion}[Razonamiento válido]
Un razonamiento con premisas $\Phi_1,\dots,\Phi_k$ y conclusión $\Psi$ es \emph{válido} si $(\Phi_1\wedge\dots\wedge\Phi_k)\Rightarrow\Psi$ es una tautología, es decir, si en toda asignación que hace verdaderas todas las premisas también es verdadera la conclusión.
\end{definicion}

\begin{proposicion}[Reglas de inferencia]\label{prop:inferencia}
Los siguientes razonamientos son válidos~\cite[§2.11]{hammack}:
\begin{center}
\begin{tabular}{@{}lll@{}}
\toprule
Nombre & Premisas & Conclusión \\
\midrule
\emph{Modus ponens} & $P\Rightarrow Q$, \ $P$ & $Q$ \\
\emph{Modus tollens} & $P\Rightarrow Q$, \ $\neg Q$ & $\neg P$ \\
Silogismo hipotético & $P\Rightarrow Q$, \ $Q\Rightarrow R$ & $P\Rightarrow R$ \\
Silogismo disyuntivo & $P\vee Q$, \ $\neg P$ & $Q$ \\
\bottomrule
\end{tabular}
\end{center}
\end{proposicion}

\begin{proof}
\emph{Modus ponens} se probó en el \cref{ej:tercero-excluido}. Para \emph{modus tollens}, si $P\Rightarrow Q$ y $\neg Q$ son verdaderas, por la \cref{prop:leyes}(v) también lo es $\neg Q\Rightarrow\neg P$, y \emph{modus ponens} da $\neg P$. Para el silogismo disyuntivo, si $\neg P$ es verdadera, $P$ es falsa, y como $P\vee Q$ es verdadera, ha de serlo $Q$. El silogismo hipotético es el \cref{eje:1-silogismo}.
\end{proof}

\begin{observacion}[Errores frecuentes]\label{obs:1-errores}
(a) Las \emph{falacias} más comunes imitan reglas válidas: de $P\Rightarrow Q$ y $Q$ \emph{no} se sigue $P$ (afirmar el consecuente), y de $P\Rightarrow Q$ y $\neg P$ \emph{no} se sigue $\neg Q$ (negar el antecedente). En ambos casos la asignación $P=\F$, $Q=\V$ hace verdaderas las premisas y falsa la conclusión.
(b) Demostrar el recíproco de un teorema no demuestra el teorema (\cref{ej:reciproco}).
(c) El «o» matemático es inclusivo; si se quiere excluir el caso «ambos», debe decirse explícitamente.
(d) $P\Rightarrow Q$ no afirma que $P$ sea verdadera ni que $P$ cause $Q$; solo excluye el caso $P$ verdadera y $Q$ falsa.
\end{observacion}

\begin{observacion}[Conexión con el curso]
Los enunciados de los módulos posteriores distinguen con cuidado entre condiciones necesarias y suficientes. Por ejemplo, en el módulo~4 el gradiente nulo es una condición \emph{necesaria} para que un punto interior sea mínimo de una función diferenciable, pero no suficiente; y el criterio de la hessiana definida positiva es \emph{suficiente}, pero no necesario. Confundir una condición con su recíproca es el error de la \cref{obs:1-errores}(b). En el módulo~13, la máscara causal de la atención es literalmente una tabla de valores de verdad: la posición $i$ puede atender a la $j$ si y solo si $j\le i$.
\end{observacion}

\begin{lecturas}
\item Hammack~\cite{hammack}, cap.~2 «Logic», §2.1–2.6 y §2.11 «Logical Inference».
\item Velleman~\cite{velleman}, §1.1 «Deductive Reasoning and Logical Connectives», §1.2 «Truth Tables» y §1.5 «The Conditional and Biconditional Connectives».
\item Lehman, Leighton y Meyer~\cite{mcs}, §3.1 «Propositions from Propositions», §3.3 «Equivalence and Validity» y §3.4 «The Algebra of Propositions».
\item Stanford CS103~\cite{cs103}, clase~3 «Propositional Logic».
\end{lecturas}

\section*{Ejercicios del tema 1}
\addcontentsline{toc}{section}{Ejercicios del tema 1}

\begin{ejercicio}{1}{eje:1-proposiciones}
Indique cuáles de los siguientes enunciados son proposiciones y, en su caso, su valor de verdad: (a) «$2^{10}>1000$»; (b) «$x^{2}\ge 0$»; (c) «el número $91$ es primo»; (d) «calcule $\sqrt{2}$ con tres decimales»; (e) «existe un número real $x$ con $x^{2}=-1$»; (f) «todo entero es par o impar».
\end{ejercicio}

\begin{ejercicio}{1}{eje:1-traducir}
Sean $P$: «$n$ es par», $Q$: «$n$ es múltiplo de~3» y $R$: «$n$ es múltiplo de~6», donde $n$ es un entero fijo. Escriba en símbolos: (a) «$n$ es múltiplo de~6 si y solo si es par y múltiplo de~3»; (b) «$n$ es par solo si es múltiplo de~6»; (c) «ser par no es suficiente para ser múltiplo de~6»; (d) «$n$ no es par ni múltiplo de~3». Indique cuáles de (a)–(c) son verdaderas para todo entero $n$.
\end{ejercicio}

\begin{ejercicio}{1}{eje:1-reciproco}
Para el enunciado «si una función $f\colon\R\to\R$ es derivable en $a$, entonces es continua en $a$», escriba el recíproco, el contrarrecíproco y el inverso, e indique cuáles son verdaderos. Para los falsos, dé una función que lo muestre.
\fuente{adaptado de Hammack~\cite{hammack}, §2.8 «More on Conditional Statements».}
\end{ejercicio}

\begin{ejercicio}{1}{eje:1-necesaria}
Para un entero $n$, clasifique cada condición como necesaria, suficiente, ambas o ninguna para que «$n$ sea múltiplo de~4»: (a) $n$ es par; (b) $n$ es múltiplo de~8; (c) $n^{2}$ es múltiplo de~16; (d) $n$ termina en~4.
\end{ejercicio}

\begin{ejercicio}{2}{eje:1-silogismo}
Construya la tabla de verdad de $\bigl((P\Rightarrow Q)\wedge(Q\Rightarrow R)\bigr)\Rightarrow(P\Rightarrow R)$ y concluya que el silogismo hipotético de la \cref{prop:inferencia} es válido.
\fuente{adaptado de Hammack~\cite{hammack}, §2.5 «Truth Tables for Statements».}
\end{ejercicio}

\begin{ejercicio}{2}{eje:1-xor}
La disyunción exclusiva $P\xor Q$ es verdadera si y solo si exactamente una de $P$ y $Q$ lo es. (a) Escriba su tabla de verdad. (b) Compruebe que $P\xor Q\eqlog(P\vee Q)\wedge\neg(P\wedge Q)$. (c) Compruebe que $P\xor Q\eqlog\neg(P\Leftrightarrow Q)$.
\end{ejercicio}

\begin{ejercicio}{2}{eje:1-equivalencias}
Decida con tablas de verdad si son equivalentes: (a) $P\Rightarrow(Q\Rightarrow R)$ y $(P\wedge Q)\Rightarrow R$; (b) $(P\Rightarrow Q)\Rightarrow R$ y $P\Rightarrow(Q\Rightarrow R)$; (c) $(P\vee Q)\Rightarrow R$ y $(P\Rightarrow R)\wedge(Q\Rightarrow R)$. Cuando no lo sean, indique una asignación que las distinga.
\fuente{adaptado de Velleman~\cite{velleman}, §1.5 «The Conditional and Biconditional Connectives».}
\end{ejercicio}

\begin{ejercicio}{2}{eje:1-algebra}
Usando solo las leyes de la \cref{prop:leyes} (sin tablas), demuestre que $\neg\bigl(P\vee(\neg P\wedge Q)\bigr)\eqlog\neg P\wedge\neg Q$ y que $(P\Rightarrow Q)\wedge(P\Rightarrow\neg Q)\eqlog\neg P$.
\fuente{adaptado de Lehman, Leighton y Meyer~\cite{mcs}, §3.4 «The Algebra of Propositions».}
\end{ejercicio}

\begin{ejercicio}{2}{eje:1-validez}
Decida si los razonamientos siguientes son válidos e identifique la regla o la falacia correspondiente. (a) «Si el modelo está sobreajustado, el error de validación sube. El error de validación sube. Luego el modelo está sobreajustado.» (b) «Si $n$ es primo y mayor que~2, entonces es impar. $n$ es par. Luego $n$ no es un primo mayor que~2.» (c) «El fallo está en los datos o en el código. No está en los datos. Luego está en el código.» (d) «Si $x>2$, entonces $x^{2}>4$. $x\le 2$. Luego $x^{2}\le 4$.»
\end{ejercicio}

\begin{ejercicio}{2}{eje:1-contar}
(a) ¿Cuántas tablas de verdad distintas puede tener una fórmula con $n$ variables? Calcule el número para $n=1,2,3$. (b) Para $n=2$, ¿cuántas de esas funciones son tautologías y cuántas contradicciones? (c) Escriba una fórmula para cada una de las cuatro funciones de una variable.
\end{ejercicio}

\begin{ejercicio}{2}{eje:1-leyes}
Complete la demostración de la \cref{prop:leyes} comprobando con tablas de verdad las leyes distributivas (iii) y las leyes (vi) y (vii).
\end{ejercicio}

\begin{ejercicio}{3}{eje:1-demorgan-n}
Sea $n\ge 2$ y sean $P_1,\dots,P_n$ proposiciones. Demuestre las leyes de De Morgan generalizadas
\[
\neg(P_1\wedge\dots\wedge P_n)\eqlog\neg P_1\vee\dots\vee\neg P_n,\qquad
\neg(P_1\vee\dots\vee P_n)\eqlog\neg P_1\wedge\dots\wedge\neg P_n,
\]
razonando directamente sobre los valores de verdad.
\end{ejercicio}

\begin{ejercicio}{3}{eje:1-completitud}
Demuestre que toda tabla de verdad con $n\ge 1$ variables es la de alguna fórmula que solo usa las conectivas $\neg$ y $\wedge$. (Indicación: construya primero una fórmula con $\neg$, $\wedge$ y $\vee$ que sea una disyunción de conjunciones, una por cada fila en la que la tabla vale $\V$, y elimine después $\vee$ con De Morgan.)
\fuente{adaptado de Lehman, Leighton y Meyer~\cite{mcs}, §3.4 «The Algebra of Propositions» (formas normales).}
\end{ejercicio}

\begin{ejercicio}{3}{eje:1-bicond-asoc}
Demuestre que el bicondicional es asociativo, $(P\Leftrightarrow Q)\Leftrightarrow R\eqlog P\Leftrightarrow(Q\Leftrightarrow R)$, probando que ambos miembros son verdaderos exactamente cuando el número de proposiciones falsas entre $P$, $Q$ y $R$ es par. Demuestre después, con una asignación concreta, que el condicional no es asociativo.
\end{ejercicio}

% =====================================================================
\chapter{Cuantificadores y negación}
\label{tema:2}
% =====================================================================

Casi ningún teorema habla de un único objeto: afirman algo de \emph{todos} los enteros, de \emph{alguna} raíz de un polinomio o de que \emph{para cada} $\varepsilon$ \emph{existe} un $\delta$. Este tema introduce los predicados y los cuantificadores, estudia el efecto del orden de los cuantificadores y da reglas sistemáticas para negar enunciados, que serán imprescindibles en las demostraciones por contrarrecíproco y por reducción al absurdo. Se siguen las §2.7–2.10 de~\cite{hammack}, el capítulo~2 de~\cite{velleman} (§2.1 y §2.2), la §3.6 de~\cite{mcs} y la guía de negación del curso CS103~\cite{cs103}.

\section{Predicados}

\begin{definicion}[Predicado]\label{def:predicado}
Sea $U$ un conjunto no vacío, llamado \emph{universo} o dominio. Un \emph{predicado} sobre $U$ es un enunciado $P(x)$ que contiene una variable $x$ y que se convierte en una proposición cada vez que se sustituye $x$ por un elemento concreto de $U$. El \emph{conjunto de verdad} de $P$ es $\{x\in U: P(x)\text{ es verdadera}\}$. Del mismo modo se definen predicados de varias variables, $P(x,y)$, con $(x,y)\in U\times W$.
\end{definicion}

\begin{ejemplo}
Con $U=\Z$, el predicado $P(n)$: «$n^{2}-1$ es múltiplo de~8» da proposiciones como $P(3)$ ($8\divides 8$, verdadera) o $P(2)$ ($8\nodivide 3$, falsa). Su conjunto de verdad es el de los enteros impares (\cref{eje:4-24} prueba una versión más fuerte).
\end{ejemplo}

\section{Cuantificador universal y existencial}

\begin{definicion}[Cuantificadores]\label{def:cuantificadores}
Sea $P(x)$ un predicado sobre $U$.
\begin{enumerate}
  \item $\forall x\in U\ P(x)$ («para todo $x$ de $U$, $P(x)$») es verdadera si $P(a)$ es verdadera para cada $a\in U$.
  \item $\exists x\in U\ P(x)$ («existe $x$ en $U$ tal que $P(x)$») es verdadera si $P(a)$ es verdadera para al menos un $a\in U$.
\end{enumerate}
Cuando el universo está claro se escribe $\forall x\,P(x)$ y $\exists x\,P(x)$. La variable $x$ está \emph{ligada} por el cuantificador: $\forall x\,P(x)$ y $\forall t\,P(t)$ dicen lo mismo.
\end{definicion}

Los cuantificadores se combinan con las conectivas. Las formas más frecuentes son la \emph{universal condicional}, $\forall x\,\bigl(P(x)\Rightarrow Q(x)\bigr)$, que traduce «todo $P$ es $Q$», y la \emph{existencial conjuntiva}, $\exists x\,\bigl(P(x)\wedge Q(x)\bigr)$, que traduce «algún $P$ es $Q$». Se abrevian como $\forall x\in A\ Q(x)$ y $\exists x\in A\ Q(x)$ cuando $P(x)$ es «$x\in A$».

\begin{ejemplo}\label{ej:traduccion}
En $U=\Z$: «todo cuadrado perfecto es no negativo» es $\forall n\,(n^{2}\ge 0)$; «hay un primo par» es $\exists p\,(p\text{ primo}\wedge p\text{ par})$; «ningún entero es par e impar a la vez» es $\neg\exists n\,(n\text{ par}\wedge n\text{ impar})$, o bien $\forall n\,\neg(n\text{ par}\wedge n\text{ impar})$.
\end{ejemplo}

\begin{observacion}
Si $A=\varnothing$, el enunciado $\forall x\in A\ Q(x)$ es verdadero, sea cual sea $Q$: no hay ningún $x\in A$ que lo incumpla. Es la versión cuantificada del condicional vacuamente verdadero del \cref{tema:1}. En cambio, $\exists x\in\varnothing\ Q(x)$ es siempre falso.
\end{observacion}

\section{Cuantificadores anidados}

Con predicados de varias variables, los cuantificadores se anidan. El orden importa cuando se mezclan cuantificadores de distinto tipo.

\begin{ejemplo}\label{ej:orden}
Sea $P(x,y)$: «$x+y=0$» con $x,y\in\Z$. El enunciado $\forall x\,\exists y\,P(x,y)$ dice que todo entero tiene un opuesto, y es verdadero: dado $x$, sirve $y=-x$. El enunciado $\exists y\,\forall x\,P(x,y)$ dice que hay un entero $y$ que es opuesto de \emph{todos} los enteros a la vez, y es falso: si existiera, $0+y=0$ y $1+y=0$ darían $y=0$ e $y=-1$. En el primer enunciado $y$ puede depender de $x$; en el segundo, no.
\end{ejemplo}

\begin{proposicion}\label{prop:orden}
Para todo predicado $P(x,y)$,
\[
\exists y\,\forall x\,P(x,y)\ \Longrightarrow\ \forall x\,\exists y\,P(x,y),
\]
y la implicación recíproca es falsa en general.
\end{proposicion}

\begin{proof}
Supóngase $\exists y\,\forall x\,P(x,y)$ y sea $y_0$ un elemento tal que $P(x,y_0)$ es verdadera para todo $x$. Sea $x$ arbitrario. Entonces $P(x,y_0)$ es verdadera, y por tanto existe $y$ (a saber, $y_0$) con $P(x,y)$. Como $x$ era arbitrario, $\forall x\,\exists y\,P(x,y)$. El \cref{ej:orden} muestra que la recíproca puede fallar.
\end{proof}

Cuantificadores del mismo tipo sí pueden intercambiarse: $\forall x\,\forall y\,P(x,y)\eqlog\forall y\,\forall x\,P(x,y)$ y $\exists x\,\exists y\,P(x,y)\eqlog\exists y\,\exists x\,P(x,y)$.

\begin{ejemplo}[Convergencia de una sucesión]\label{ej:limite}
Que una sucesión real $(a_n)_{n\in\N}$ converja a $L$ se escribe
\[
\forall\varepsilon>0\ \exists N\in\N\ \forall n\in\N\ \bigl(n\ge N\Rightarrow\abs{a_n-L}<\varepsilon\bigr).
\]
El índice $N$ puede depender de $\varepsilon$: cuanto más exigente es la tolerancia, más lejos hay que ir en la sucesión. Intercambiar los dos primeros cuantificadores daría un enunciado mucho más fuerte y casi nunca cierto (\cref{eje:2-arquimedes}).
\end{ejemplo}

\section{Negación de enunciados cuantificados}

\begin{teorema}[Negación de cuantificadores]\label{teo:negacion}
Para todo predicado $P(x)$ sobre un universo $U$,
\[
\neg\,\forall x\,P(x)\ \eqlog\ \exists x\,\neg P(x),\qquad
\neg\,\exists x\,P(x)\ \eqlog\ \forall x\,\neg P(x).
\]
\end{teorema}

\begin{proof}
$\neg\,\forall x\,P(x)$ es verdadera si y solo si no es cierto que $P(a)$ valga para todo $a\in U$, es decir, si y solo si hay algún $a\in U$ con $P(a)$ falsa, que es la condición para que $\exists x\,\neg P(x)$ sea verdadera. Para la segunda equivalencia se aplica la primera al predicado $\neg P$ y se niegan ambos miembros: $\neg\,\forall x\,\neg P(x)\eqlog\exists x\,\neg\neg P(x)\eqlog\exists x\,P(x)$; negando, $\forall x\,\neg P(x)\eqlog\neg\,\exists x\,P(x)$.
\end{proof}

El \cref{teo:negacion} y las leyes de la \cref{prop:leyes} dan un procedimiento mecánico: la negación «atraviesa» los cuantificadores cambiando $\forall$ por $\exists$ y viceversa, y al llegar a la parte proposicional se aplican De Morgan y la negación del condicional, $\neg(P\Rightarrow Q)\eqlog P\wedge\neg Q$. Con universos acotados, el cuantificador conserva su dominio: $\neg\,\forall x\in A\ Q(x)\eqlog\exists x\in A\ \neg Q(x)$.

\begin{ejemplo}\label{ej:negar-limite}
Se niega la definición de convergencia del \cref{ej:limite}:
\begin{align*}
&\neg\,\forall\varepsilon>0\ \exists N\ \forall n\ \bigl(n\ge N\Rightarrow\abs{a_n-L}<\varepsilon\bigr)\\
\eqlog\ &\exists\varepsilon>0\ \forall N\ \exists n\ \neg\bigl(n\ge N\Rightarrow\abs{a_n-L}<\varepsilon\bigr)\\
\eqlog\ &\exists\varepsilon>0\ \forall N\ \exists n\ \bigl(n\ge N\wedge\abs{a_n-L}\ge\varepsilon\bigr).
\end{align*}
Es decir, $(a_n)$ \emph{no} converge a $L$ si hay una tolerancia $\varepsilon$ tal que, por lejos que se mire, siempre aparece algún término a distancia al menos $\varepsilon$ de $L$. Con $a_n=(-1)^{n}$ y $L=1$ sirve $\varepsilon=1$: para cualquier $N$, el término impar $n=2N+1\ge N$ cumple $\abs{a_n-1}=2\ge 1$.
\end{ejemplo}

\section{Traducción y unicidad}

Las expresiones del lenguaje natural ocultan a menudo los cuantificadores. «Los primos mayores que~2 son impares» es universal; «un cuadrado puede ser par» es existencial; «solo los múltiplos de~6 son múltiplos de~2 y de~3» significa «si $n$ es múltiplo de~2 y de~3, entonces es múltiplo de~6». Antes de demostrar un enunciado conviene escribirlo con cuantificadores explícitos~\cite[§2.9]{hammack}.

\begin{definicion}[Existencia y unicidad]\label{def:unicidad}
$\exists!\,x\,P(x)$ («existe un único $x$ tal que $P(x)$») abrevia
\[
\exists x\,\Bigl(P(x)\wedge\forall y\,\bigl(P(y)\Rightarrow y=x\bigr)\Bigr).
\]
\end{definicion}

Así, $\exists!$ contiene dos afirmaciones: hay al menos un elemento con la propiedad (\emph{existencia}) y dos elementos con la propiedad son iguales (\emph{unicidad}). El \cref{tema:6} explica cómo se demuestra cada parte.

\begin{observacion}[Errores frecuentes]\label{obs:2-errores}
(a) Negar «todos los $x$ cumplen $P$» como «ningún $x$ cumple $P$»: la negación correcta es «algún $x$ no cumple $P$».
(b) Negar «$\forall x\,(P(x)\Rightarrow Q(x))$» como otro condicional: la negación es $\exists x\,(P(x)\wedge\neg Q(x))$, una conjunción.
(c) Cambiar el orden de $\forall$ y $\exists$ al copiar una definición (\cref{ej:orden}).
(d) Olvidar el universo: $\exists x\,(x^{2}=2)$ es verdadero en $\R$ y falso en $\Q$.
\end{observacion}

\begin{observacion}[Conexión con el curso]
La definición $\varepsilon$–$N$ del \cref{ej:limite} y su análoga $\varepsilon$–$\delta$ para funciones son la base de los límites y derivadas del módulo~4, de los resultados de convergencia de los algoritmos de optimización del módulo~7 y del análisis de la softmax con temperatura del módulo~1, donde se estudia el límite cuando la temperatura tiende a~0. Saber negar estas definiciones es lo que permite demostrar que un algoritmo \emph{no} converge o que una función \emph{no} es continua.
\end{observacion}

\begin{lecturas}
\item Hammack~\cite{hammack}, §2.7 «Quantifiers», §2.9 «Translating English to Symbolic Logic» y §2.10 «Negating Statements».
\item Velleman~\cite{velleman}, §2.1 «Quantifiers» y §2.2 «Equivalences Involving Quantifiers».
\item Lehman, Leighton y Meyer~\cite{mcs}, §1.2 «Predicates» y §3.6 «Predicate Formulas».
\item Stanford CS103~\cite{cs103}, clases~4 y~5 «First-Order Logic» y la guía «Guide to Negation».
\end{lecturas}

\section*{Ejercicios del tema 2}
\addcontentsline{toc}{section}{Ejercicios del tema 2}

\begin{ejercicio}{1}{eje:2-valores}
Determine el valor de verdad: (a) $\forall x\in\R\ (x^{2}\ge 0)$; (b) $\exists x\in\Q\ (x^{2}=2)$; (c) $\exists x\in\R\ (x^{2}=2)$; (d) $\forall n\in\N\ \exists m\in\N\ (m<n)$; (e) $\exists m\in\N\ \forall n\in\N\ (m\le n)$; (f) $\forall x\in\varnothing\ (x\neq x)$.
\end{ejercicio}

\begin{ejercicio}{1}{eje:2-simbolos}
Escriba con cuantificadores, sobre el universo indicado: (a) «todo entero par mayor que~2 es suma de dos primos» ($\Z$); (b) «hay un número real que es mayor que todos los enteros» y decida si es verdadero; (c) «ningún cuadrado de un entero termina en~2»; (d) «la ecuación $x^{2}+1=0$ no tiene solución real».
\end{ejercicio}

\begin{ejercicio}{1}{eje:2-lenguaje}
Exprese en lenguaje natural y determine el valor de verdad: (a) $\forall x\in\R\ \exists y\in\R\ (y>x)$; (b) $\exists y\in\R\ \forall x\in\R\ (y>x)$; (c) $\forall x\in\R\ \bigl(x\neq 0\Rightarrow\exists y\in\R\ (xy=1)\bigr)$; (d) $\exists x\in\Z\ \forall y\in\Z\ (xy=y)$.
\end{ejercicio}

\begin{ejercicio}{2}{eje:2-negar}
Niegue los enunciados siguientes de modo que el símbolo $\neg$ no aparezca delante de ningún cuantificador ni paréntesis, y decida cuál es verdadero, el original o su negación: (a) $\forall x\in\R\ \exists y\in\R\ (x<y<x+1)$; (b) $\exists n\in\N\ \forall m\in\N\ (n\divides m)$; (c) $\forall x\in\R\ \bigl(x^{2}>4\Rightarrow x>2\bigr)$; (d) $\forall n\in\Z\ \bigl(n\text{ primo}\Rightarrow n\text{ impar}\vee n=2\bigr)$.
\fuente{adaptado de Hammack~\cite{hammack}, §2.10 «Negating Statements».}
\end{ejercicio}

\begin{ejercicio}{2}{eje:2-continuidad}
Una función $f\colon\R\to\R$ es continua en $a$ si $\forall\varepsilon>0\ \exists\delta>0\ \forall x\in\R\ \bigl(\abs{x-a}<\delta\Rightarrow\abs{f(x)-f(a)}<\varepsilon\bigr)$. (a) Escriba la negación sin $\neg$ delante de cuantificadores. (b) Use la negación para demostrar que la función $f(x)=0$ si $x<0$ y $f(x)=1$ si $x\ge 0$ no es continua en $a=0$.
\fuente{adaptado de Velleman~\cite{velleman}, §2.2 «Equivalences Involving Quantifiers».}
\end{ejercicio}

\begin{ejercicio}{2}{eje:2-orden}
Para cada predicado, determine el valor de verdad de $\forall x\,\exists y\,P(x,y)$ y de $\exists y\,\forall x\,P(x,y)$: (a) $P(x,y)$: «$x+y=0$» en $\N_0$; (b) $P(x,y)$: «$xy=1$» en $\Q\setminus\{0\}$; (c) $P(x,y)$: «$x\le y$» en $\N$; (d) $P(x,y)$: «$y\le x$» en $\N$.
\end{ejercicio}

\begin{ejercicio}{2}{eje:2-distribuir}
Decida cuáles de las siguientes equivalencias son ciertas para todos los predicados $P$ y $Q$ sobre cualquier universo, y dé un contraejemplo para las que no lo sean: (a) $\forall x\,(P(x)\wedge Q(x))\eqlog\forall x\,P(x)\wedge\forall x\,Q(x)$; (b) $\forall x\,(P(x)\vee Q(x))\eqlog\forall x\,P(x)\vee\forall x\,Q(x)$; (c) $\exists x\,(P(x)\vee Q(x))\eqlog\exists x\,P(x)\vee\exists x\,Q(x)$; (d) $\exists x\,(P(x)\wedge Q(x))\eqlog\exists x\,P(x)\wedge\exists x\,Q(x)$.
\fuente{adaptado de Velleman~\cite{velleman}, §2.2 «Equivalences Involving Quantifiers».}
\end{ejercicio}

\begin{ejercicio}{3}{eje:2-distribuir-dem}
Demuestre las equivalencias (a) y (c) del \cref{eje:2-distribuir}. Demuestre además que $\forall x\,P(x)\vee\forall x\,Q(x)\Rightarrow\forall x\,(P(x)\vee Q(x))$ es siempre verdadera.
\end{ejercicio}

\begin{ejercicio}{3}{eje:2-vacio}
Demuestre que, si el universo $U$ es no vacío, $\forall x\,P(x)\Rightarrow\exists x\,P(x)$ para todo predicado $P$. Explique por qué la implicación falla si se permite $U=\varnothing$.
\end{ejercicio}

\begin{ejercicio}{2}{eje:2-unicidad}
(a) Escriba $\exists!\,x\,P(x)$ sin el símbolo $\exists!$ de una forma distinta de la \cref{def:unicidad}, como conjunción de un enunciado de existencia y otro de unicidad. (b) Niegue $\exists!\,x\,P(x)$ y explique en palabras las dos maneras en que puede fallar. (c) Decida si $\exists!\,x\in\R\ (x^{2}=4)$ y $\exists!\,x\in\N\ (x^{2}=4)$ son verdaderos.
\end{ejercicio}

\begin{ejercicio}{2}{eje:2-acotada}
Una función $f\colon\R\to\R$ es \emph{acotada} si $\exists M\in\R\ \forall x\in\R\ (\abs{f(x)}\le M)$, y es \emph{creciente} si $\forall x\,\forall y\ (x<y\Rightarrow f(x)\le f(y))$. (a) Niegue ambas definiciones. (b) Demuestre, usando la negación, que $f(x)=x^{2}$ no es acotada ni creciente.
\end{ejercicio}

\begin{ejercicio}{3}{eje:2-arquimedes}
Se admite la \emph{propiedad arquimediana}: para todo número real $x$ existe $N\in\N$ con $N>x$. Demuestre que (a) $\forall\varepsilon>0\ \exists N\in\N\ \forall n\in\N\ (n\ge N\Rightarrow 1/n<\varepsilon)$ es verdadero, y que (b) $\exists N\in\N\ \forall\varepsilon>0\ \forall n\in\N\ (n\ge N\Rightarrow 1/n<\varepsilon)$ es falso.
\end{ejercicio}

\begin{ejercicio}{2}{eje:2-critica}
Se pretende demostrar que «todo número natural es impar» así: «Por reducción al absurdo, supóngase que todo número natural es par. Entonces $137$ es par, pero $137=2\cdot 68+1$ es impar. Contradicción; luego todo natural es impar.» Identifique el error lógico, escriba la negación correcta del enunciado y decida si este es verdadero.
\fuente{adaptado de Stanford CS103~\cite{cs103}, \emph{Problem Set}~1, problema~9 «Proof Critiques».}
\end{ejercicio}

\begin{ejercicio}{3}{eje:2-finito}
Sea $U=\{a_1,\dots,a_n\}$ un universo finito. Justifique que $\forall x\,P(x)\eqlog P(a_1)\wedge\dots\wedge P(a_n)$ y que $\exists x\,P(x)\eqlog P(a_1)\vee\dots\vee P(a_n)$, y deduzca el \cref{teo:negacion} para universos finitos a partir de las leyes de De Morgan generalizadas del \cref{eje:1-demorgan-n}.
\end{ejercicio}

% =====================================================================
\chapter{Demostración directa y por contrarrecíproco}
\label{tema:3}
% =====================================================================

La mayoría de los teoremas tienen la forma $\forall x\,(P(x)\Rightarrow Q(x))$. Este tema presenta las dos técnicas básicas para demostrarlos: la demostración directa, que parte de $P(x)$ y llega a $Q(x)$, y la demostración por contrarrecíproco, que parte de $\neg Q(x)$ y llega a $\neg P(x)$. Como banco de pruebas se usa la aritmética de los enteros (paridad, divisibilidad y congruencias), igual que en los capítulos~4 y~5 de~\cite{hammack}, la §3.2 de~\cite{velleman}, las §1.5–1.6 de~\cite{mcs} y el primer \emph{problem set} de CS103~\cite{cs103}.

\section{Definiciones de partida}

Una demostración solo puede usar lo que está definido o demostrado antes. Se admiten como conocidas las propiedades algebraicas habituales de $\Z$ (la suma, la resta y el producto de enteros son enteros; asociatividad, conmutatividad, distributividad) y de $\R$ (incluidas las reglas de las desigualdades), y se fijan las definiciones siguientes~\cite[§4.2]{hammack}.

\begin{definicion}[Paridad]\label{def:paridad}
Un entero $n$ es \emph{par} si existe $k\in\Z$ con $n=2k$, e \emph{impar} si existe $k\in\Z$ con $n=2k+1$. Dos enteros tienen la \emph{misma paridad} si ambos son pares o ambos impares.
\end{definicion}

\begin{definicion}[Divisibilidad]\label{def:divide}
Sean $a,b\in\Z$. Se dice que $a$ \emph{divide} a $b$, y se escribe $a\divides b$, si existe $c\in\Z$ con $b=ac$. Se dice también que $b$ es \emph{múltiplo} de $a$ y que $a$ es \emph{divisor} de $b$. Un entero $p\ge 2$ es \emph{primo} si sus únicos divisores positivos son $1$ y $p$; un entero $n\ge 2$ que no es primo es \emph{compuesto}.
\end{definicion}

\begin{teorema}[División euclídea]\label{teo:division}
Sean $a\in\Z$ y $b\in\N$. Existen enteros únicos $q$ y $r$ con $a=qb+r$ y $0\le r<b$.
\end{teorema}

La existencia se demuestra en el \cref{tema:5} (\cref{prop:division-existencia}) y la unicidad en el \cref{tema:6} (\cref{prop:division-unicidad}). Con $b=2$, el teorema dice que todo entero es par o impar, y no ambas cosas; este hecho se usará libremente.

\section{Demostración directa}

\begin{definicion}[Demostración directa]\label{def:directa}
Para demostrar $\forall x\,(P(x)\Rightarrow Q(x))$ de forma \emph{directa}, se toma un $x$ arbitrario, se supone $P(x)$ y, mediante definiciones, resultados previos y reglas de inferencia, se deduce $Q(x)$.
\end{definicion}

El esquema es siempre el mismo: «Sea $x$ [del universo] tal que $P(x)$. [\dots] Por tanto, $Q(x)$.» Antes de empezar conviene escribir qué significa la tesis según las definiciones, para saber a qué expresión hay que llegar~\cite[§4.3]{hammack}.

\begin{proposicion}\label{prop:impar-cuadrado}
Si $n$ es un entero impar, entonces $n^{2}$ es impar.
\end{proposicion}

\begin{proof}
Sea $n$ un entero impar. Por definición existe $k\in\Z$ con $n=2k+1$. Entonces
\[
n^{2}=4k^{2}+4k+1=2(2k^{2}+2k)+1.
\]
Como $m=2k^{2}+2k$ es un entero, $n^{2}=2m+1$ es impar.
\end{proof}

Obsérvese que la demostración introduce un nombre nuevo, $m$, para el entero que da la forma $2m+1$: el objetivo era encontrar \emph{algún} entero con esa propiedad, y la definición de impar no exige que sea el mismo $k$ de la hipótesis.

\begin{proposicion}[Propiedades de la divisibilidad]\label{prop:divisibilidad}
Sean $a,b,c,x,y\in\Z$.
\begin{enumerate}
  \item Si $a\divides b$ y $b\divides c$, entonces $a\divides c$.
  \item Si $a\divides b$ y $a\divides c$, entonces $a\divides(bx+cy)$.
\end{enumerate}
\end{proposicion}

\begin{proof}
(i) Supóngase $a\divides b$ y $b\divides c$. Existen $d,e\in\Z$ con $b=ad$ y $c=be$. Sustituyendo, $c=(ad)e=a(de)$ con $de\in\Z$, luego $a\divides c$.

(ii) Supóngase $a\divides b$ y $a\divides c$, y sean $d,e\in\Z$ con $b=ad$ y $c=ae$. Entonces $bx+cy=adx+aey=a(dx+ey)$, y $dx+ey\in\Z$, luego $a\divides(bx+cy)$.
\end{proof}

\begin{ejemplo}[Desigualdad con números reales]\label{ej:amgm}
Se demuestra que, si $x,y\in\R$ son positivos, entonces $\sqrt{xy}\le\frac{x+y}{2}$. Una forma de \emph{descubrir} la prueba es trabajar hacia atrás: la desigualdad equivale a $2\sqrt{xy}\le x+y$, es decir, a $0\le x-2\sqrt{xy}+y=(\sqrt x-\sqrt y)^{2}$, que es cierta. La demostración se \emph{redacta}, sin embargo, hacia delante: «Sean $x,y>0$. Como el cuadrado de un número real es no negativo, $0\le(\sqrt x-\sqrt y)^{2}=x-2\sqrt{xy}+y$. Sumando $2\sqrt{xy}$ y dividiendo por~2 se obtiene $\sqrt{xy}\le(x+y)/2$.» Cada paso de la redacción es una consecuencia de los anteriores, y no al revés.
\end{ejemplo}

\section{Demostración por contrarrecíproco}

\begin{definicion}[Demostración por contrarrecíproco]\label{def:contrarreciproco}
Para demostrar $\forall x\,(P(x)\Rightarrow Q(x))$ por \emph{contrarrecíproco}, se demuestra de forma directa $\forall x\,(\neg Q(x)\Rightarrow\neg P(x))$. Es legítimo porque ambos enunciados son lógicamente equivalentes (\cref{prop:leyes}(v)).
\end{definicion}

La técnica es útil cuando la negación de la tesis aporta información más manejable que la hipótesis. Un indicio típico es que la tesis afirme algo sobre un número ($n$ es par) y la hipótesis, algo sobre una expresión de ese número ($n^{2}$ es par)~\cite[§5.1]{hammack}.

\begin{proposicion}\label{prop:cuadrado-par}
Sea $n\in\Z$. Si $n^{2}$ es par, entonces $n$ es par.
\end{proposicion}

\begin{proof}
Se demuestra el contrarrecíproco: si $n$ no es par, entonces $n^{2}$ no es par. Si $n$ no es par, es impar, y por la \cref{prop:impar-cuadrado} $n^{2}$ es impar, luego no es par.
\end{proof}

Un intento directo habría partido de $n^{2}=2k$ y necesitaría extraer una raíz cuadrada, que no conduce a ninguna parte. El contrarrecíproco convierte el problema en el cálculo ya hecho.

\begin{ejemplo}\label{ej:contra-real}
Sea $x\in\R$. Si $x^{3}-x>0$, entonces $x>-1$. El contrarrecíproco dice: si $x\le -1$, entonces $x^{3}-x\le 0$. En efecto, $x^{3}-x=x(x-1)(x+1)$; para $x\le -1$ se tiene $x<0$, $x-1<0$ y $x+1\le 0$, y el producto de dos factores negativos y uno no positivo es no positivo.
\end{ejemplo}

\section{Congruencias}

\begin{definicion}[Congruencia]\label{def:congruencia}
Sea $n\in\N$. Dos enteros $a$ y $b$ son \emph{congruentes módulo $n$}, y se escribe $a\equiv b\pmod n$, si $n\divides(a-b)$.
\end{definicion}

Por la división euclídea, todo entero es congruente módulo $n$ con exactamente uno de $0,1,\dots,n-1$, su resto. Por ejemplo, $n$ es par si y solo si $n\equiv 0\pmod 2$.

\begin{proposicion}\label{prop:congruencias}
Sea $n\in\N$ y sean $a,b,c,d\in\Z$ con $a\equiv b\pmod n$ y $c\equiv d\pmod n$. Entonces
\begin{enumerate}
  \item $a+c\equiv b+d\pmod n$;
  \item $ac\equiv bd\pmod n$;
  \item $a^{k}\equiv b^{k}\pmod n$ para todo $k\in\N$.
\end{enumerate}
\end{proposicion}

\begin{proof}
Por hipótesis existen $s,t\in\Z$ con $a-b=ns$ y $c-d=nt$.
(i) $(a+c)-(b+d)=ns+nt=n(s+t)$.
(ii) $ac-bd=ac-bc+bc-bd=(a-b)c+b(c-d)=n(sc+bt)$.
(iii) Se aplica (ii) con $c=a$ y $d=b$ repetidamente: $a^{2}\equiv b^{2}$, $a^{3}=a^{2}a\equiv b^{2}b=b^{3}$, etc.; la forma rigurosa de este «etc.» es la inducción del \cref{tema:5}.
\end{proof}

\begin{ejemplo}\label{ej:resto-potencia}
Se calcula el resto de dividir $2^{100}$ entre~7. Como $2^{3}=8\equiv 1\pmod 7$, la \cref{prop:congruencias}(iii) da $2^{99}=(2^{3})^{33}\equiv 1^{33}=1$, y por (ii), $2^{100}=2^{99}\cdot 2\equiv 2\pmod 7$. El resto es~2.
\end{ejemplo}

\section{Bicondicionales y cadenas de equivalencias}

\begin{definicion}[Demostración de un bicondicional]
Para demostrar $P\Leftrightarrow Q$ se demuestran por separado $P\Rightarrow Q$ y $Q\Rightarrow P$, cada una con la técnica más conveniente (\cref{prop:leyes}(vii)). Para demostrar que varias afirmaciones $P_1,\dots,P_k$ son equivalentes basta una cadena cíclica $P_1\Rightarrow P_2\Rightarrow\dots\Rightarrow P_k\Rightarrow P_1$~\cite[§7.2]{hammack}.
\end{definicion}

\begin{proposicion}\label{prop:par-equiv}
Para $n\in\Z$, son equivalentes: (a) $n$ es par; (b) $n^{2}$ es par; (c) $n^{2}+2n$ es múltiplo de~4.
\end{proposicion}

\begin{proof}
(a)$\Rightarrow$(c): si $n=2k$, entonces $n^{2}+2n=4k^{2}+4k=4(k^{2}+k)$.
(c)$\Rightarrow$(b): si $n^{2}+2n=4m$, entonces $n^{2}=4m-2n=2(2m-n)$ es par.
(b)$\Rightarrow$(a): es la \cref{prop:cuadrado-par}.
\end{proof}

\begin{observacion}[Errores frecuentes]\label{obs:3-errores}
(a) Suponer la tesis: una «demostración» que empieza escribiendo lo que se quiere probar y llega a algo cierto no demuestra nada, salvo que cada paso sea reversible y se redacte en el sentido correcto (\cref{ej:amgm}).
(b) Reutilizar una letra: si $a$ y $b$ son impares, no puede escribirse $a=2k+1$ y $b=2k+1$ con el mismo $k$, pues eso supone $a=b$.
(c) Comprobar unos cuantos casos ($n=1,3,5$) no demuestra un enunciado universal.
(d) Demostrar el recíproco $Q\Rightarrow P$ en lugar del contrarrecíproco $\neg Q\Rightarrow\neg P$.
\end{observacion}

\begin{observacion}[Conexión con el curso]
Una función $f$ es inyectiva si $f(x)=f(y)\Rightarrow x=y$, lo que por contrarrecíproco equivale a $x\neq y\Rightarrow f(x)\neq f(y)$. Las dos formas se usan en el curso: en el módulo~1 se prueba de forma directa que un tokenizador $\tau$ con decodificador $\delta$ cumple $\tau(s)=\tau(s')\Rightarrow s=s'$, aplicando $\delta$ a ambos lados; en el módulo~12 se demuestra la inyectividad de la codificación UTF-8 de manera análoga. Las congruencias reaparecen en la aritmética binaria y hexadecimal del módulo~12 y en el análisis de la coma flotante del módulo~11.
\end{observacion}

\begin{lecturas}
\item Hammack~\cite{hammack}, cap.~4 «Direct Proof» (§4.1–4.3) y cap.~5 «Contrapositive Proof» (§5.1–5.2).
\item Velleman~\cite{velleman}, §3.1 «Proof Strategies» y §3.2 «Proofs Involving Negations and Conditionals».
\item Lehman, Leighton y Meyer~\cite{mcs}, §1.5 «Proving an Implication» y §1.6 «Proving an “If and Only If”».
\item Stanford CS103~\cite{cs103}, clase~1 «Mathematical Proofs» y la guía «Guide to Proofs».
\end{lecturas}

\section*{Ejercicios del tema 3}
\addcontentsline{toc}{section}{Ejercicios del tema 3}

\begin{ejercicio}{1}{eje:3-hipotesis}
Para cada enunciado, identifique hipótesis y tesis, escriba el contrarrecíproco y diga qué técnica parece más sencilla: (a) «si $a\divides b$ y $b\divides c$, entonces $a\divides c$»; (b) «si $n^{2}$ es impar, entonces $n$ es impar»; (c) «si $x\in\R$ y $x^{2}+x<0$, entonces $x<0$»; (d) «si $xy$ es irracional, entonces $x$ o $y$ es irracional».
\end{ejercicio}

\begin{ejercicio}{1}{eje:3-esquema}
Explique qué está mal en esta demostración de «si $n$ es par, entonces $n^{2}$ es múltiplo de~4»: «Sea $n^{2}=4k$. Entonces $n=2\sqrt{k}$, que es par. Por tanto, el resultado es cierto.» Redáctela correctamente.
\end{ejercicio}

\begin{ejercicio}{3}{eje:3-cubo}
Demuestre que, si $x$ es un entero impar, entonces $x^{3}$ es impar.
\fuente{adaptado de Hammack~\cite{hammack}, cap.~4, ejercicio~2.}
\end{ejercicio}

\begin{ejercicio}{3}{eje:3-polinomio}
Demuestre que, si $a$ es un entero impar, entonces $a^{2}+3a+5$ es impar.
\fuente{adaptado de Hammack~\cite{hammack}, cap.~4, ejercicio~3.}
\end{ejercicio}

\begin{ejercicio}{3}{eje:3-cuadrados}
Sean $a,b\in\Z$. Demuestre que, si $a\divides b$, entonces $a^{2}\divides b^{2}$.
\fuente{adaptado de Hammack~\cite{hammack}, cap.~4, ejercicio~7.}
\end{ejercicio}

\begin{ejercicio}{3}{eje:3-producto}
Sean $a,b,c,d\in\Z$. Demuestre que, si $a\divides b$ y $c\divides d$, entonces $ac\divides bd$.
\fuente{adaptado de Hammack~\cite{hammack}, cap.~4, ejercicio~11.}
\end{ejercicio}

\begin{ejercicio}{3}{eje:3-desigualdad}
Demuestre que, si $x\in\R$ y $0<x<4$, entonces $\dfrac{4}{x(4-x)}\ge 1$.
\fuente{adaptado de Hammack~\cite{hammack}, cap.~4, ejercicio~12.}
\end{ejercicio}

\begin{ejercicio}{3}{eje:3-diferencia}
Demuestre que todo entero impar es diferencia de los cuadrados de dos enteros.
\fuente{adaptado de Hammack~\cite{hammack}, cap.~4, ejercicio~26.}
\end{ejercicio}

\begin{ejercicio}{3}{eje:3-contra-cubo}
Sea $n\in\Z$. Demuestre por contrarrecíproco que, si $n^{3}+5$ es impar, entonces $n$ es par.
\end{ejercicio}

\begin{ejercicio}{3}{eje:3-pitagoras}
Sean $a,b,c\in\Z$. Demuestre que, si $a^{2}+b^{2}=c^{2}$, entonces al menos uno de los tres números es par.
\fuente{adaptado de Stanford CS103~\cite{cs103}, \emph{Problem Set}~1, problema~5 «Writing Proofs by Contrapositive».}
\end{ejercicio}

\begin{ejercicio}{3}{eje:3-equivalencia}
Sea $k\in\N$ y defínase, para $x,y\in\Z$, $x\equiv_k y$ si existe $q\in\Z$ con $x=y+kq$. Demuestre que la relación es (a) reflexiva: $x\equiv_k x$; (b) simétrica: $x\equiv_k y\Rightarrow y\equiv_k x$; (c) transitiva: $x\equiv_k y$ y $y\equiv_k z$ implican $x\equiv_k z$.
\fuente{adaptado de Stanford CS103~\cite{cs103}, \emph{Problem Set}~1, problema~8.}
\end{ejercicio}

\begin{ejercicio}{3}{eje:3-real}
Sean $x,y\in\R$. Demuestre que, si $y^{3}+yx^{2}\le x^{3}+xy^{2}$, entonces $y\le x$.
\end{ejercicio}

\begin{ejercicio}{3}{eje:3-bicond}
Sea $n\in\Z$. Demuestre que $4\divides n^{2}$ si y solo si $n$ es par.
\end{ejercicio}

\begin{ejercicio}{2}{eje:3-restos}
Usando la \cref{prop:congruencias}, calcule sin calculadora: (a) el resto de dividir $3^{100}$ entre~7; (b) la última cifra de $7^{2026}$; (c) el resto de dividir $5^{40}+6^{40}$ entre~11.
\end{ejercicio}

\begin{ejercicio}{2}{eje:3-division}
(a) Calcule el cociente y el resto de la división euclídea de $2026$ entre~$7$ y de $-17$ entre~$5$. (b) Escriba la definición que hay que comprobar para decidir si $7\divides 2023$ y $7\divides 2026$, y decídalo. (c) Calcule el resto de $2026$ módulo~$9$ y compruebe que coincide con el resto de la suma de sus cifras.
\end{ejercicio}

% =====================================================================
\chapter{Reducción al absurdo y demostración por casos}
\label{tema:4}
% =====================================================================

Algunas afirmaciones no tienen la forma de un condicional, o la tienen pero ni la hipótesis ni la negación de la tesis dan un punto de partida útil. Para ellas se dispone de la reducción al absurdo: suponer lo contrario de lo que se quiere probar y llegar a una contradicción. Este tema la presenta con dos de los teoremas más antiguos de las matemáticas, la irracionalidad de $\sqrt2$ y la infinitud de los números primos, y estudia también la demostración por casos. Se siguen el capítulo~6 y las §4.4–4.5 de~\cite{hammack}, las §1.7–1.8 de~\cite{mcs}, la §3.5 de~\cite{velleman} y el primer \emph{problem set} de CS103~\cite{cs103}.

\section{Reducción al absurdo}

\begin{definicion}[Reducción al absurdo]\label{def:absurdo}
Para demostrar una proposición $P$ por \emph{reducción al absurdo} (o por contradicción) se supone $\neg P$ y se deduce una contradicción, es decir, una proposición de la forma $C\wedge\neg C$. Se concluye $P$.
\end{definicion}

La justificación es que $(\neg P\Rightarrow(C\wedge\neg C))\Rightarrow P$ es una tautología: si $P$ fuera falsa, $\neg P$ sería verdadera, y el condicional $\neg P\Rightarrow(C\wedge\neg C)$, con hipótesis verdadera y tesis falsa, sería falso. Para un enunciado condicional $P\Rightarrow Q$, la negación es $P\wedge\neg Q$ (\cref{prop:leyes}(vi)): se suponen a la vez la hipótesis y la negación de la tesis~\cite[§6.2]{hammack}.

\begin{ejemplo}\label{ej:no-mayor}
No existe un número entero mayor que todos los demás. Supóngase, por reducción al absurdo, que existe $M\in\Z$ con $n\le M$ para todo $n\in\Z$. Tomando $n=M+1\in\Z$ se obtiene $M+1\le M$, es decir, $1\le 0$, que es falso. Por tanto, tal $M$ no existe.
\end{ejemplo}

\begin{ejemplo}\label{ej:no-solucion}
No existen enteros $a$ y $b$ con $18a+6b=1$. Supóngase que sí existen. Entonces $1=18a+6b=2(9a+3b)$ sería par, mientras que $1=2\cdot 0+1$ es impar; ningún entero es par e impar a la vez (\cref{teo:division}). Contradicción.
\end{ejemplo}

\section{La irracionalidad de \texorpdfstring{$\sqrt2$}{raíz de 2}}

Un número real es \emph{racional} si se escribe como $a/b$ con $a,b\in\Z$ y $b\neq 0$, e \emph{irracional} en caso contrario. Todo racional puede escribirse en \emph{forma irreducible}, $a/b$ con $b>0$ y sin divisores comunes mayores que~1: basta dividir $a$ y $b$ por su máximo común divisor. (Alternativamente, se toma la representación con $b$ positivo mínimo, que existe por el principio del buen orden del \cref{tema:5}.)

\begin{teorema}\label{teo:raiz2}
$\sqrt2$ es irracional.
\end{teorema}

\begin{proof}
Supóngase, por reducción al absurdo, que $\sqrt2$ es racional, y escríbase $\sqrt2=a/b$ con $a,b\in\Z$, $b>0$ y la fracción irreducible. Elevando al cuadrado, $2=a^{2}/b^{2}$, es decir,
\[
a^{2}=2b^{2}.
\]
Así, $a^{2}$ es par, y por la \cref{prop:cuadrado-par} $a$ es par: $a=2c$ con $c\in\Z$. Sustituyendo, $4c^{2}=2b^{2}$, luego $b^{2}=2c^{2}$ es par, y de nuevo $b$ es par. Entonces $2$ divide a $a$ y a $b$, en contra de que $a/b$ sea irreducible. Por tanto, $\sqrt2$ no es racional.
\end{proof}

La demostración sigue la del teorema~1.8.1 de~\cite{mcs}. Su estructura se repite en los \cref{eje:4-raiz3,eje:4-raiz-cubica,eje:4-log}: suponer una fracción irreducible, obtener una ecuación entre enteros y encontrar un divisor común.

\section{La infinitud de los números primos}

\begin{lema}\label{lem:divisor-primo}
Todo entero $n\ge 2$ tiene al menos un divisor primo.
\end{lema}

\begin{proof}
El conjunto $D=\{d\in\Z: 2\le d\le n,\ d\divides n\}$ es finito y no vacío, pues $n\in D$. Sea $p$ su menor elemento. Si $p$ no fuera primo, tendría un divisor $e$ con $2\le e<p$; por la \cref{prop:divisibilidad}(i), $e\divides n$, luego $e\in D$ y $e<p$, en contra de la minimalidad de $p$. Por tanto, $p$ es un divisor primo de $n$.
\end{proof}

\begin{teorema}[Euclides]\label{teo:euclides}
Existen infinitos números primos.
\end{teorema}

\begin{proof}
Supóngase, por reducción al absurdo, que hay solo un número finito de primos, $p_1,p_2,\dots,p_k$. Sea
\[
N=p_1p_2\cdots p_k+1.
\]
Como $N\ge 2+1\ge 2$, el \cref{lem:divisor-primo} da un primo $p$ que divide a $N$, y por la hipótesis $p=p_i$ para algún $i$. Entonces $p_i$ divide a $N$ y a $p_1\cdots p_k$, y por la \cref{prop:divisibilidad}(ii) divide a su diferencia, $N-p_1\cdots p_k=1$. Pero ningún primo divide a~1, porque $p_i\ge 2$. Contradicción.
\end{proof}

\begin{observacion}
La prueba \emph{no} afirma que $N$ sea primo. Con los seis primeros primos, $N=2\cdot 3\cdot 5\cdot 7\cdot 11\cdot 13+1=30\,031=59\cdot 509$ es compuesto; lo que la prueba garantiza es que sus factores primos no están en la lista (\cref{eje:4-euclides-n}).
\end{observacion}

\section{Demostración por casos}

\begin{definicion}[Demostración por casos]\label{def:casos}
Para demostrar $Q$ a partir de $P_1\vee P_2\vee\dots\vee P_k$ se demuestran por separado $P_1\Rightarrow Q$, \dots, $P_k\Rightarrow Q$. Es válido porque $(P_1\vee\dots\vee P_k)\Rightarrow Q\eqlog(P_1\Rightarrow Q)\wedge\dots\wedge(P_k\Rightarrow Q)$, generalización del \cref{eje:1-equivalencias}(c).
\end{definicion}

Lo esencial es que los casos sean \emph{exhaustivos}: su disyunción debe ser verdadera. Las divisiones más habituales son «$n$ par o impar», los restos módulo $m$ y el signo de un número real~\cite[§4.4]{hammack}.

\begin{proposicion}\label{prop:consecutivos}
Para todo $n\in\Z$, el número $n^{2}+n$ es par.
\end{proposicion}

\begin{proof}
Todo entero es par o impar. \emph{Caso 1:} $n=2k$. Entonces $n^{2}+n=4k^{2}+2k=2(2k^{2}+k)$ es par. \emph{Caso 2:} $n=2k+1$. Entonces $n^{2}+n=(2k+1)(2k+2)=2(2k+1)(k+1)$ es par. En ambos casos $n^{2}+n$ es par.
\end{proof}

\begin{proposicion}\label{prop:cuadrado-mod4}
Para todo $n\in\Z$, $n^{2}\equiv 0\pmod 4$ o $n^{2}\equiv 1\pmod 4$.
\end{proposicion}

\begin{proof}
Si $n=2k$, $n^{2}=4k^{2}\equiv 0$. Si $n=2k+1$, $n^{2}=4(k^{2}+k)+1\equiv 1\pmod 4$.
\end{proof}

\begin{teorema}[Desigualdad triangular]\label{teo:triangular}
Para todos $x,y\in\R$, $\abs{x+y}\le\abs{x}+\abs{y}$.
\end{teorema}

\begin{proof}
Recuérdese que $\abs{t}=t$ si $t\ge 0$ y $\abs{t}=-t$ si $t<0$; en ambos casos $-\abs{t}\le t\le\abs{t}$. Se distinguen dos casos según el signo de $x+y$.
\emph{Caso 1:} $x+y\ge 0$. Entonces $\abs{x+y}=x+y\le\abs{x}+\abs{y}$.
\emph{Caso 2:} $x+y<0$. Entonces $\abs{x+y}=-x-y=(-x)+(-y)\le\abs{x}+\abs{y}$.
\end{proof}

Cuando dos casos se tratan con el mismo argumento intercambiando los papeles de las variables, se escribe «sin pérdida de generalidad, supóngase $x\le y$» y se trata solo uno~\cite[§4.5]{hammack}. Debe usarse solo cuando la simetría es real.

\begin{observacion}[Errores frecuentes]\label{obs:4-errores}
(a) El «sándwich de contradicción»: suponer $\neg P$, demostrar $P$ directamente sin usar la suposición y declarar una contradicción. Esa prueba es directa y la reducción al absurdo sobra~\cite{cs103}.
(b) Negar mal: la negación de «$\forall n$, $n$ es impar» no es «$\forall n$, $n$ es par» (\cref{eje:2-critica}).
(c) Casos no exhaustivos: dividir un real en «positivo» y «negativo» olvida el cero.
(d) En una reducción al absurdo, señalar con claridad qué se supone, cuál es la contradicción y qué se concluye.
\end{observacion}

\begin{observacion}[Conexión con el curso]
El módulo~4 demuestra por reducción al absurdo, con el teorema de Taylor, que en un mínimo local la hessiana es semidefinida positiva, y el módulo~6 usa un argumento por contradicción sobre los códigos prefijo. Las demostraciones por casos aparecen en cuanto intervienen funciones definidas a trozos: el valor absoluto en las cotas de error del módulo~17, la función ReLU, $\max(0,x)$, y su derivada en los módulos~9 y~11, y la desigualdad triangular en las normas del módulo~2.
\end{observacion}

\begin{lecturas}
\item Hammack~\cite{hammack}, cap.~6 «Proof by Contradiction» (§6.1–6.4) y §4.4 «Using Cases», §4.5 «Treating Similar Cases».
\item Lehman, Leighton y Meyer~\cite{mcs}, §1.7 «Proof by Cases» y §1.8 «Proof by Contradiction» (teorema~1.8.1).
\item Velleman~\cite{velleman}, §3.5 «Proofs Involving Disjunctions».
\item Stanford CS103~\cite{cs103}, clase~2 «Indirect Proofs».
\end{lecturas}

\section*{Ejercicios del tema 4}
\addcontentsline{toc}{section}{Ejercicios del tema 4}

\begin{ejercicio}{1}{eje:4-suponer}
Escriba qué se supone al comenzar una reducción al absurdo de: (a) «$\sqrt5$ es irracional»; (b) «para todo $n\in\N$, $n^{2}+1$ no es múltiplo de~3»; (c) «si $a$ y $b$ son enteros con $ab$ impar, entonces $a$ y $b$ son impares»; (d) «existe un primo mayor que $10^{100}$».
\end{ejercicio}

\begin{ejercicio}{3}{eje:4-raiz3}
Demuestre que, si $n\in\Z$ y $3\divides n^{2}$, entonces $3\divides n$ (por contrarrecíproco y casos). Deduzca que $\sqrt3$ es irracional.
\fuente{adaptado de Hammack~\cite{hammack}, cap.~6 «Proof by Contradiction», y de Lehman, Leighton y Meyer~\cite{mcs}, teorema~1.8.1.}
\end{ejercicio}

\begin{ejercicio}{3}{eje:4-raiz-cubica}
Demuestre que $\sqrt[3]{2}$ es irracional.
\end{ejercicio}

\begin{ejercicio}{3}{eje:4-log}
Demuestre que $\log_{2}3$ es irracional.
\fuente{adaptado de Lehman, Leighton y Meyer~\cite{mcs}, §1.8 «Proof by Contradiction» (problemas de la sección).}
\end{ejercicio}

\begin{ejercicio}{3}{eje:4-racional-irracional}
Sean $a$ un número racional y $b$ un número irracional. Demuestre que $a+b$ es irracional y que, si $a\neq 0$, también $ab$ es irracional. ¿Qué ocurre con $ab$ si $a=0$?
\end{ejercicio}

\begin{ejercicio}{3}{eje:4-diofantica}
Demuestre que no existen enteros $a$ y $b$ con $21a+30b=1$. ¿Existen enteros con $21a+30b=3$?
\fuente{adaptado de Hammack~\cite{hammack}, cap.~6 «Proof by Contradiction».}
\end{ejercicio}

\begin{ejercicio}{3}{eje:4-mod4}
Demuestre que, para todo $n\in\Z$, $n^{2}+2$ no es múltiplo de~4. Demuestre también que ningún entero de la forma $4k+3$ es suma de dos cuadrados de enteros.
\end{ejercicio}

\begin{ejercicio}{3}{eje:4-24}
Sea $p$ un número primo mayor que~3. Demuestre que $24\divides p^{2}-1$.
\end{ejercicio}

\begin{ejercicio}{2}{eje:4-maximo}
Demuestre por casos que, para todos $x,y\in\R$, $\max(x,y)=\dfrac{x+y+\abs{x-y}}{2}$ y $\min(x,y)=\dfrac{x+y-\abs{x-y}}{2}$. Compruebe las fórmulas con $(x,y)=(3,-5)$.
\end{ejercicio}

\begin{ejercicio}{2}{eje:4-restos-cuadrados}
Construya, por casos según el resto de $n$, las tablas de $n^{2}$ módulo~3, módulo~5 y módulo~8. Deduzca: (a) $n^{2}+1$ nunca es múltiplo de~3; (b) ningún cuadrado es congruente con $2$ ni con $3$ módulo~5; (c) el cuadrado de todo entero impar es congruente con~$1$ módulo~8.
\end{ejercicio}

\begin{ejercicio}{3}{eje:4-triangular-inversa}
Demuestre que, para todos $x,y\in\R$, $\bigl\lvert\abs{x}-\abs{y}\bigr\rvert\le\abs{x-y}$. (Indicación: aplique el \cref{teo:triangular} a $x=(x-y)+y$ y a $y=(y-x)+x$.)
\end{ejercicio}

\begin{ejercicio}{3}{eje:4-producto-par}
Sean $m,n\in\Z$. Demuestre por reducción al absurdo que, si $mn$ es par y $m$ es impar, entonces $n$ es par.
\fuente{adaptado de Stanford CS103~\cite{cs103}, \emph{Problem Set}~1, problema~6 «Writing Proofs by Contradiction».}
\end{ejercicio}

\begin{ejercicio}{3}{eje:4-cubos}
(a) Demuestre que todo múltiplo de~6 es suma de cuatro cubos de enteros. (b) Demuestre que $x^{3}\equiv x\pmod 6$ para todo $x\in\Z$. (c) Deduzca que todo entero es suma de cinco cubos de enteros y escriba $7$ de esa forma.
\fuente{adaptado de Stanford CS103~\cite{cs103}, \emph{Problem Set}~1, problema~10 «Sums of Cubes».}
\end{ejercicio}

\begin{ejercicio}{2}{eje:4-euclides-n}
Para $k=1,\dots,6$, sea $N_k=p_1\cdots p_k+1$, donde $p_1,\dots,p_k$ son los $k$ primeros primos. Calcule $N_k$, decida cuáles son primos y factorice $N_6$. Explique por qué el resultado no contradice la demostración del \cref{teo:euclides}.
\end{ejercicio}

\begin{ejercicio}{3}{eje:4-4k3}
Demuestre que existen infinitos primos de la forma $4k+3$. (Indicación: suponga que $p_1,\dots,p_m$ son todos, considere $N=4p_1\cdots p_m-1$ y demuestre que algún divisor primo de $N$ es de la forma $4k+3$.)
\end{ejercicio}

\begin{ejercicio}{3}{eje:4-raiz-primo}
Sea $p$ un número primo. Admitiendo el \emph{lema de Euclides} (si $p\divides ab$, entonces $p\divides a$ o $p\divides b$), demuestre que $\sqrt{p}$ es irracional. ¿En qué paso falla la demostración si se intenta con $\sqrt4$?
\end{ejercicio}

% =====================================================================
\chapter{Inducción simple y fuerte; buen orden}
\label{tema:5}
% =====================================================================

Muchos enunciados del curso dependen de un número natural: la longitud de una secuencia, el número de capas de una red, el paso de un algoritmo iterativo. La inducción matemática permite demostrarlos para todos los valores a la vez, y es la técnica de demostración más usada en informática. Este tema presenta la inducción simple, la inducción fuerte y el principio del buen orden, y muestra que son tres formas de la misma idea. Se siguen el capítulo~10 de~\cite{hammack}, los capítulos~2 y~5 de~\cite{mcs} y las §6.1, 6.2 y 6.4 de~\cite{velleman}.

\section{El principio de inducción}

\begin{teorema}[Principio de inducción]\label{teo:induccion}
Sea $P(n)$ un predicado sobre $\N_0$. Si
\begin{enumerate}
  \item (\emph{caso base}) $P(0)$ es verdadera, y
  \item (\emph{paso de inducción}) para todo $n\in\N_0$, $P(n)\Rightarrow P(n+1)$,
\end{enumerate}
entonces $P(n)$ es verdadera para todo $n\in\N_0$.
\end{teorema}

La imagen habitual es la de una fila infinita de fichas de dominó: la primera cae, y cada una que cae derriba a la siguiente. En el paso de inducción, la suposición $P(n)$ se llama \emph{hipótesis de inducción}; no se está suponiendo lo que se quiere demostrar, sino probando un condicional para un $n$ arbitrario. Si el caso base es $P(n_0)$ con $n_0\in\Z$, la conclusión vale para todo $n\ge n_0$. El \cref{teo:induccion} se toma como axioma de los números naturales o se deduce del principio del buen orden (\cref{prop:bo-induccion}).

\begin{ejemplo}[Suma de los primeros naturales]\label{ej:gauss}
Se demuestra que, para todo $n\in\N$, $1+2+\dots+n=\dfrac{n(n+1)}{2}$.

\emph{Demostración.} Sea $P(n)$ la igualdad del enunciado. \emph{Caso base:} para $n=1$, el miembro izquierdo vale~1 y el derecho $1\cdot 2/2=1$. \emph{Paso de inducción:} sea $n\ge 1$ y supóngase $P(n)$. Entonces
\[
1+\dots+n+(n+1)\overset{P(n)}{=}\frac{n(n+1)}{2}+(n+1)=\frac{(n+1)(n+2)}{2},
\]
que es $P(n+1)$. Por el principio de inducción, $P(n)$ vale para todo $n\in\N$. \hfill$\blacksquare$
\end{ejemplo}

\begin{ejemplo}[Serie geométrica finita]\label{ej:geometrica}
Para $r\in\R$, $r\neq 1$, y $n\in\N_0$:
\[
\sum_{k=0}^{n}r^{k}=\frac{r^{n+1}-1}{r-1}.
\]
\emph{Caso base:} $n=0$ da $1=(r-1)/(r-1)$. \emph{Paso:} si vale para $n$, entonces
$\sum_{k=0}^{n+1}r^{k}=\frac{r^{n+1}-1}{r-1}+r^{n+1}=\frac{r^{n+1}-1+r^{n+2}-r^{n+1}}{r-1}=\frac{r^{n+2}-1}{r-1}$. \hfill$\blacksquare$
\end{ejemplo}

\begin{ejemplo}[Una desigualdad]\label{ej:desig-induccion}
Para todo $n\in\N_0$, $2^{n}\ge n+1$. \emph{Caso base:} $2^{0}=1\ge 1$. \emph{Paso:} si $2^{n}\ge n+1$, entonces $2^{n+1}=2\cdot 2^{n}\ge 2(n+1)=n+2+n\ge n+2$. \hfill$\blacksquare$
\end{ejemplo}

En las desigualdades, el paso de inducción suele necesitar una estimación intermedia, como aquí $2n+2\ge n+2$. Conviene escribir primero qué hay que obtener, $P(n+1)$, y buscar cómo llegar desde $P(n)$.

\begin{ejemplo}[Divisibilidad]\label{ej:div-induccion}
Para todo $n\in\N_0$, $3\divides 4^{n}-1$. \emph{Caso base:} $4^{0}-1=0=3\cdot 0$. \emph{Paso:} si $4^{n}-1=3m$, entonces $4^{n+1}-1=4\cdot 4^{n}-1=4(3m+1)-1=3(4m+1)$. \hfill$\blacksquare$
\end{ejemplo}

\section{Inducción fuerte}

\begin{teorema}[Inducción fuerte]\label{teo:induccion-fuerte}
Sea $P(n)$ un predicado sobre los enteros $n\ge n_0$. Si $P(n_0)$ es verdadera y, para todo $n\ge n_0$,
\[
\bigl(P(n_0)\wedge P(n_0+1)\wedge\dots\wedge P(n)\bigr)\Rightarrow P(n+1),
\]
entonces $P(n)$ es verdadera para todo $n\ge n_0$.
\end{teorema}

\begin{proof}
Sea $Q(n)$ el predicado «$P(n_0),\dots,P(n)$ son verdaderas». $Q(n_0)$ es $P(n_0)$, verdadera. Si $Q(n)$ es verdadera, la hipótesis da $P(n+1)$, y por tanto $Q(n+1)$. Por inducción simple, $Q(n)$ vale para todo $n\ge n_0$, y en particular $P(n)$.
\end{proof}

La inducción fuerte es útil cuando $P(n+1)$ depende no de $P(n)$, sino de casos anteriores más lejanos, como al descomponer un número en factores o en las recurrencias de varios términos.

\begin{teorema}[Existencia de la factorización en primos]\label{teo:factorizacion}
Todo entero $n\ge 2$ es primo o producto de primos.
\end{teorema}

\begin{proof}
Por inducción fuerte sobre $n$. \emph{Caso base:} $2$ es primo. \emph{Paso:} sea $n\ge 2$ y supóngase el enunciado cierto para $2,3,\dots,n$. Si $n+1$ es primo, no hay nada que probar. Si no lo es, $n+1=ab$ con $2\le a,b\le n$. Por la hipótesis de inducción, $a$ y $b$ son primos o productos de primos, y entonces $n+1=ab$ es producto de primos.
\end{proof}

\begin{ejemplo}[Varios casos base]\label{ej:sellos}
Todo entero $n\ge 8$ se escribe como $n=3a+5b$ con $a,b\in\N_0$. \emph{Casos base:} $8=3+5$, $9=3\cdot 3$, $10=5\cdot 2$. \emph{Paso:} sea $n\ge 10$ y supóngase el enunciado para $8,\dots,n$. Entonces $n-2\ge 8$, luego $n-2=3a+5b$ y $n+1=3(a+1)+5b$. \hfill$\blacksquare$

Se necesitan tres casos base porque el paso usa el caso $n-2$, que debe ser al menos~8; con un solo caso base, el argumento no cubriría $9$ ni $10$.
\end{ejemplo}

\section{El principio del buen orden}

\begin{teorema}[Principio del buen orden]\label{teo:buen-orden}
Todo subconjunto no vacío de $\N_0$ tiene un elemento mínimo.
\end{teorema}

\begin{proof}
Sea $S\subseteq\N_0$ sin elemento mínimo; se demuestra que $S=\varnothing$. Sea $P(n)$: «ninguno de $0,1,\dots,n$ pertenece a $S$». $P(0)$ es cierta: si $0\in S$, sería su mínimo. Si $P(n)$ es cierta y $n+1\in S$, entonces $n+1$ sería el mínimo de $S$, porque ningún número menor está en $S$; luego $n+1\notin S$ y vale $P(n+1)$. Por inducción, $P(n)$ vale para todo $n$, y ningún natural está en $S$.
\end{proof}

\begin{proposicion}\label{prop:bo-induccion}
El principio del buen orden implica el principio de inducción.
\end{proposicion}

\begin{proof}
Sea $P$ con $P(0)$ y $P(n)\Rightarrow P(n+1)$ para todo $n$. Supóngase, por reducción al absurdo, que el conjunto $C=\{n\in\N_0:\neg P(n)\}$ de contraejemplos es no vacío, y sea $m$ su mínimo. Como $P(0)$ es cierta, $m\ge 1$, y $m-1\in\N_0$ no está en $C$ por ser menor que $m$; luego $P(m-1)$ es cierta. El paso de inducción da $P(m)$, en contra de $m\in C$.
\end{proof}

La demostración anterior es el modelo de las \emph{pruebas por el menor contraejemplo}~\cite[§2.2]{mcs}, \cite[§10.3]{hammack}: para demostrar $\forall n\,P(n)$, se supone que el conjunto de contraejemplos no es vacío, se toma el menor y se construye un contraejemplo aún menor, o se llega a otra contradicción.

\begin{proposicion}[Existencia en la división euclídea]\label{prop:division-existencia}
Sean $a\in\Z$ y $b\in\N$. Existen $q,r\in\Z$ con $a=qb+r$ y $0\le r<b$.
\end{proposicion}

\begin{proof}
Sea $S=\{a-qb: q\in\Z,\ a-qb\ge 0\}\subseteq\N_0$. Es no vacío: con $q=-\abs{a}$ se tiene $a+\abs{a}b\ge a+\abs{a}\ge 0$, porque $b\ge 1$. Por el buen orden tiene un mínimo $r=a-qb\ge 0$. Si fuera $r\ge b$, entonces $r-b=a-(q+1)b\ge 0$ estaría en $S$ y sería menor que $r$, en contra de la minimalidad. Luego $0\le r<b$.
\end{proof}

\section{Definiciones recursivas}

Los sumatorios, los productorios, las potencias y muchas sucesiones se definen \emph{recursivamente}: se da el primer valor y una regla para obtener cada valor a partir de los anteriores. Por ejemplo, $\sum_{k=1}^{0}a_k=0$ y $\sum_{k=1}^{n+1}a_k=\sum_{k=1}^{n}a_k+a_{n+1}$; la sucesión de Fibonacci se define por $F_1=F_2=1$ y $F_{n+2}=F_{n+1}+F_n$. La inducción es la forma natural de demostrar propiedades de objetos definidos así: el caso base comprueba la definición inicial y el paso usa la regla recursiva~\cite[§10.5]{hammack}.

\begin{observacion}[Errores frecuentes]\label{obs:5-errores}
(a) Omitir el caso base. Con $P(n)$: «$n=n+1$», el paso $P(n)\Rightarrow P(n+1)$ es correcto (sumando~1), pero $P(0)$ es falsa y el enunciado también.
(b) Un paso que no vale para todo $n$. La célebre «demostración» de que todos los caballos tienen el mismo color usa que dos grupos de $n$ caballos se solapan, lo que falla al pasar de $n=1$ a $n=2$ (\cref{eje:5-caballos}).
(c) Usar $P(n+1)$ dentro del paso: hay que llegar a $P(n+1)$, no partir de ella.
(d) En inducción fuerte, no incluir todos los casos base que el paso necesita (\cref{ej:sellos}).
\end{observacion}

\begin{observacion}[Conexión con el curso]
La inducción aparece en todo el curso. En el módulo~1, la regla de la cadena de la probabilidad se demuestra por inducción sobre la longitud de la secuencia, y también la normalización de un modelo de bigramas sin símbolo de fin. En el módulo~7, la expresión cerrada de las medias exponenciales de Adam se obtiene por inducción y se suma con la serie geométrica del \cref{ej:geometrica}. En el módulo~10, la corrección del modo directo de diferenciación automática se prueba por inducción sobre el orden topológico del grafo; en el módulo~12, la de BPE, por inducción sobre el número de fusiones; en el módulo~13, que la máscara causal impide que la salida en la posición $t$ dependa de posiciones posteriores; y en el módulo~17, la corrección del \emph{all-reduce} en anillo y de FlashAttention, por inducción sobre el número de pasos.
\end{observacion}

\begin{lecturas}
\item Hammack~\cite{hammack}, cap.~10 «Mathematical Induction» (§10.1–10.5).
\item Lehman, Leighton y Meyer~\cite{mcs}, cap.~2 «The Well Ordering Principle» (§2.1–2.3) y cap.~5 «Induction» (§5.1–5.3).
\item Velleman~\cite{velleman}, §6.1 «Proof by Mathematical Induction», §6.2 «More Examples» y §6.4 «Strong Induction».
\item MIT 6.042J~\cite{mit6042}, unidad~1, sesiones 3 (buen orden) y~8–9 (inducción).
\end{lecturas}

\section*{Ejercicios del tema 5}
\addcontentsline{toc}{section}{Ejercicios del tema 5}

\begin{ejercicio}{1}{eje:5-plantilla}
Para el enunciado «para todo $n\ge 1$, $1+3+5+\dots+(2n-1)=n^{2}$», escriba: (a) el predicado $P(n)$; (b) $P(1)$ y $P(n+1)$; (c) qué hay que demostrar exactamente en el paso de inducción. Después, complete la demostración.
\end{ejercicio}

\begin{ejercicio}{3}{eje:5-cuadrados}
Demuestre por inducción que, para todo $n\in\N$, $\displaystyle\sum_{k=1}^{n}k^{2}=\frac{n(n+1)(2n+1)}{6}$.
\fuente{adaptado de Hammack~\cite{hammack}, cap.~10 «Mathematical Induction».}
\end{ejercicio}

\begin{ejercicio}{3}{eje:5-cubos}
Demuestre por inducción que, para todo $n\in\N$, $\displaystyle\sum_{k=1}^{n}k^{3}=\Bigl(\frac{n(n+1)}{2}\Bigr)^{2}$.
\end{ejercicio}

\begin{ejercicio}{3}{eje:5-telescopica}
Demuestre por inducción que, para todo $n\in\N$, $\displaystyle\sum_{k=1}^{n}\frac{1}{k(k+1)}=\frac{n}{n+1}$. Dé después una segunda demostración sin inducción, usando $\frac{1}{k(k+1)}=\frac1k-\frac1{k+1}$.
\end{ejercicio}

\begin{ejercicio}{3}{eje:5-potencia}
Demuestre que $2^{n}>n^{2}$ para todo entero $n\ge 5$. ¿Para qué valores $n\in\{1,2,3,4\}$ es cierta la desigualdad?
\end{ejercicio}

\begin{ejercicio}{3}{eje:5-div9}
Demuestre que, para todo $n\in\N_0$, $9\divides 4^{n}+6n-1$.
\fuente{adaptado de Lehman, Leighton y Meyer~\cite{mcs}, §5.1 «Ordinary Induction» (problemas de la sección).}
\end{ejercicio}

\begin{ejercicio}{3}{eje:5-bernoulli}
Demuestre la desigualdad de Bernoulli: si $x\in\R$ y $x\ge -1$, entonces $(1+x)^{n}\ge 1+nx$ para todo $n\in\N_0$. Indique dónde se usa la hipótesis $x\ge -1$.
\end{ejercicio}

\begin{ejercicio}{3}{eje:5-sellos}
Demuestre por inducción fuerte que todo entero $n\ge 12$ se escribe como $n=4a+5b$ con $a,b\in\N_0$. ¿Qué enteros positivos menores que~12 no admiten esa forma?
\fuente{adaptado de Lehman, Leighton y Meyer~\cite{mcs}, §5.2 «Strong Induction».}
\end{ejercicio}

\begin{ejercicio}{3}{eje:5-fib-suma}
Sea $(F_n)$ la sucesión de Fibonacci ($F_1=F_2=1$, $F_{n+2}=F_{n+1}+F_n$). Demuestre que $\sum_{k=1}^{n}F_k=F_{n+2}-1$ para todo $n\in\N$.
\fuente{adaptado de Hammack~\cite{hammack}, §10.5 «Fibonacci Numbers».}
\end{ejercicio}

\begin{ejercicio}{3}{eje:5-fib-cota}
Demuestre por inducción fuerte que $F_n\le(5/3)^{n}$ para todo $n\in\N$. Explique por qué hacen falta dos casos base.
\fuente{adaptado de Hammack~\cite{hammack}, §10.5 «Fibonacci Numbers».}
\end{ejercicio}

\begin{ejercicio}{3}{eje:5-binario}
Demuestre por inducción fuerte que todo $n\in\N$ es suma de potencias de~2 distintas (es decir, tiene una representación binaria). (Indicación: distinga $n$ par y $n$ impar.)
\end{ejercicio}

\begin{ejercicio}{3}{eje:5-bo-geometrica}
Demuestre con el principio del buen orden (por el menor contraejemplo) que $\sum_{k=1}^{n}(2k-1)=n^{2}$ para todo $n\in\N$.
\fuente{adaptado de Lehman, Leighton y Meyer~\cite{mcs}, §2.2 «Template for Well Ordering Proofs».}
\end{ejercicio}

\begin{ejercicio}{2}{eje:5-caballos}
Se «demuestra» que, en todo conjunto de $n\ge 1$ caballos, todos tienen el mismo color: para $n=1$ es obvio; si vale para $n$, dado un conjunto de $n+1$ caballos se quita el primero y quedan $n$ del mismo color; se quita el último y quedan otros $n$ del mismo color; como ambos grupos comparten caballos, los $n+1$ tienen el mismo color. Señale con precisión el valor de $n$ para el que falla el paso y por qué.
\end{ejercicio}

\begin{ejercicio}{2}{eje:5-factorial}
(a) Calcule $S_n=\sum_{k=1}^{n}k\cdot k!$ para $n=1,\dots,5$ y conjeture una fórmula cerrada comparando con $(n+1)!$. (b) Demuestre la conjetura por inducción.
\end{ejercicio}

\begin{ejercicio}{2}{eje:5-sumas}
(a) Usando el \cref{ej:gauss}, obtenga una fórmula cerrada para $\sum_{k=1}^{n}(3k-1)$ y demuéstrela después por inducción. (b) Usando el \cref{ej:geometrica}, calcule $\sum_{k=0}^{7}2^{k}$ e interprete el resultado como el mayor número que cabe en un byte.
\end{ejercicio}

\begin{ejercicio}{3}{eje:5-log}
Sean $p_1,\dots,p_n$ números reales positivos. Demuestre por inducción sobre $n$, a partir de $\ln(ab)=\ln a+\ln b$, que $\ln\bigl(\prod_{t=1}^{n}p_t\bigr)=\sum_{t=1}^{n}\ln p_t$. Explique por qué esta identidad permite, en el módulo~1, convertir la probabilidad de una secuencia (un producto de muchos factores menores que~1) en una suma de log-probabilidades.
\end{ejercicio}

% =====================================================================
\chapter{Existencia, unicidad, contraejemplos y redacción}
\label{tema:6}
% =====================================================================

Los temas anteriores se han centrado en enunciados universales. Este tema trata los enunciados existenciales y de unicidad, explica cómo se refuta un enunciado falso y termina con unas pautas para redactar demostraciones claras, que resumen todo el módulo. Se siguen los capítulos~7 y~9 y la §5.3 de~\cite{hammack}, la §3.6 de~\cite{velleman}, la §1.9 de~\cite{mcs}, la lista de comprobación de CS103~\cite{cs103} y el método de Polya~\cite{polya}.

\section{Demostraciones de existencia}

Para demostrar $\exists x\,P(x)$ basta encontrar \emph{un} elemento $a$, llamado \emph{testigo}, y comprobar $P(a)$. Una demostración de este tipo es \emph{constructiva}. A veces se prueba la existencia sin exhibir el testigo, por ejemplo por reducción al absurdo o examinando casos sin saber cuál se da; la demostración es entonces \emph{no constructiva}~\cite[§7.3–7.4]{hammack}.

\begin{ejemplo}[Existencia constructiva]\label{ej:1729}
Existe un entero positivo que es suma de dos cubos positivos de dos maneras distintas: $1729=1^{3}+12^{3}=9^{3}+10^{3}$. La demostración consiste en dar el número y hacer las cuentas: $1+1728=1729$ y $729+1000=1729$.
\end{ejemplo}

\begin{proposicion}[Existencia no constructiva]\label{prop:irracional-potencia}
Existen números irracionales $a$ y $b$ tales que $a^{b}$ es racional.
\end{proposicion}

\begin{proof}
Sea $c=\sqrt2^{\sqrt2}$. Se distinguen dos casos. Si $c$ es racional, basta tomar $a=b=\sqrt2$, irracionales por el \cref{teo:raiz2}. Si $c$ es irracional, se toman $a=c$ y $b=\sqrt2$, y entonces
\[
a^{b}=\bigl(\sqrt2^{\sqrt2}\bigr)^{\sqrt2}=\sqrt2^{\sqrt2\cdot\sqrt2}=\sqrt2^{\,2}=2,
\]
que es racional. En ambos casos existen $a$ y $b$ con la propiedad.
\end{proof}

La prueba no dice cuál de los dos pares sirve. (Se sabe, por un teorema profundo de Gelfond y Schneider, que $c$ es irracional, pero la demostración anterior no lo necesita.)

\section{Demostraciones de unicidad}

Para demostrar $\exists!\,x\,P(x)$ se prueban dos cosas (\cref{def:unicidad}): la \emph{existencia}, $\exists x\,P(x)$, y la \emph{unicidad}, que se redacta así: «sean $x$ e $y$ tales que $P(x)$ y $P(y)$; entonces $x=y$». Una alternativa equivalente es suponer $x\neq y$ y llegar a una contradicción~\cite[§3.6]{velleman}.

\begin{proposicion}[Unicidad en la división euclídea]\label{prop:division-unicidad}
Sean $a\in\Z$ y $b\in\N$. Si $a=qb+r=q'b+r'$ con $q,q',r,r'\in\Z$ y $0\le r,r'<b$, entonces $q=q'$ y $r=r'$.
\end{proposicion}

\begin{proof}
Restando las dos expresiones, $(q-q')b=r'-r$. De $0\le r,r'<b$ se sigue $-b<r'-r<b$, es decir, $\abs{q-q'}\,b<b$, luego $\abs{q-q'}<1$. Como $q-q'$ es entero, $q=q'$, y entonces $r'-r=0$.
\end{proof}

Junto con la \cref{prop:division-existencia}, esto completa la demostración del \cref{teo:division}.

\begin{ejemplo}\label{ej:unico-real}
La ecuación $x^{3}+x-2=0$ tiene una única solución real. \emph{Existencia:} $x=1$ cumple $1+1-2=0$. \emph{Unicidad:} $x^{3}+x-2=(x-1)(x^{2}+x+2)$, y el segundo factor no se anula en $\R$ porque su discriminante es $1-8=-7<0$. Si $x^{3}+x-2=0$, el producto se anula y, como el segundo factor no es cero, $x-1=0$.
\end{ejemplo}

\section{Contraejemplos y refutaciones}

Para \emph{refutar} un enunciado se demuestra su negación~\cite[cap.~9]{hammack}. Según la forma del enunciado:
\begin{itemize}[leftmargin=1.4em]
  \item $\forall x\,P(x)$ se refuta con un \emph{contraejemplo}: un $a$ concreto con $\neg P(a)$.
  \item $\forall x\,(P(x)\Rightarrow Q(x))$ se refuta con un $a$ que cumple $P(a)$ y no cumple $Q(a)$.
  \item $\exists x\,P(x)$ se refuta demostrando $\forall x\,\neg P(x)$, lo que exige una demostración general, no un ejemplo.
\end{itemize}

\begin{ejemplo}[El polinomio de Euler]\label{ej:euler}
El polinomio $f(n)=n^{2}+n+41$ toma valores primos para $n=0,1,\dots,39$, lo que podría sugerir que es primo para todo $n\in\N_0$. El enunciado es falso: $f(40)=1600+40+41=1681=41^{2}$. Muchas comprobaciones favorables no demuestran un enunciado universal; un único contraejemplo lo refuta.
\end{ejemplo}

\begin{ejemplo}\label{ej:divide-producto}
«Si $a\divides bc$, entonces $a\divides b$ o $a\divides c$» es falso para enteros cualesquiera: con $a=6$, $b=4$, $c=9$ se tiene $6\divides 36$, pero $6\nodivide 4$ y $6\nodivide 9$. (Es cierto cuando $a$ es primo; es el lema de Euclides del \cref{eje:4-raiz-primo}.)
\end{ejemplo}

\section{Cómo redactar una demostración}

Una demostración es un texto dirigido a un lector, no una lista de cálculos. Las pautas siguientes resumen las de la §5.3 de~\cite{hammack}, la §1.9 de~\cite{mcs} y la lista de comprobación de CS103~\cite{cs103}.

\begin{enumerate}[label=\arabic*.]
  \item \emph{Anunciar el plan.} Indicar qué se supone, qué se quiere demostrar y con qué técnica («Se procede por contrarrecíproco», «Por inducción sobre $n$»).
  \item \emph{Introducir cada variable.} Toda letra debe presentarse antes de usarse, con su universo: «Sea $n$ un entero impar»; «existe $k\in\Z$ con\dots». No reutilizar una letra para objetos distintos.
  \item \emph{Usar las definiciones, no repetirlas.} Desplegar la definición en el punto donde se necesita, aplicada a los objetos concretos.
  \item \emph{Escribir frases completas.} Las fórmulas forman parte de las frases. Los símbolos $\forall$, $\exists$, $\Rightarrow$ son para enunciados formales, no abreviaturas en la prosa.
  \item \emph{Que cada frase aporte.} Eliminar lo que no se usa después; justificar cada paso no evidente citando el resultado empleado.
  \item \emph{Terminar con la conclusión.} Decir explícitamente que se ha demostrado lo que se buscaba y cerrar con $\blacksquare$.
\end{enumerate}

Para \emph{encontrar} la demostración, antes de redactarla, sirve el método de cuatro fases de Polya~\cite{polya}: comprender el problema (datos, incógnita, condiciones; escribir el enunciado con cuantificadores), concebir un plan (¿se conoce un problema relacionado?, ¿sirve un caso particular o el contrarrecíproco?), ejecutarlo comprobando cada paso, y examinar la solución obtenida (¿se usaron todas las hipótesis?, ¿hay un argumento más directo?).

\begin{ejemplo}[Borrador y redacción final]\label{ej:redaccion}
Enunciado: «si $a$ y $b$ son enteros impares, $a+b$ es par». Un borrador típico es «$a=2k+1$, $b=2k+1$, $a+b=4k+2=2(2k+1)$ par». Tiene dos defectos: usa la misma $k$ para dos números distintos y no presenta las variables. Redacción correcta: «Sean $a$ y $b$ enteros impares. Por definición existen $k,\ell\in\Z$ con $a=2k+1$ y $b=2\ell+1$. Entonces $a+b=2k+2\ell+2=2(k+\ell+1)$, y como $k+\ell+1$ es entero, $a+b$ es par. $\blacksquare$»
\end{ejemplo}

\begin{observacion}[Errores frecuentes]\label{obs:6-errores}
(a) «Demostrar» un enunciado universal con ejemplos (\cref{ej:euler}).
(b) Refutar un enunciado existencial con un ejemplo en que falla: que $P(5)$ sea falsa no impide que exista otro $x$ con $P(x)$.
(c) En una prueba de unicidad, suponer de entrada que los dos elementos son distintos y olvidar después usar esa suposición.
(d) Escribir una cadena de igualdades sin indicar cuál es el punto de partida y cuál el de llegada.
\end{observacion}

\begin{observacion}[Conexión con el curso]
Los resultados de existencia y unicidad son frecuentes en el curso: el sistema de ecuaciones normales de mínimos cuadrados tiene solución única cuando las columnas de la matriz de diseño son linealmente independientes (módulos~2 y~8), y una función estrictamente convexa tiene a lo sumo un mínimo (módulo~7). Los contraejemplos aparecen para marcar los límites de un resultado: el mínimo de dos funciones convexas no es convexo (módulo~7), el criterio de la hessiana no decide si hay un autovalor nulo (módulo~4) y la fórmula de la longitud tras una fusión de BPE falla cuando el par se solapa consigo mismo (módulo~12).
\end{observacion}

\begin{lecturas}
\item Hammack~\cite{hammack}, cap.~7 «Proving Non-Conditional Statements» (§7.3–7.4), cap.~9 «Disproof» y §5.3 «Mathematical Writing».
\item Velleman~\cite{velleman}, §3.6 «Existence and Uniqueness Proofs» y §3.7 «More Examples of Proofs».
\item Lehman, Leighton y Meyer~\cite{mcs}, §1.9 «Good Proofs in Practice».
\item Stanford CS103~\cite{cs103}, «Proofwriting Checklist».
\item Polya~\cite{polya}, primera parte (las cuatro fases de la resolución de problemas).
\end{lecturas}

\section*{Ejercicios del tema 6}
\addcontentsline{toc}{section}{Ejercicios del tema 6}

\begin{ejercicio}{1}{eje:6-clasificar}
Para cada enunciado, indique si es universal, existencial o de existencia y unicidad, y qué haría falta para demostrarlo y para refutarlo: (a) «todo primo mayor que~2 es impar»; (b) «hay un primo entre $100$ y $110$»; (c) «la ecuación $2x+3=7$ tiene una única solución real»; (d) «ningún entero cumple $n^{2}=3$».
\end{ejercicio}

\begin{ejercicio}{2}{eje:6-contraejemplos}
Refute con un contraejemplo: (a) para todo $x\in\R$, $x^{2}>x$; (b) para todos $a,b\in\R$, $\sqrt{a^{2}+b^{2}}=a+b$; (c) si $n\in\N$ y $n^{2}$ es múltiplo de~4, entonces $n$ es múltiplo de~4; (d) para todo $n\in\N$, $2^{n}+1$ es primo.
\fuente{adaptado de Hammack~\cite{hammack}, §9.1 «Counterexamples».}
\end{ejercicio}

\begin{ejercicio}{2}{eje:6-demostrar-refutar}
Demuestre o refute: (a) la suma de dos números irracionales es irracional; (b) el producto de dos números racionales es racional; (c) existe $n\in\N$ tal que $n^{2}+n+1$ es par; (d) para todo $x\in\R$ existe $y\in\R$ con $y^{2}=x$.
\fuente{adaptado de Hammack~\cite{hammack}, cap.~9 «Disproof» (§9.1–9.2).}
\end{ejercicio}

\begin{ejercicio}{3}{eje:6-testigos}
Demuestre: (a) existen números reales $a$ y $b$ con $\floor{a}\cdot\ceil{b}\neq\floor{ab}$; (b) existen naturales $a>b>c>d>0$ con $a^{2}+b^{2}+c^{2}+d^{2}=137$.
\fuente{adaptado de Stanford CS103~\cite{cs103}, \emph{Problem Set}~1, problema~7.}
\end{ejercicio}

\begin{ejercicio}{3}{eje:6-irracional-entre}
Sean $r<s$ números racionales. Demuestre que existe un número irracional $x$ con $r<x<s$.
\end{ejercicio}

\begin{ejercicio}{3}{eje:6-paridad-unica}
Demuestre que para cada $n\in\Z$ existe un único $k\in\Z$ tal que $n=2k$ o $n=2k+1$, y que ningún entero es a la vez par e impar. (Puede usar el \cref{teo:division}; para la segunda parte dé además una demostración directa por reducción al absurdo.)
\end{ejercicio}

\begin{ejercicio}{3}{eje:6-terna}
Demuestre que existe un único número primo $p$ tal que $p$, $p+2$ y $p+4$ son todos primos.
\end{ejercicio}

\begin{ejercicio}{3}{eje:6-mersenne}
Demuestre o refute: (a) si $n$ es primo, entonces $2^{n}-1$ es primo; (b) si $2^{n}-1$ es primo, entonces $n$ es primo (con $n\ge 2$). (Indicación para (b): por contrarrecíproco, use que $x^{m}-1=(x-1)(x^{m-1}+\dots+x+1)$.)
\end{ejercicio}

\begin{ejercicio}{2}{eje:6-promedio}
Sean $x_1,\dots,x_{10}$ números reales con $x_1+\dots+x_{10}=100$. Demuestre que existe $i$ con $x_i\ge 10$. ¿Es constructiva la demostración? ¿Existe necesariamente $i$ con $x_i>10$?
\end{ejercicio}

\begin{ejercicio}{2}{eje:6-critica-1}
Analice esta prueba de «la suma de un entero par y uno impar es impar»: «Sea $x=2k$ par e $y=2k+1$ impar. Entonces $x+y=4k+1$. Por tanto, $x+y$ es impar.» Identifique los defectos de redacción y de razonamiento, decida si el teorema es cierto y redacte una demostración correcta.
\fuente{adaptado de Stanford CS103~\cite{cs103}, \emph{Problem Set}~1, problema~9 «Proof Critiques».}
\end{ejercicio}

\begin{ejercicio}{2}{eje:6-critica-2}
Analice esta prueba de «para todo $n\in\Z$, si $n\equiv 0\pmod 2$, entonces $n\equiv 0\pmod 4$»: «Se procede por contrarrecíproco. Supóngase $n\equiv 0\pmod 4$. Entonces $n=4k=2(2k)$, luego $n\equiv 0\pmod 2$.» Indique qué se ha demostrado en realidad y si el enunciado es cierto.
\fuente{adaptado de Stanford CS103~\cite{cs103}, \emph{Problem Set}~1, problema~9 «Proof Critiques».}
\end{ejercicio}

\begin{ejercicio}{2}{eje:6-uno-dos}
Encuentre el error: «Sean $a=b$. Entonces $a^{2}=ab$, luego $a^{2}-b^{2}=ab-b^{2}$, es decir, $(a+b)(a-b)=b(a-b)$. Simplificando, $a+b=b$; como $a=b$, $2b=b$ y, dividiendo por $b$, $2=1$.»
\end{ejercicio}

\begin{ejercicio}{3}{eje:6-redactar}
El siguiente texto contiene las ideas de una demostración, pero no está redactado como tal: «$a\divides b$, $b\divides a$, $a,b>0$ $\Rightarrow$ $b=ak$, $a=bl$ $\Rightarrow$ $a=akl$ $\Rightarrow$ $kl=1$ $\Rightarrow$ $k=l=1$ $\Rightarrow$ $a=b$.» Redáctelo como una demostración completa de «si $a,b\in\N$, $a\divides b$ y $b\divides a$, entonces $a=b$», justificando especialmente el paso $kl=1\Rightarrow k=l=1$.
\fuente{adaptado de Hammack~\cite{hammack}, §5.3 «Mathematical Writing».}
\end{ejercicio}

\begin{ejercicio}{3}{eje:6-unica-raiz}
Sean $a,b,c\in\R$ con $a\neq 0$. Demuestre que la ecuación $ax^{2}+bx+c=0$ tiene una única solución real si y solo si $b^{2}=4ac$, y que en ese caso la solución es $x=-b/(2a)$.
\end{ejercicio}


\appendix
% =====================================================================
\chapter{Soluciones de los ejercicios}
\label{ap:soluciones}
% =====================================================================

Cada solución se presenta en cinco partes: datos, planteamiento, desarrollo, resultado y comprobación. En los ejercicios de demostración ($\star\star\star$), los pasos del desarrollo, leídos en orden, forman la demostración redactada tal como debería entregarse: cada paso es una frase completa que introduce sus variables y justifica lo que afirma. La comprobación indica una forma independiente de verificar el resultado o señala dónde se ha usado cada hipótesis.

% ---------------------------------------------------------------------
\section{Tema 1}

\begin{solucion}{eje:1-proposiciones}
\datos seis enunciados.
\planteamiento según la \cref{def:proposicion}, un enunciado es proposición si es verdadero o falso, sin depender de variables libres.
\pasos
\begin{enumerate}[label=(\alph*)]
  \item $2^{10}=1024>1000$: proposición \textbf{verdadera}.
  \item Contiene la variable libre $x$: es un predicado, \textbf{no} una proposición (aunque $\forall x\in\R\ (x^{2}\ge 0)$ sí lo sería, y verdadera).
  \item $91=7\cdot 13$: proposición \textbf{falsa}.
  \item Es una orden: \textbf{no} es proposición.
  \item Proposición \textbf{falsa}, porque $x^{2}\ge 0>-1$ para todo real $x$.
  \item Proposición \textbf{verdadera} (\cref{teo:division} con $b=2$).
\end{enumerate}
\resultado proposiciones: (a) V, (c) F, (e) F, (f) V; no son proposiciones (b) y (d).
\comprobacion (b) se convierte en proposición al fijar $x$ o al cuantificarla, como se verá en el \cref{tema:2}.
\end{solucion}

\begin{solucion}{eje:1-traducir}
\datos $P$: $n$ par; $Q$: $3\divides n$; $R$: $6\divides n$.
\planteamiento «$A$ solo si $B$» es $A\Rightarrow B$; «$A$ no es suficiente para $B$» es $\neg(A\Rightarrow B)$.
\pasos
\begin{enumerate}[label=(\alph*)]
  \item $R\Leftrightarrow(P\wedge Q)$. Es verdadera para todo $n$: $6\divides n$ si y solo si $2\divides n$ y $3\divides n$.
  \item $P\Rightarrow R$. Falsa en general: $n=2$ es par y no es múltiplo de~6.
  \item $\neg(P\Rightarrow R)\eqlog P\wedge\neg R$. No es verdadera para todo $n$ (falla con $n=6$); lo que es cierto es que existe un $n$ que la cumple ($n=2$).
  \item $\neg P\wedge\neg Q$, o equivalentemente $\neg(P\vee Q)$ por De Morgan.
\end{enumerate}
\resultado (a) $R\Leftrightarrow(P\wedge Q)$, siempre verdadera; (b) $P\Rightarrow R$, no siempre; (c) $P\wedge\neg R$, no siempre; (d) $\neg P\wedge\neg Q$.
\comprobacion con $n=2$: $P=\V$, $Q=\F$, $R=\F$; (a) da $\F\Leftrightarrow\F=\V$, (b) da $\V\Rightarrow\F=\F$, (c) da $\V$.
\end{solucion}

\begin{solucion}{eje:1-reciproco}
\datos $P$: «$f$ es derivable en $a$»; $Q$: «$f$ es continua en $a$». Se sabe (bachillerato) que $P\Rightarrow Q$ es verdadero.
\planteamiento aplicar la \cref{def:reciproco} y la \cref{prop:leyes}(v).
\pasos
\begin{enumerate}[label=\arabic*.]
  \item Recíproco $Q\Rightarrow P$: «si $f$ es continua en $a$, es derivable en $a$». Falso: $f(x)=\abs{x}$ es continua en $0$ y no es derivable allí (las derivadas laterales valen $-1$ y $1$).
  \item Contrarrecíproco $\neg Q\Rightarrow\neg P$: «si $f$ no es continua en $a$, no es derivable en $a$». Verdadero, por ser equivalente al enunciado original.
  \item Inverso $\neg P\Rightarrow\neg Q$: «si $f$ no es derivable en $a$, no es continua en $a$». Es el contrarrecíproco del recíproco, luego equivalente a él, y también falso con $f(x)=\abs{x}$.
\end{enumerate}
\resultado \textbf{contrarrecíproco verdadero; recíproco e inverso falsos} (contraejemplo $f(x)=\abs{x}$, $a=0$).
\comprobacion el recíproco y el inverso siempre tienen el mismo valor de verdad, pues $Q\Rightarrow P\eqlog\neg P\Rightarrow\neg Q$.
\end{solucion}

\begin{solucion}{eje:1-necesaria}
\datos $T$: «$4\divides n$».
\planteamiento $C$ es suficiente si $C\Rightarrow T$; necesaria si $T\Rightarrow C$. Se busca demostración o contraejemplo para cada implicación.
\pasos
\begin{enumerate}[label=(\alph*)]
  \item $T\Rightarrow$ par: si $n=4k$, $n=2(2k)$. Par $\Rightarrow T$ falla con $n=2$. \textbf{Necesaria, no suficiente.}
  \item $8\divides n\Rightarrow T$: $n=8k=4(2k)$. $T\Rightarrow 8\divides n$ falla con $n=4$. \textbf{Suficiente, no necesaria.}
  \item $T\Rightarrow 16\divides n^{2}$: $n=4k$ da $n^{2}=16k^{2}$. Recíproco: si $n$ es impar, $n^{2}$ es impar; si $n=2m$ con $m$ impar, $n^{2}=4m^{2}$ con $m^{2}$ impar, que no es múltiplo de~16 (sería $4m^{2}=16j$, $m^{2}=4j$ par). Por contrarrecíproco, $16\divides n^{2}\Rightarrow n=2m$ con $m$ par $\Rightarrow 4\divides n$. \textbf{Necesaria y suficiente.}
  \item «Termina en~4» no implica $T$ ($14$) y $T$ no la implica ($8$). \textbf{Ninguna.}
\end{enumerate}
\resultado (a) necesaria; (b) suficiente; (c) necesaria y suficiente; (d) ninguna.
\comprobacion valores: $n=12$ cumple $T$, es par, $144=16\cdot 9$, no termina en~4 y no es múltiplo de~8, coherente con todo lo anterior.
\end{solucion}

\begin{solucion}{eje:1-silogismo}
\datos $\Phi=\bigl((P\Rightarrow Q)\wedge(Q\Rightarrow R)\bigr)\Rightarrow(P\Rightarrow R)$.
\planteamiento construir las $2^{3}=8$ filas (\cref{prop:filas}).
\pasos
\begin{enumerate}[label=\arabic*.]
  \item Tabla, con $H=(P\Rightarrow Q)\wedge(Q\Rightarrow R)$:
  \begin{center}\small
  \begin{tabular}{@{}ccc|ccccc@{}}
  \toprule
  $P$&$Q$&$R$&$P\Rightarrow Q$&$Q\Rightarrow R$&$H$&$P\Rightarrow R$&$\Phi$\\
  \midrule
  V&V&V&V&V&V&V&V\\
  V&V&F&V&F&F&F&V\\
  V&F&V&F&V&F&V&V\\
  V&F&F&F&V&F&F&V\\
  F&V&V&V&V&V&V&V\\
  F&V&F&V&F&F&V&V\\
  F&F&V&V&V&V&V&V\\
  F&F&F&V&V&V&V&V\\
  \bottomrule
  \end{tabular}
  \end{center}
  \item La última columna es siempre $\V$: $\Phi$ es una tautología.
\end{enumerate}
\resultado \textbf{$\Phi$ es una tautología}; por definición, el silogismo hipotético es un razonamiento válido.
\comprobacion sin tabla: $\Phi$ solo podría ser falsa con $P\Rightarrow R$ falsa, es decir, $P=\V$, $R=\F$; entonces $Q\Rightarrow R$ verdadera exige $Q=\F$, y $P\Rightarrow Q$ sería $\V\Rightarrow\F=\F$, luego $H=\F$ y $\Phi=\V$.
\end{solucion}

\begin{solucion}{eje:1-xor}
\datos definición de $\xor$.
\planteamiento comparar columnas de las tablas.
\pasos
\begin{enumerate}[label=\arabic*.]
  \item Tabla (filas $PQ=$ VV, VF, FV, FF):
  \begin{center}\small
  \begin{tabular}{@{}cc|ccccc@{}}
  \toprule
  $P$&$Q$&$P\xor Q$&$P\vee Q$&$\neg(P\wedge Q)$&$(P\vee Q)\wedge\neg(P\wedge Q)$&$\neg(P\Leftrightarrow Q)$\\
  \midrule
  V&V&F&V&F&F&F\\
  V&F&V&V&V&V&V\\
  F&V&V&V&V&V&V\\
  F&F&F&F&V&F&F\\
  \bottomrule
  \end{tabular}
  \end{center}
  \item Las columnas 3, 6 y 7 coinciden fila a fila.
\end{enumerate}
\resultado (a) $P\xor Q$ vale $\F,\V,\V,\F$; (b) y (c) son \textbf{equivalencias ciertas}.
\comprobacion intuitivamente, «exactamente una» es «al menos una» y «no las dos», y también «no tienen el mismo valor».
\end{solucion}

\begin{solucion}{eje:1-equivalencias}
\datos tres pares de fórmulas en $P,Q,R$.
\planteamiento basta estudiar las filas en que alguna fórmula es falsa.
\pasos
\begin{enumerate}[label=(\alph*)]
  \item $P\Rightarrow(Q\Rightarrow R)$ es falsa solo si $P=\V$ y $Q\Rightarrow R$ falsa, es decir, $P=Q=\V$, $R=\F$. $(P\wedge Q)\Rightarrow R$ es falsa solo si $P\wedge Q=\V$ y $R=\F$: la misma fila. \textbf{Equivalentes.}
  \item $P\Rightarrow(Q\Rightarrow R)$ es falsa solo en $(\V,\V,\F)$. $(P\Rightarrow Q)\Rightarrow R$ es falsa cuando $R=\F$ y $P\Rightarrow Q=\V$, es decir, en $(\V,\V,\F)$, $(\F,\V,\F)$ y $(\F,\F,\F)$. \textbf{No equivalentes}: en $P=Q=R=\F$, la primera es $\V$ y la segunda $\F$.
  \item $(P\vee Q)\Rightarrow R$ es falsa si $R=\F$ y $P\vee Q=\V$: filas $(\V,\V,\F)$, $(\V,\F,\F)$, $(\F,\V,\F)$. $(P\Rightarrow R)\wedge(Q\Rightarrow R)$ es falsa si alguna de las dos lo es: $R=\F$ con $P=\V$ o con $Q=\V$, las mismas tres filas. \textbf{Equivalentes.}
\end{enumerate}
\resultado (a) sí; (b) no (asignación $P=Q=R=\F$, también $P=\F,Q=\V,R=\F$); (c) sí.
\comprobacion (a) es la «exportación»: demostrar «si $P$ y $Q$, entonces $R$» equivale a suponer $P$ y luego demostrar «si $Q$, entonces $R$». (c) justifica las demostraciones por casos (\cref{def:casos}).
\end{solucion}

\begin{solucion}{eje:1-algebra}
\datos leyes (i)–(vii) de la \cref{prop:leyes}.
\planteamiento transformar paso a paso indicando la ley usada.
\pasos
\begin{enumerate}[label=\arabic*.]
  \item Distributiva: $P\vee(\neg P\wedge Q)\eqlog(P\vee\neg P)\wedge(P\vee Q)$.
  \item $P\vee\neg P$ es tautología, y $\V\wedge A\eqlog A$: queda $P\vee Q$.
  \item Luego $\neg\bigl(P\vee(\neg P\wedge Q)\bigr)\eqlog\neg(P\vee Q)\eqlog\neg P\wedge\neg Q$ por De Morgan.
  \item Por (iv): $(P\Rightarrow Q)\wedge(P\Rightarrow\neg Q)\eqlog(\neg P\vee Q)\wedge(\neg P\vee\neg Q)$.
  \item Distributiva (sacando $\neg P$): $\eqlog\neg P\vee(Q\wedge\neg Q)$.
  \item $Q\wedge\neg Q$ es contradicción y $A\vee\F\eqlog A$: queda $\neg P$.
\end{enumerate}
\resultado \textbf{ambas equivalencias se cumplen}.
\comprobacion la segunda dice que, si $P$ implicara a la vez $Q$ y $\neg Q$, $P$ sería falsa: es el fundamento de la reducción al absurdo (\cref{tema:4}). Por tabla: con $P=\V$ el miembro izquierdo exige $Q$ y $\neg Q$, luego vale $\F=\neg P$; con $P=\F$ ambos condicionales son verdaderos y $\neg P=\V$.
\end{solucion}

\begin{solucion}{eje:1-validez}
\datos cuatro razonamientos.
\planteamiento identificar la forma lógica y compararla con la \cref{prop:inferencia} y la \cref{obs:1-errores}.
\pasos
\begin{enumerate}[label=(\alph*)]
  \item $P\Rightarrow Q$, $Q$; luego $P$. Es la falacia de \textbf{afirmar el consecuente}: el error de validación puede subir por otras causas. \textbf{No válido.}
  \item $P\Rightarrow Q$ con $P$: «$n$ primo mayor que~2», $Q$: «$n$ impar»; premisa $\neg Q$; conclusión $\neg P$. \textbf{Modus tollens, válido.}
  \item $P\vee Q$, $\neg P$; luego $Q$. \textbf{Silogismo disyuntivo, válido.}
  \item $P\Rightarrow Q$, $\neg P$; luego $\neg Q$. Falacia de \textbf{negar el antecedente}; contraejemplo: $x=-3\le 2$ y $x^{2}=9>4$. \textbf{No válido.}
\end{enumerate}
\resultado válidos (b) y (c); no válidos (a) y (d).
\comprobacion en (a) y (d), la asignación $P=\F$, $Q=\V$ hace verdaderas las premisas y falsa la conclusión.
\end{solucion}

\begin{solucion}{eje:1-contar}
\datos $n$ variables.
\planteamiento una tabla asigna un valor a cada una de las $2^{n}$ filas (\cref{prop:filas}).
\pasos
\begin{enumerate}[label=\arabic*.]
  \item Cada una de las $2^{n}$ filas recibe $\V$ o $\F$ de forma independiente: $2^{2^{n}}$ tablas posibles.
  \item $n=1$: $2^{2}=4$; $n=2$: $2^{4}=16$; $n=3$: $2^{8}=256$.
  \item Para $n=2$, solo una tabla es toda $\V$ (tautología) y solo una es toda $\F$ (contradicción).
  \item Las cuatro funciones de una variable: constante $\V$ ($P\vee\neg P$), constante $\F$ ($P\wedge\neg P$), identidad ($P$) y negación ($\neg P$).
\end{enumerate}
\resultado (a) $2^{2^{n}}$: \textbf{4, 16, 256}; (b) \textbf{una tautología y una contradicción}; (c) $P\vee\neg P$, $P\wedge\neg P$, $P$, $\neg P$.
\comprobacion la \cref{tab:conectivas} y el \cref{eje:1-xor} muestran ya 6 de las 16 funciones de dos variables: $\wedge$, $\vee$, $\Rightarrow$, $\Leftrightarrow$, $\xor$ y la proyección $P$.
\end{solucion}

\begin{solucion}{eje:1-leyes}
\datos leyes (iii), (vi) y (vii).
\planteamiento tablas de verdad; para (iii) basta razonar por casos según $P$.
\pasos
\begin{enumerate}[label=\arabic*.]
  \item (iii), primera ley. Si $P=\F$, ambos miembros valen $\F$. Si $P=\V$, el izquierdo vale $Q\vee R$ y el derecho $(\V\wedge Q)\vee(\V\wedge R)=Q\vee R$. Coinciden en las 8 filas.
  \item (iii), segunda ley. Si $P=\V$, ambos valen $\V$. Si $P=\F$, el izquierdo vale $Q\wedge R$ y el derecho $(\F\vee Q)\wedge(\F\vee R)=Q\wedge R$.
  \item (vi). $\neg(P\Rightarrow Q)$ es verdadera solo en la fila $P=\V,Q=\F$ (\cref{tab:conectivas}); $P\wedge\neg Q$ también. En las demás filas ambas son falsas.
  \item (vii). $(P\Rightarrow Q)\wedge(Q\Rightarrow P)$: en VV vale $\V\wedge\V=\V$; en VF, $\F\wedge\V=\F$; en FV, $\V\wedge\F=\F$; en FF, $\V\wedge\V=\V$. Es la columna de $P\Leftrightarrow Q$.
\end{enumerate}
\resultado las tres leyes \textbf{se cumplen} en todas las asignaciones.
\comprobacion (vi) se deduce también de (iv) y De Morgan: $\neg(\neg P\vee Q)\eqlog\neg\neg P\wedge\neg Q\eqlog P\wedge\neg Q$.
\end{solucion}

\begin{solucion}{eje:1-demorgan-n}
\datos $n\ge 2$ proposiciones $P_1,\dots,P_n$.
\planteamiento dos fórmulas son equivalentes si son verdaderas exactamente en las mismas asignaciones; se describe ese conjunto para cada miembro.
\pasos
\begin{enumerate}[label=\arabic*.]
  \item Sea una asignación cualquiera. Por definición de la conjunción (aplicada repetidamente), $P_1\wedge\dots\wedge P_n$ es verdadera si y solo si todas las $P_i$ son verdaderas.
  \item Por tanto, $\neg(P_1\wedge\dots\wedge P_n)$ es verdadera si y solo si \emph{no} todas las $P_i$ son verdaderas, es decir, si y solo si existe un índice $i$ con $P_i$ falsa.
  \item Por otra parte, $\neg P_1\vee\dots\vee\neg P_n$ es verdadera si y solo si alguna $\neg P_i$ es verdadera, es decir, si y solo si existe $i$ con $P_i$ falsa.
  \item Las dos condiciones de los pasos 2 y 3 son la misma, luego ambas fórmulas tienen el mismo valor en cada asignación: son equivalentes.
  \item Para la segunda ley, se aplica la primera a las proposiciones $\neg P_1,\dots,\neg P_n$: $\neg(\neg P_1\wedge\dots\wedge\neg P_n)\eqlog\neg\neg P_1\vee\dots\vee\neg\neg P_n\eqlog P_1\vee\dots\vee P_n$. Negando ambos miembros y usando la doble negación se obtiene $\neg P_1\wedge\dots\wedge\neg P_n\eqlog\neg(P_1\vee\dots\vee P_n)$. $\blacksquare$
\end{enumerate}
\resultado \textbf{las leyes de De Morgan valen para cualquier número finito de proposiciones}.
\comprobacion con $n=2$ se recupera la \cref{prop:leyes}(ii). Una prueba alternativa es por inducción sobre $n$ (\cref{tema:5}), escribiendo $P_1\wedge\dots\wedge P_{n+1}=(P_1\wedge\dots\wedge P_n)\wedge P_{n+1}$.
\end{solucion}

\begin{solucion}{eje:1-completitud}
\datos una tabla de verdad $T$ con variables $P_1,\dots,P_n$.
\planteamiento construir la forma normal disyuntiva y eliminar $\vee$ con De Morgan.
\pasos
\begin{enumerate}[label=\arabic*.]
  \item Para cada asignación $a=(a_1,\dots,a_n)$ se define el \emph{literal} $L_i^{a}=P_i$ si $a_i=\V$ y $L_i^{a}=\neg P_i$ si $a_i=\F$, y la conjunción $C_a=L_1^{a}\wedge\dots\wedge L_n^{a}$. Por construcción, $C_a$ es verdadera en la asignación $a$ y falsa en todas las demás, pues en cualquier otra asignación algún literal es falso.
  \item Sea $A$ el conjunto de asignaciones en las que $T$ vale $\V$. Si $A=\varnothing$, la fórmula $P_1\wedge\neg P_1$ tiene la tabla $T$. Si $A\neq\varnothing$, sea $\Phi=\bigvee_{a\in A}C_a$.
  \item $\Phi$ tiene la tabla $T$: en una asignación $b\in A$, el término $C_b$ es verdadero y $\Phi$ también; en una asignación $b\notin A$, todos los $C_a$ con $a\in A$ son falsos, porque $a\neq b$, y $\Phi$ es falsa.
  \item Se elimina $\vee$: por De Morgan generalizado (\cref{eje:1-demorgan-n}) y la doble negación, $\bigvee_{a\in A}C_a\eqlog\neg\bigwedge_{a\in A}\neg C_a$, que solo contiene $\neg$ y $\wedge$.
  \item Por tanto, toda tabla de verdad es la de una fórmula con solo $\neg$ y $\wedge$. $\blacksquare$
\end{enumerate}
\resultado \textbf{$\{\neg,\wedge\}$ es un conjunto completo de conectivas}.
\comprobacion para $T=P\xor Q$, $A=\{(\V,\F),(\F,\V)\}$ y $\Phi=(P\wedge\neg Q)\vee(\neg P\wedge Q)$, que coincide con la tabla del \cref{eje:1-xor}; su forma con $\neg,\wedge$ es $\neg\bigl(\neg(P\wedge\neg Q)\wedge\neg(\neg P\wedge Q)\bigr)$.
\end{solucion}

\begin{solucion}{eje:1-bicond-asoc}
\datos $L=(P\Leftrightarrow Q)\Leftrightarrow R$, $D=P\Leftrightarrow(Q\Leftrightarrow R)$.
\planteamiento caracterizar cuándo es verdadero un bicondicional en términos del número de falsas.
\pasos
\begin{enumerate}[label=\arabic*.]
  \item \emph{Lema.} $A\Leftrightarrow B$ es verdadero si y solo si el número de falsas entre $A$ y $B$ es par (0 o 2). En efecto, es verdadero cuando ambas tienen el mismo valor, es decir, con 0 falsas (ambas $\V$) o 2 falsas (ambas $\F$).
  \item Sea $f$ el número de falsas entre $P,Q,R$. Si $R=\V$, $L$ tiene el valor de $P\Leftrightarrow Q$, verdadero si y solo si el número de falsas entre $P,Q$ es par, es decir (como $R$ no aporta falsas), si y solo si $f$ es par.
  \item Si $R=\F$, $L$ es verdadera si y solo si $P\Leftrightarrow Q$ es falsa, es decir, si el número de falsas entre $P,Q$ es impar; sumando la falsa $R$, si y solo si $f$ es par.
  \item El mismo argumento, separando según el valor de $P$, prueba que $D$ es verdadera si y solo si $f$ es par.
  \item Ambas fórmulas son verdaderas exactamente en las asignaciones con $f$ par; luego $L\eqlog D$. $\blacksquare$
  \item El condicional no es asociativo: con $P=Q=R=\F$, $(P\Rightarrow Q)\Rightarrow R=\V\Rightarrow\F=\F$, mientras que $P\Rightarrow(Q\Rightarrow R)=\F\Rightarrow\V=\V$.
\end{enumerate}
\resultado \textbf{$\Leftrightarrow$ es asociativo; $\Rightarrow$ no lo es} (asignación $P=Q=R=\F$).
\comprobacion fila $P=\V$, $Q=\F$, $R=\F$ ($f=2$, par): $L=(\F)\Leftrightarrow\F=\V$ y $D=\V\Leftrightarrow(\F\Leftrightarrow\F)=\V\Leftrightarrow\V=\V$. Se han comprobado además las 8 filas por ordenador.
\end{solucion}

% ---------------------------------------------------------------------
\section{Tema 2}

\begin{solucion}{eje:2-valores}
\datos seis enunciados cuantificados; $\N=\{1,2,\dots\}$.
\planteamiento un universal se prueba en general y se refuta con un contraejemplo; un existencial se prueba con un testigo.
\pasos
\begin{enumerate}[label=(\alph*)]
  \item \textbf{V}: el cuadrado de un real es no negativo.
  \item \textbf{F}: $\sqrt2\notin\Q$ (\cref{teo:raiz2}), y las únicas soluciones reales son $\pm\sqrt2$.
  \item \textbf{V}: testigo $x=\sqrt2$.
  \item \textbf{F}: contraejemplo $n=1$, pues no hay $m\in\N$ con $m<1$.
  \item \textbf{V}: testigo $m=1$, que cumple $1\le n$ para todo $n\in\N$.
  \item \textbf{V}: universal sobre el conjunto vacío (vacuamente verdadero).
\end{enumerate}
\resultado V, F, V, F, V, V.
\comprobacion (d) y (e) muestran que en $\N$ hay mínimo; con $\Z$ en lugar de $\N$ los valores se invertirían.
\end{solucion}

\begin{solucion}{eje:2-simbolos}
\datos cuatro enunciados.
\planteamiento identificar universo, cuantificadores y forma (universal condicional, existencial conjuntiva).
\pasos
\begin{enumerate}[label=(\alph*)]
  \item $\forall n\in\Z\ \bigl((n\text{ par}\wedge n>2)\Rightarrow\exists p\,\exists q\,(p,q\text{ primos}\wedge n=p+q)\bigr)$.
  \item $\exists x\in\R\ \forall n\in\Z\ (x>n)$. \textbf{Falso}: dado $x$, el entero $n=\ceil{x}$ cumple $n\ge x$, luego no es cierto que $x>n$.
  \item $\forall n\in\Z\ \bigl(n^{2}\not\equiv 2\pmod{10}\bigr)$, o $\neg\exists n\in\Z\ (n^{2}\equiv 2\pmod{10})$.
  \item $\neg\exists x\in\R\ (x^{2}+1=0)$, equivalente a $\forall x\in\R\ (x^{2}+1\neq 0)$.
\end{enumerate}
\resultado las fórmulas anteriores; (b) es falso.
\comprobacion (c) es verdadero: los cuadrados de $0,\dots,9$ terminan en $0,1,4,9,6,5,6,9,4,1$, y la última cifra de $n^{2}$ solo depende de la de $n$ (\cref{prop:congruencias}).
\end{solucion}

\begin{solucion}{eje:2-lenguaje}
\datos cuatro fórmulas.
\planteamiento leer de izquierda a derecha, teniendo presente que una variable cuantificada después puede depender de las anteriores.
\pasos
\begin{enumerate}[label=(\alph*)]
  \item «Para todo real hay otro mayor». \textbf{V}: $y=x+1$.
  \item «Hay un real mayor que todos los reales». \textbf{F}: tomando $x=y$ sería $y>y$.
  \item «Todo real no nulo tiene inverso». \textbf{V}: $y=1/x$.
  \item «Hay un entero que, multiplicado por cualquier entero, lo deja igual». \textbf{V}: $x=1$.
\end{enumerate}
\resultado V, F, V, V.
\comprobacion (a) y (b) son el par $\forall\exists$/$\exists\forall$ de la \cref{prop:orden} con $P(x,y)$: «$y>x$».
\end{solucion}

\begin{solucion}{eje:2-negar}
\datos cuatro enunciados.
\planteamiento aplicar el \cref{teo:negacion}, De Morgan y $\neg(A\Rightarrow B)\eqlog A\wedge\neg B$.
\pasos
\begin{enumerate}[label=(\alph*)]
  \item Negación: $\exists x\in\R\ \forall y\in\R\ (y\le x\vee y\ge x+1)$. El original es \textbf{verdadero}: $y=x+\tfrac12$.
  \item Negación: $\forall n\in\N\ \exists m\in\N\ (n\nodivide m)$. El original es \textbf{verdadero}: $n=1$ divide a todo $m$.
  \item Negación: $\exists x\in\R\ (x^{2}>4\wedge x\le 2)$. La \textbf{negación es verdadera}: $x=-3$.
  \item Negación: $\exists n\in\Z\ (n\text{ primo}\wedge n\text{ par}\wedge n\neq 2)$. El original es \textbf{verdadero}: un primo par es divisible por~2, luego igual a~2.
\end{enumerate}
\resultado verdaderos: original en (a), (b), (d); negación en (c).
\comprobacion en (d), $\neg(A\vee B)\eqlog\neg A\wedge\neg B$ convierte «impar o igual a~2» en «par y distinto de~2».
\end{solucion}

\begin{solucion}{eje:2-continuidad}
\datos definición $\varepsilon$–$\delta$; $f(x)=0$ si $x<0$, $f(x)=1$ si $x\ge 0$; $a=0$.
\planteamiento negar como en el \cref{ej:negar-limite} y exhibir $\varepsilon$ y, para cada $\delta$, un $x$.
\pasos
\begin{enumerate}[label=\arabic*.]
  \item Negación: $\exists\varepsilon>0\ \forall\delta>0\ \exists x\in\R\ \bigl(\abs{x-a}<\delta\wedge\abs{f(x)-f(a)}\ge\varepsilon\bigr)$.
  \item Se toma $\varepsilon=\tfrac12$. Sea $\delta>0$ arbitrario y sea $x=-\delta/2$.
  \item Entonces $\abs{x-0}=\delta/2<\delta$ y, como $x<0$, $\abs{f(x)-f(0)}=\abs{0-1}=1\ge\tfrac12$.
  \item Como $\delta$ era arbitrario, se cumple la negación: $f$ no es continua en $0$. $\blacksquare$
\end{enumerate}
\resultado negación del paso~1; \textbf{$f$ no es continua en $0$} (con $\varepsilon=\tfrac12$ y $x=-\delta/2$).
\comprobacion cualquier $\varepsilon\in(0,1]$ sirve; con $\varepsilon>1$ no, porque $\abs{f(x)-f(0)}\le 1$.
\end{solucion}

\begin{solucion}{eje:2-orden}
\datos cuatro predicados y universos.
\planteamiento para $\forall x\exists y$, dar $y$ en función de $x$; para $\exists y\forall x$, un $y$ fijo o refutarlo.
\pasos
\begin{enumerate}[label=(\alph*)]
  \item $\N_0$, $x+y=0$. $\forall\exists$: \textbf{F} ($x=1$ exigiría $y=-1\notin\N_0$). $\exists\forall$: \textbf{F} (implica el anterior, \cref{prop:orden}).
  \item $\Q\setminus\{0\}$, $xy=1$. $\forall\exists$: \textbf{V} ($y=1/x$). $\exists\forall$: \textbf{F} ($x=1$ da $y=1$ y $x=2$ da $y=\tfrac12$).
  \item $\N$, $x\le y$. $\forall\exists$: \textbf{V} ($y=x$). $\exists\forall$: \textbf{F} (sería un máximo de $\N$; $x=y+1$ lo refuta).
  \item $\N$, $y\le x$. $\forall\exists$: \textbf{V} ($y=1$). $\exists\forall$: \textbf{V} ($y=1$, el mínimo de $\N$).
\end{enumerate}
\resultado (a) F, F; (b) V, F; (c) V, F; (d) V, V.
\comprobacion en ningún caso se tiene $\exists\forall$ verdadero y $\forall\exists$ falso, como exige la \cref{prop:orden}.
\end{solucion}

\begin{solucion}{eje:2-distribuir}
\datos cuatro posibles equivalencias.
\planteamiento para las falsas, universo $\Z$, $P(x)$: «$x$ par», $Q(x)$: «$x$ impar».
\pasos
\begin{enumerate}[label=(\alph*)]
  \item \textbf{Cierta} (demostración en el \cref{eje:2-distribuir-dem}).
  \item \textbf{Falsa}: $\forall x\,(P(x)\vee Q(x))$ es verdadera (todo entero es par o impar), pero $\forall x\,P(x)$ y $\forall x\,Q(x)$ son falsas.
  \item \textbf{Cierta} (\cref{eje:2-distribuir-dem}).
  \item \textbf{Falsa}: $\exists x\,P(x)$ y $\exists x\,Q(x)$ son verdaderas (testigos 2 y 3), pero $\exists x\,(P(x)\wedge Q(x))$ es falsa.
\end{enumerate}
\resultado $\forall$ distribuye sobre $\wedge$ y $\exists$ sobre $\vee$; no al revés.
\comprobacion (b) y (d) son negación una de la otra: aplicando el \cref{teo:negacion} a (b) con $\neg P$, $\neg Q$ se obtiene (d), luego una falla si y solo si falla la otra.
\end{solucion}

\begin{solucion}{eje:2-distribuir-dem}
\datos predicados $P,Q$ sobre un universo $U$.
\planteamiento probar cada bicondicional por sus dos implicaciones, de forma directa.
\pasos
\begin{enumerate}[label=\arabic*.]
  \item (a, $\Rightarrow$) Supóngase $\forall x\,(P(x)\wedge Q(x))$. Sea $a\in U$ arbitrario. Entonces $P(a)\wedge Q(a)$, en particular $P(a)$; como $a$ es arbitrario, $\forall x\,P(x)$. Del mismo modo, $\forall x\,Q(x)$. Luego $\forall x\,P(x)\wedge\forall x\,Q(x)$.
  \item (a, $\Leftarrow$) Supóngase $\forall x\,P(x)$ y $\forall x\,Q(x)$. Sea $a\in U$ arbitrario. Por la primera, $P(a)$; por la segunda, $Q(a)$; luego $P(a)\wedge Q(a)$. Como $a$ es arbitrario, $\forall x\,(P(x)\wedge Q(x))$.
  \item (c, $\Rightarrow$) Supóngase $\exists x\,(P(x)\vee Q(x))$ y sea $a$ un testigo: $P(a)\vee Q(a)$. Si $P(a)$, entonces $\exists x\,P(x)$; si $Q(a)$, entonces $\exists x\,Q(x)$. En ambos casos, $\exists x\,P(x)\vee\exists x\,Q(x)$.
  \item (c, $\Leftarrow$) Supóngase $\exists x\,P(x)\vee\exists x\,Q(x)$. Si $\exists x\,P(x)$ con testigo $a$, entonces $P(a)\vee Q(a)$ y $a$ es testigo de $\exists x\,(P(x)\vee Q(x))$. El otro caso es análogo.
  \item (Implicación adicional) Supóngase $\forall x\,P(x)\vee\forall x\,Q(x)$ y sea $a\in U$. Si vale $\forall x\,P(x)$, entonces $P(a)$, luego $P(a)\vee Q(a)$; en el otro caso, $Q(a)$ y también $P(a)\vee Q(a)$. Como $a$ es arbitrario, $\forall x\,(P(x)\vee Q(x))$. $\blacksquare$
\end{enumerate}
\resultado \textbf{(a) y (c) son equivalencias; $\forall x\,P\vee\forall x\,Q\Rightarrow\forall x\,(P\vee Q)$ siempre}, aunque su recíproca falla (\cref{eje:2-distribuir}(b)).
\comprobacion los pasos 3, 4 y 5 son demostraciones por casos (\cref{def:casos}); la palabra «arbitrario» marca dónde se demuestra un universal.
\end{solucion}

\begin{solucion}{eje:2-vacio}
\datos $U\neq\varnothing$; predicado $P$.
\planteamiento demostración directa, usando un elemento de $U$ como testigo.
\pasos
\begin{enumerate}[label=\arabic*.]
  \item Supóngase $\forall x\in U\ P(x)$. Como $U$ es no vacío, existe algún $a\in U$.
  \item Por la hipótesis aplicada a $a$, $P(a)$ es verdadera.
  \item Luego $a$ es testigo de $\exists x\in U\ P(x)$. $\blacksquare$
  \item Si $U=\varnothing$, $\forall x\,P(x)$ es vacuamente verdadero y $\exists x\,P(x)$ es falso; el condicional sería $\V\Rightarrow\F$, falso. El paso~1 es el que falla: no hay ningún $a$.
\end{enumerate}
\resultado \textbf{la implicación vale si $U\neq\varnothing$} y falla para $U=\varnothing$.
\comprobacion la \cref{def:predicado} exige universos no vacíos precisamente para evitar este caso.
\end{solucion}

\begin{solucion}{eje:2-unicidad}
\datos \cref{def:unicidad}.
\planteamiento separar existencia y unicidad; negar con el \cref{teo:negacion}.
\pasos
\begin{enumerate}[label=(\alph*)]
  \item $\exists x\,P(x)\ \wedge\ \forall x\,\forall y\,\bigl((P(x)\wedge P(y))\Rightarrow x=y\bigr)$.
  \item Negación: $\neg\exists x\,P(x)\vee\exists x\,\exists y\,(P(x)\wedge P(y)\wedge x\neq y)$, es decir, $\forall x\,\neg P(x)\ \vee\ \exists x\,\exists y\,(P(x)\wedge P(y)\wedge x\neq y)$. Puede fallar porque no haya ningún elemento con la propiedad o porque haya al menos dos distintos.
  \item En $\R$, $x^{2}=4$ tiene dos soluciones, $2$ y $-2$: \textbf{falso}. En $\N$, solo $x=2$: \textbf{verdadero}.
\end{enumerate}
\resultado las fórmulas de (a) y (b); (c) falso en $\R$ y verdadero en $\N$.
\comprobacion la fórmula de (a) es equivalente a la de la \cref{def:unicidad}: si $a$ cumple $P(a)\wedge\forall y\,(P(y)\Rightarrow y=a)$, dos elementos $x,y$ con la propiedad son ambos iguales a $a$.
\end{solucion}

\begin{solucion}{eje:2-acotada}
\datos definiciones de acotada y creciente; $f(x)=x^{2}$.
\planteamiento negar y dar testigos.
\pasos
\begin{enumerate}[label=\arabic*.]
  \item No acotada: $\forall M\in\R\ \exists x\in\R\ (\abs{f(x)}>M)$.
  \item No creciente: $\exists x\,\exists y\ (x<y\wedge f(x)>f(y))$.
  \item $f$ no es acotada: sea $M\in\R$ arbitrario y $x=\abs{M}+1$. Entonces $x\ge 1$, luego $x^{2}\ge x=\abs{M}+1>M$.
  \item $f$ no es creciente: $x=-2<y=0$ y $f(-2)=4>0=f(0)$. $\blacksquare$
\end{enumerate}
\resultado negaciones de los pasos 1 y 2; $x^{2}$ \textbf{no es acotada ni creciente}.
\comprobacion la negación de «no acotada» requiere un $x$ \emph{para cada} $M$: el $x$ del paso~3 depende de $M$, como debe ser.
\end{solucion}

\begin{solucion}{eje:2-arquimedes}
\datos propiedad arquimediana.
\planteamiento (a) dado $\varepsilon$, elegir $N>1/\varepsilon$; (b) negar y exhibir, para cada $N$, un $\varepsilon$ y un $n$.
\pasos
\begin{enumerate}[label=\arabic*.]
  \item (a) Sea $\varepsilon>0$. Por la propiedad arquimediana aplicada a $x=1/\varepsilon$, existe $N\in\N$ con $N>1/\varepsilon$. Sea $n\in\N$ con $n\ge N$. Entonces $n>1/\varepsilon$, y como $n>0$ y $\varepsilon>0$, $1/n<\varepsilon$. Esto prueba (a).
  \item (b) La negación es $\forall N\in\N\ \exists\varepsilon>0\ \exists n\in\N\ (n\ge N\wedge 1/n\ge\varepsilon)$. Sea $N\in\N$ arbitrario. Se toman $n=N$ y $\varepsilon=1/N>0$. Entonces $n\ge N$ y $1/n=\varepsilon\ge\varepsilon$.
  \item Se ha probado la negación de (b), luego (b) es falso. $\blacksquare$
\end{enumerate}
\resultado \textbf{(a) verdadero; (b) falso}. En (a), $N$ depende de $\varepsilon$; (b) pediría un $N$ válido para todas las tolerancias.
\comprobacion con $\varepsilon=0{,}01$, (a) da $N=101$, y $1/101\approx 0{,}0099<0{,}01$.
\end{solucion}

\begin{solucion}{eje:2-critica}
\datos «todo natural es impar» y la prueba propuesta.
\planteamiento aplicar el \cref{teo:negacion}.
\pasos
\begin{enumerate}[label=\arabic*.]
  \item La negación de $\forall n\in\N\ (n\text{ impar})$ es $\exists n\in\N\ (n\text{ no impar})$, es decir, «algún natural es par», no «todos los naturales son pares».
  \item La prueba refuta un enunciado («todos son pares») que no es la negación del original; la contradicción obtenida no dice nada sobre él.
  \item El enunciado es \textbf{falso}: $2\in\N$ es par.
\end{enumerate}
\resultado error de \textbf{negación} (\cref{obs:2-errores}(a)); el teorema es falso, contraejemplo $n=2$.
\comprobacion el mismo «argumento» probaría, cambiando 137 por 2, que todo natural es par: señal de que no prueba nada.
\end{solucion}

\begin{solucion}{eje:2-finito}
\datos $U=\{a_1,\dots,a_n\}$.
\planteamiento traducir la \cref{def:cuantificadores} y usar el \cref{eje:1-demorgan-n}.
\pasos
\begin{enumerate}[label=\arabic*.]
  \item $\forall x\,P(x)$ es verdadera si y solo si $P(a)$ lo es para cada $a\in U$, es decir, si y solo si $P(a_1),\dots,P(a_n)$ son todas verdaderas, que es la condición de verdad de $P(a_1)\wedge\dots\wedge P(a_n)$.
  \item $\exists x\,P(x)$ es verdadera si y solo si alguna $P(a_i)$ lo es, condición de verdad de $P(a_1)\vee\dots\vee P(a_n)$.
  \item Entonces $\neg\forall x\,P(x)\eqlog\neg\bigl(P(a_1)\wedge\dots\wedge P(a_n)\bigr)\eqlog\neg P(a_1)\vee\dots\vee\neg P(a_n)\eqlog\exists x\,\neg P(x)$, usando los pasos 1 y 2 (este último con el predicado $\neg P$) y De Morgan generalizado.
  \item La segunda ley se obtiene igual con la otra ley de De Morgan. $\blacksquare$
\end{enumerate}
\resultado en universos finitos, $\forall$ es una conjunción, $\exists$ una disyunción, y el \cref{teo:negacion} es De Morgan.
\comprobacion la demostración del \cref{teo:negacion} en el texto no usa la finitud, y vale para universos infinitos como $\N$ o $\R$.
\end{solucion}

% ---------------------------------------------------------------------
\section{Tema 3}

\begin{solucion}{eje:3-hipotesis}
\datos cuatro condicionales.
\planteamiento separar hipótesis y tesis, aplicar la \cref{def:reciproco} y valorar qué punto de partida es más manejable.
\pasos
\begin{enumerate}[label=(\alph*)]
  \item Hipótesis: $a\divides b$ y $b\divides c$; tesis: $a\divides c$. Contrarrecíproco: si $a\nodivide c$, entonces $a\nodivide b$ o $b\nodivide c$. \textbf{Directa}: las hipótesis dan igualdades $b=ad$, $c=be$ (\cref{prop:divisibilidad}).
  \item Hipótesis: $n^{2}$ impar; tesis: $n$ impar. Contrarrecíproco: si $n$ es par, $n^{2}$ es par. \textbf{Contrarrecíproco}: partir de $n=2k$ es inmediato; partir de $n^{2}=2k+1$ no.
  \item Hipótesis: $x^{2}+x<0$; tesis: $x<0$. Contrarrecíproco: si $x\ge 0$, entonces $x^{2}+x\ge 0$. \textbf{Contrarrecíproco}: suma de no negativos.
  \item Hipótesis: $xy$ irracional; tesis: $x$ o $y$ irracional. Contrarrecíproco: si $x$ e $y$ son racionales, $xy$ es racional. \textbf{Contrarrecíproco}: «irracional» es una negación y su contrario, «racional», da una fracción con la que operar.
\end{enumerate}
\resultado directa en (a); contrarrecíproco en (b), (c) y (d).
\comprobacion en (d) la negación de «$x$ o $y$ irracional» es «$x$ e $y$ racionales» por De Morgan.
\end{solucion}

\begin{solucion}{eje:3-esquema}
\datos la prueba propuesta.
\planteamiento comparar con la \cref{def:directa}.
\pasos
\begin{enumerate}[label=\arabic*.]
  \item Defecto 1: parte de la tesis ($n^{2}=4k$) en lugar de la hipótesis ($n$ par): demuestra, si acaso, el recíproco.
  \item Defecto 2: $2\sqrt{k}$ no es necesariamente un entero, de modo que «$n=2\sqrt k$ es par» no está justificado.
  \item Redacción correcta: «Sea $n$ un entero par. Por definición existe $k\in\Z$ con $n=2k$. Entonces $n^{2}=4k^{2}$, y como $k^{2}\in\Z$, $n^{2}$ es múltiplo de~4. $\blacksquare$»
\end{enumerate}
\resultado error de dirección (se supone la tesis) y paso no justificado; prueba correcta en el paso~3.
\comprobacion con $n=6$: $n^{2}=36=4\cdot 9$.
\end{solucion}

\begin{solucion}{eje:3-cubo}
\datos $x\in\Z$ impar.
\planteamiento demostración directa: desplegar la definición y llegar a la forma $2m+1$.
\pasos
\begin{enumerate}[label=\arabic*.]
  \item Sea $x$ un entero impar. Por la \cref{def:paridad}, existe $k\in\Z$ con $x=2k+1$.
  \item Entonces $x^{3}=(2k+1)^{3}=8k^{3}+12k^{2}+6k+1=2(4k^{3}+6k^{2}+3k)+1$.
  \item Sea $m=4k^{3}+6k^{2}+3k$, que es entero por ser suma de productos de enteros. Entonces $x^{3}=2m+1$, luego $x^{3}$ es impar. $\blacksquare$
\end{enumerate}
\resultado \textbf{si $x$ es impar, $x^{3}$ es impar.}
\comprobacion $x=3$: $k=1$, $m=4+6+3=13$, y $2\cdot 13+1=27=3^{3}$.
\end{solucion}

\begin{solucion}{eje:3-polinomio}
\datos $a\in\Z$ impar.
\planteamiento directa, como en el \cref{eje:3-cubo}.
\pasos
\begin{enumerate}[label=\arabic*.]
  \item Sea $a$ impar; existe $k\in\Z$ con $a=2k+1$.
  \item $a^{2}+3a+5=(4k^{2}+4k+1)+(6k+3)+5=4k^{2}+10k+9=2(2k^{2}+5k+4)+1$.
  \item Con $m=2k^{2}+5k+4\in\Z$, $a^{2}+3a+5=2m+1$ es impar. $\blacksquare$
\end{enumerate}
\resultado \textbf{$a^{2}+3a+5$ es impar si $a$ es impar.}
\comprobacion $a=1$: $1+3+5=9$ impar, y la fórmula da $k=0$, $m=4$, $2\cdot 4+1=9$. Otra vía: $a^{2}$ y $3a$ son impares, su suma es par, y par más~5 es impar.
\end{solucion}

\begin{solucion}{eje:3-cuadrados}
\datos $a,b\in\Z$ con $a\divides b$.
\planteamiento directa con la \cref{def:divide}.
\pasos
\begin{enumerate}[label=\arabic*.]
  \item Supóngase $a\divides b$. Existe $c\in\Z$ con $b=ac$.
  \item Elevando al cuadrado, $b^{2}=a^{2}c^{2}$.
  \item Como $c^{2}\in\Z$, la igualdad muestra que $a^{2}\divides b^{2}$. $\blacksquare$
\end{enumerate}
\resultado \textbf{$a\divides b\Rightarrow a^{2}\divides b^{2}$.}
\comprobacion $3\divides 12$ y $9\divides 144=9\cdot 16$. El recíproco también es cierto, pero su demostración requiere la factorización única en primos.
\end{solucion}

\begin{solucion}{eje:3-producto}
\datos $a\divides b$, $c\divides d$.
\planteamiento directa; usar letras distintas para los dos cocientes.
\pasos
\begin{enumerate}[label=\arabic*.]
  \item Supóngase $a\divides b$ y $c\divides d$. Existen $k,\ell\in\Z$ con $b=ak$ y $d=c\ell$.
  \item Multiplicando, $bd=(ak)(c\ell)=(ac)(k\ell)$.
  \item Como $k\ell\in\Z$, se concluye $ac\divides bd$. $\blacksquare$
\end{enumerate}
\resultado \textbf{$ac\divides bd$.}
\comprobacion $2\divides 6$ y $5\divides 15$ dan $10\divides 90$; en efecto, $90=10\cdot 9$ y $k\ell=3\cdot 3=9$.
\end{solucion}

\begin{solucion}{eje:3-desigualdad}
\datos $x\in\R$, $0<x<4$.
\planteamiento buscar hacia atrás: $\frac{4}{x(4-x)}\ge 1\iff 4\ge x(4-x)$ (por ser $x(4-x)>0$) $\iff x^{2}-4x+4\ge 0\iff(x-2)^{2}\ge 0$. Después se redacta hacia delante.
\pasos
\begin{enumerate}[label=\arabic*.]
  \item Sea $x\in\R$ con $0<x<4$. Como el cuadrado de un real es no negativo, $0\le(x-2)^{2}=x^{2}-4x+4$.
  \item Sumando $4x-x^{2}$ a ambos miembros, $x(4-x)=4x-x^{2}\le 4$.
  \item De $0<x$ y $x<4$ se obtiene $x>0$ y $4-x>0$, luego $x(4-x)>0$.
  \item Dividiendo la desigualdad del paso~2 por el número positivo $x(4-x)$, se obtiene $1\le\dfrac{4}{x(4-x)}$. $\blacksquare$
\end{enumerate}
\resultado \textbf{$\dfrac{4}{x(4-x)}\ge 1$}, con igualdad solo si $x=2$.
\comprobacion $x=1$: $4/3\ge 1$; $x=2$: $4/4=1$. La hipótesis $0<x<4$ se usa en el paso~3: sin ella el denominador podría ser negativo o nulo.
\end{solucion}

\begin{solucion}{eje:3-diferencia}
\datos $n\in\Z$ impar.
\planteamiento buscar $a,b$ con $a^{2}-b^{2}=(a-b)(a+b)=n$; con $a-b=1$ basta $a+b=n$.
\pasos
\begin{enumerate}[label=\arabic*.]
  \item Sea $n$ un entero impar; existe $k\in\Z$ con $n=2k+1$.
  \item Sean $a=k+1$ y $b=k$, que son enteros.
  \item Entonces $a^{2}-b^{2}=(k+1)^{2}-k^{2}=k^{2}+2k+1-k^{2}=2k+1=n$.
  \item Por tanto, $n$ es diferencia de los cuadrados de dos enteros. $\blacksquare$
\end{enumerate}
\resultado \textbf{$2k+1=(k+1)^{2}-k^{2}$.}
\comprobacion $n=9$: $k=4$, $25-16=9$; $n=-3$: $k=-2$, $(-1)^{2}-(-2)^{2}=1-4=-3$.
\end{solucion}

\begin{solucion}{eje:3-contra-cubo}
\datos $n\in\Z$.
\planteamiento contrarrecíproco: si $n$ es impar, $n^{3}+5$ es par.
\pasos
\begin{enumerate}[label=\arabic*.]
  \item Se demuestra el contrarrecíproco. Sea $n$ un entero no par, es decir, impar.
  \item Por el \cref{eje:3-cubo}, $n^{3}$ es impar: $n^{3}=2m+1$ con $m\in\Z$.
  \item Entonces $n^{3}+5=2m+6=2(m+3)$, que es par, es decir, no es impar.
  \item Se ha probado «si $n$ no es par, $n^{3}+5$ no es impar», que equivale al enunciado. $\blacksquare$
\end{enumerate}
\resultado \textbf{si $n^{3}+5$ es impar, $n$ es par.}
\comprobacion $n=2$: $13$ impar y $n$ par; $n=1$: $6$ par, coherente con el contrarrecíproco.
\end{solucion}

\begin{solucion}{eje:3-pitagoras}
\datos $a,b,c\in\Z$ con $a^{2}+b^{2}=c^{2}$.
\planteamiento contrarrecíproco: si $a,b,c$ son todos impares, $a^{2}+b^{2}\neq c^{2}$.
\pasos
\begin{enumerate}[label=\arabic*.]
  \item La negación de «al menos uno es par» es «los tres son impares». Se demuestra: si $a$, $b$ y $c$ son impares, entonces $a^{2}+b^{2}\neq c^{2}$.
  \item Sean $a,b,c$ impares. Por la \cref{prop:impar-cuadrado}, $a^{2}$, $b^{2}$ y $c^{2}$ son impares: existen $p,q\in\Z$ con $a^{2}=2p+1$ y $b^{2}=2q+1$.
  \item Entonces $a^{2}+b^{2}=2(p+q+1)$ es par, mientras que $c^{2}$ es impar.
  \item Como ningún entero es par e impar a la vez, $a^{2}+b^{2}\neq c^{2}$. $\blacksquare$
\end{enumerate}
\resultado \textbf{en toda terna pitagórica entera hay al menos un número par.}
\comprobacion $(3,4,5)$: $4$ es par; $(5,12,13)$: $12$ es par. No hace falta desplegar «$c^{2}$ es impar»: basta citar la \cref{prop:impar-cuadrado}.
\end{solucion}

\begin{solucion}{eje:3-equivalencia}
\datos $x\equiv_k y\iff\exists q\in\Z\ (x=y+kq)$.
\planteamiento en cada apartado, exhibir el entero $q$ que pide la definición.
\pasos
\begin{enumerate}[label=(\alph*)]
  \item Sea $x\in\Z$. Con $q=0$, $x=x+k\cdot 0$; luego $x\equiv_k x$.
  \item Supóngase $x\equiv_k y$: existe $q\in\Z$ con $x=y+kq$. Entonces $y=x+k(-q)$ y $-q\in\Z$, luego $y\equiv_k x$.
  \item Supóngase $x\equiv_k y$ y $y\equiv_k z$: existen $q,r\in\Z$ con $x=y+kq$ e $y=z+kr$. Sustituyendo, $x=z+kr+kq=z+k(q+r)$ con $q+r\in\Z$, luego $x\equiv_k z$. $\blacksquare$
\end{enumerate}
\resultado \textbf{$\equiv_k$ es reflexiva, simétrica y transitiva} (una relación de equivalencia).
\comprobacion $\equiv_k$ coincide con la congruencia de la \cref{def:congruencia}: $x=y+kq\iff k\divides(x-y)$. Ejemplo: $17\equiv_5 2$ ($q=3$) y $2\equiv_5 17$ ($q=-3$).
\end{solucion}

\begin{solucion}{eje:3-real}
\datos $x,y\in\R$ con $y^{3}+yx^{2}\le x^{3}+xy^{2}$.
\planteamiento contrarrecíproco: si $y>x$, entonces $y^{3}+yx^{2}>x^{3}+xy^{2}$. Factorizar la diferencia.
\pasos
\begin{enumerate}[label=\arabic*.]
  \item Se demuestra el contrarrecíproco. Sean $x,y\in\R$ con $y>x$.
  \item Se calcula $(y^{3}+yx^{2})-(x^{3}+xy^{2})=(y^{3}-x^{3})-xy(y-x)=(y-x)(y^{2}+xy+x^{2})-xy(y-x)=(y-x)(x^{2}+y^{2})$.
  \item Como $y>x$, $y-x>0$. Además, $x$ e $y$ no son ambos nulos (son distintos), luego $x^{2}+y^{2}>0$.
  \item El producto de dos positivos es positivo: $y^{3}+yx^{2}>x^{3}+xy^{2}$, que es la negación de la hipótesis. $\blacksquare$
\end{enumerate}
\resultado \textbf{la hipótesis implica $y\le x$.}
\comprobacion $x=1$, $y=2$: $8+2=10>1+4=5$, y $(y-x)(x^{2}+y^{2})=1\cdot 5=5=10-5$.
\end{solucion}

\begin{solucion}{eje:3-bicond}
\datos $n\in\Z$.
\planteamiento dos implicaciones: ($\Leftarrow$) directa; ($\Rightarrow$) por contrarrecíproco.
\pasos
\begin{enumerate}[label=\arabic*.]
  \item ($\Leftarrow$) Sea $n$ par, $n=2k$ con $k\in\Z$. Entonces $n^{2}=4k^{2}$ y $4\divides n^{2}$.
  \item ($\Rightarrow$) Por contrarrecíproco: sea $n$ impar. Por la \cref{prop:impar-cuadrado}, $n^{2}$ es impar.
  \item Si fuera $4\divides n^{2}$, sería $n^{2}=4j=2(2j)$ par, imposible. Luego $4\nodivide n^{2}$.
  \item Se han probado ambas implicaciones, luego $4\divides n^{2}\iff n$ par. $\blacksquare$
\end{enumerate}
\resultado \textbf{$4\divides n^{2}$ si y solo si $n$ es par.}
\comprobacion $n=6$: $36=4\cdot 9$; $n=5$: $25$ no es múltiplo de~4. Coherente con la \cref{prop:cuadrado-mod4}.
\end{solucion}

\begin{solucion}{eje:3-restos}
\datos $3^{100}\bmod 7$, $7^{2026}\bmod 10$, $(5^{40}+6^{40})\bmod 11$.
\planteamiento encontrar una potencia congruente con $\pm1$ y aplicar la \cref{prop:congruencias}.
\pasos
\begin{enumerate}[label=(\alph*)]
  \item $3^{3}=27\equiv -1\pmod 7$, luego $3^{6}\equiv 1$. Como $100=6\cdot 16+4$, $3^{100}\equiv 3^{4}=81=7\cdot 11+4\equiv 4$.
  \item $7^{2}=49\equiv -1\pmod{10}$, luego $7^{4}\equiv 1$. Como $2026=4\cdot 506+2$, $7^{2026}\equiv 7^{2}\equiv 9\pmod{10}$.
  \item $5^{2}=25\equiv 3$, $5^{5}=5^{4}\cdot 5\equiv 9\cdot 5=45\equiv 1\pmod{11}$, luego $5^{40}=(5^{5})^{8}\equiv 1$. $6^{2}=36\equiv 3$, $6^{5}=6^{4}\cdot 6\equiv 9\cdot 6=54\equiv -1$, luego $6^{40}=(6^{5})^{8}\equiv 1$. La suma es $\equiv 2$.
\end{enumerate}
\resultado (a) \textbf{4}; (b) \textbf{9}; (c) \textbf{2}.
\comprobacion cálculo exacto con Python: \texttt{pow(3,100,7)=4}, \texttt{pow(7,2026,10)=9}, $(5^{40}+6^{40})\bmod 11=2$.
\end{solucion}

\begin{solucion}{eje:3-division}
\datos $2026$ y $7$; $-17$ y $5$; $2023$, $2026$; módulo~9.
\planteamiento buscar $q$ con $qb\le a<(q+1)b$ y tomar $r=a-qb$ (\cref{teo:division}); aplicar la \cref{def:divide}.
\pasos
\begin{enumerate}[label=\arabic*.]
  \item $7\cdot 289=2023\le 2026<2030=7\cdot 290$: $q=289$, $r=2026-2023=3$.
  \item $5\cdot(-4)=-20\le -17<-15=5\cdot(-3)$: $q=-4$, $r=-17+20=3$. (No vale $q=-3$, $r=-2$, porque el resto debe ser no negativo.)
  \item $7\divides 2023$ si existe $c\in\Z$ con $2023=7c$: sí, $c=289$. $7\divides 2026$ exigiría $2026=7c$, pero por la unicidad de la división $2026=7\cdot 289+3$ con resto $3\neq 0$: no.
  \item $2026=9\cdot 225+1$, resto $1$. Suma de cifras: $2+0+2+6=10=9+1$, resto $1$. Coinciden porque $10\equiv 1\pmod 9$, luego $10^{k}\equiv 1$ y cada cifra $d$ en la posición $k$ aporta $d\cdot 10^{k}\equiv d$ (\cref{prop:congruencias}).
\end{enumerate}
\resultado (a) $2026=7\cdot 289+3$; $-17=5\cdot(-4)+3$. (b) $7\divides 2023$, $7\nodivide 2026$. (c) resto $1$ en ambos casos.
\comprobacion Python: \texttt{divmod(2026,7)=(289,3)}, \texttt{divmod(-17,5)=(-4,3)}, \texttt{2026\%9=1}.
\end{solucion}

% ---------------------------------------------------------------------
\section{Tema 4}

\begin{solucion}{eje:4-suponer}
\datos cuatro enunciados.
\planteamiento se supone la negación (\cref{def:absurdo}); para un condicional, hipótesis y negación de la tesis.
\pasos
\begin{enumerate}[label=(\alph*)]
  \item Se supone que $\sqrt5$ es racional: $\sqrt5=a/b$ con $a,b\in\Z$, $b>0$, fracción irreducible.
  \item Se supone que existe $n\in\N$ con $3\divides n^{2}+1$.
  \item Se suponen $a,b\in\Z$ con $ab$ impar y con $a$ par o $b$ par.
  \item Se supone que todo primo es menor o igual que $10^{100}$.
\end{enumerate}
\resultado las cuatro suposiciones anteriores.
\comprobacion en (c), la negación de «$a$ y $b$ impares» es «$a$ par o $b$ par» (De Morgan); en (b), la negación de $\forall n\,\neg Q(n)$ es $\exists n\,Q(n)$.
\end{solucion}

\begin{solucion}{eje:4-raiz3}
\datos $n\in\Z$; se quiere $3\divides n^{2}\Rightarrow 3\divides n$ y $\sqrt3\notin\Q$.
\planteamiento lema por contrarrecíproco con los restos módulo~3; teorema imitando el \cref{teo:raiz2}.
\pasos
\begin{enumerate}[label=\arabic*.]
  \item \emph{Lema.} Se prueba el contrarrecíproco: si $3\nodivide n$, entonces $3\nodivide n^{2}$. Si $3\nodivide n$, por la división euclídea $n=3k+1$ o $n=3k+2$ con $k\in\Z$.
  \item Caso $n=3k+1$: $n^{2}=9k^{2}+6k+1=3(3k^{2}+2k)+1$. Caso $n=3k+2$: $n^{2}=9k^{2}+12k+4=3(3k^{2}+4k+1)+1$. En ambos, $n^{2}$ deja resto~1 al dividir por~3, luego $3\nodivide n^{2}$ (por la unicidad del resto).
  \item \emph{Teorema.} Supóngase, por reducción al absurdo, que $\sqrt3=a/b$ con $a,b\in\Z$, $b>0$ y fracción irreducible. Entonces $a^{2}=3b^{2}$, luego $3\divides a^{2}$ y, por el lema, $3\divides a$: $a=3c$.
  \item Sustituyendo, $9c^{2}=3b^{2}$, es decir, $b^{2}=3c^{2}$, luego $3\divides b^{2}$ y $3\divides b$.
  \item Así $3$ divide a $a$ y a $b$, en contra de que la fracción sea irreducible. Por tanto, $\sqrt3$ es irracional. $\blacksquare$
\end{enumerate}
\resultado \textbf{$3\divides n^{2}\Rightarrow 3\divides n$, y $\sqrt3$ es irracional.}
\comprobacion los cuadrados de $0,1,2$ son $0,1,4$, con restos $0,1,1$ módulo~3: solo los múltiplos de~3 tienen cuadrado múltiplo de~3.
\end{solucion}

\begin{solucion}{eje:4-raiz-cubica}
\datos $\sqrt[3]{2}$.
\planteamiento reducción al absurdo; se necesita «$n^{3}$ par $\Rightarrow n$ par», que es el contrarrecíproco del \cref{eje:3-cubo}.
\pasos
\begin{enumerate}[label=\arabic*.]
  \item \emph{Lema.} Si $n^{3}$ es par, $n$ es par: si $n$ fuera impar, $n^{3}$ sería impar por el \cref{eje:3-cubo}.
  \item Supóngase, por reducción al absurdo, que $\sqrt[3]{2}=a/b$ con $a,b\in\Z$, $b>0$ y fracción irreducible. Elevando al cubo, $a^{3}=2b^{3}$.
  \item Entonces $a^{3}$ es par y, por el lema, $a=2c$ con $c\in\Z$. Sustituyendo, $8c^{3}=2b^{3}$, es decir, $b^{3}=4c^{3}=2(2c^{3})$ es par, y $b$ es par.
  \item Así $2$ divide a $a$ y a $b$, contradicción con la irreducibilidad. Luego $\sqrt[3]{2}$ es irracional. $\blacksquare$
\end{enumerate}
\resultado \textbf{$\sqrt[3]{2}$ es irracional.}
\comprobacion la prueba es la del \cref{teo:raiz2} cambiando cuadrados por cubos; numéricamente $\sqrt[3]2\approx 1{,}2599$ y ninguna fracción de denominador pequeño cumple $a^{3}=2b^{3}$ (por ejemplo, $5^{3}=125\neq 2\cdot 4^{3}=128$).
\end{solucion}

\begin{solucion}{eje:4-log}
\datos $\log_2 3$, el número $x$ con $2^{x}=3$.
\planteamiento suponer $x=a/b$ y obtener una igualdad de enteros con paridades distintas.
\pasos
\begin{enumerate}[label=\arabic*.]
  \item Como $2^{1}<3<2^{2}$, se tiene $1<\log_2 3<2$; en particular, $\log_2 3>0$.
  \item Supóngase, por reducción al absurdo, que $\log_2 3=a/b$ con $a,b\in\N$ (ambos positivos, pues el cociente lo es y se puede tomar $b>0$).
  \item Entonces $2^{a/b}=3$ y, elevando a $b$, $2^{a}=3^{b}$.
  \item Como $a\ge 1$, $2^{a}$ es par; como $3^{b}$ es producto de impares, es impar (\cref{prop:congruencias}(iii) con $3\equiv 1\pmod 2$). Un entero no puede ser par e impar: contradicción.
  \item Por tanto, $\log_2 3$ es irracional. $\blacksquare$
\end{enumerate}
\resultado \textbf{$\log_2 3$ es irracional.}
\comprobacion $\log_2 3\approx 1{,}585$; la hipótesis $a\ge 1$ (paso~2) es esencial, pues $2^{0}=3^{0}$.
\end{solucion}

\begin{solucion}{eje:4-racional-irracional}
\datos $a\in\Q$, $b\notin\Q$.
\planteamiento reducción al absurdo: si $a+b$ (o $ab$) fuera racional, $b$ se despejaría como racional.
\pasos
\begin{enumerate}[label=\arabic*.]
  \item Escríbase $a=m/n$ con $m,n\in\Z$, $n\neq 0$.
  \item \emph{Suma.} Supóngase que $a+b=p/q$ con $p,q\in\Z$, $q\neq 0$. Entonces $b=\frac{p}{q}-\frac{m}{n}=\frac{pn-mq}{qn}$, con $pn-mq\in\Z$ y $qn\neq 0$; luego $b$ es racional, contradicción.
  \item \emph{Producto}, con $a\neq 0$ (es decir, $m\neq 0$). Supóngase $ab=p/q$. Entonces $b=\frac{p}{q}\cdot\frac{n}{m}=\frac{pn}{qm}$ con $qm\neq 0$, luego $b$ es racional, contradicción.
  \item Si $a=0$, $ab=0$ es racional: la hipótesis $a\neq 0$ es necesaria. $\blacksquare$
\end{enumerate}
\resultado \textbf{$a+b$ es irracional; $ab$ es irracional si $a\neq 0$}, y racional ($=0$) si $a=0$.
\comprobacion ejemplos: $1+\sqrt2$ y $3\sqrt2$ son irracionales; $0\cdot\sqrt2=0$.
\end{solucion}

\begin{solucion}{eje:4-diofantica}
\datos $21a+30b=1$ y $21a+30b=3$.
\planteamiento $3$ divide a $21$ y a $30$, luego a toda combinación $21a+30b$ (\cref{prop:divisibilidad}(ii)).
\pasos
\begin{enumerate}[label=\arabic*.]
  \item Supóngase, por reducción al absurdo, que existen $a,b\in\Z$ con $21a+30b=1$.
  \item Como $21=3\cdot 7$ y $30=3\cdot 10$, $21a+30b=3(7a+10b)$, luego $3\divides 1$.
  \item Pero los únicos divisores de~1 son $\pm1$, y $3\neq\pm1$. Contradicción; no existen tales enteros. $\blacksquare$
  \item Para el valor~3 sí hay solución: $a=3$, $b=-2$, pues $63-60=3$.
\end{enumerate}
\resultado \textbf{no hay solución para~1; sí para~3} (por ejemplo, $(a,b)=(3,-2)$).
\comprobacion $\operatorname{mcd}(21,30)=3$: las combinaciones enteras de $21$ y $30$ son exactamente los múltiplos de~3 (identidad de Bézout, que excede este módulo).
\end{solucion}

\begin{solucion}{eje:4-mod4}
\datos $n\in\Z$; enteros $4k+3$.
\planteamiento \cref{prop:cuadrado-mod4}: todo cuadrado es $\equiv 0$ o $1\pmod 4$.
\pasos
\begin{enumerate}[label=\arabic*.]
  \item Sea $n\in\Z$. Por la \cref{prop:cuadrado-mod4}, $n^{2}\equiv 0$ o $n^{2}\equiv 1\pmod 4$.
  \item Caso $n^{2}\equiv 0$: $n^{2}+2\equiv 2\pmod 4$. Caso $n^{2}\equiv 1$: $n^{2}+2\equiv 3\pmod 4$. En ningún caso es $\equiv 0$, luego $4\nodivide n^{2}+2$.
  \item Sean ahora $x,y\in\Z$. Por la misma proposición, $x^{2}+y^{2}$ es congruente con $0+0$, $0+1$, $1+0$ o $1+1$, es decir, con $0$, $1$ o $2\pmod 4$.
  \item Un entero $4k+3$ es $\equiv 3\pmod 4$, que no está entre esos valores. Por la unicidad del resto, $x^{2}+y^{2}\neq 4k+3$. $\blacksquare$
\end{enumerate}
\resultado \textbf{$4\nodivide n^{2}+2$ y ningún entero $4k+3$ es suma de dos cuadrados.}
\comprobacion $3$, $7$, $11$, $15$ no son sumas de dos cuadrados; $5=1+4$ y $13=4+9$ sí, y son $\equiv 1\pmod 4$.
\end{solucion}

\begin{solucion}{eje:4-24}
\datos $p$ primo, $p>3$.
\planteamiento $p^{2}-1=(p-1)(p+1)$; probar por separado $8\divides p^{2}-1$ y $3\divides p^{2}-1$.
\pasos
\begin{enumerate}[label=\arabic*.]
  \item Sea $p>3$ primo. No es par (sería divisible por~2 y distinto de~2) ni múltiplo de~3 (sería divisible por~3 y distinto de~3).
  \item \emph{Factor 8.} Como $p$ es impar, $p=2k+1$ y $p^{2}-1=(p-1)(p+1)=2k(2k+2)=4k(k+1)$. Por la \cref{prop:consecutivos}, $k(k+1)=k^{2}+k$ es par, $k(k+1)=2j$, luego $p^{2}-1=8j$.
  \item \emph{Factor 3.} Como $3\nodivide p$, $p=3m+1$ o $p=3m+2$. En el primer caso, $3\divides p-1$; en el segundo, $p+1=3(m+1)$ y $3\divides p+1$. En ambos, $3\divides(p-1)(p+1)=p^{2}-1$.
  \item Por los pasos 2 y 3, $3\divides 8j$. Como también $3\divides 9j$, la \cref{prop:divisibilidad}(ii) da $3\divides 9j-8j=j$. Así $j=3i$ con $i\in\Z$ y $p^{2}-1=24i$. $\blacksquare$
\end{enumerate}
\resultado \textbf{$24\divides p^{2}-1$ para todo primo $p>3$.}
\comprobacion $5^{2}-1=24$, $7^{2}-1=48$, $11^{2}-1=120=24\cdot 5$; se ha comprobado con Python para todos los primos entre $5$ y $10^{4}$. Para $p=3$ falla ($8$), lo que muestra que la hipótesis $p>3$ se usa (paso~3).
\end{solucion}

\begin{solucion}{eje:4-maximo}
\datos $x,y\in\R$.
\planteamiento dos casos según el signo de $x-y$.
\pasos
\begin{enumerate}[label=\arabic*.]
  \item Caso $x\ge y$: $\abs{x-y}=x-y$, luego $\frac{x+y+\abs{x-y}}{2}=\frac{2x}{2}=x=\max(x,y)$ y $\frac{x+y-\abs{x-y}}{2}=y=\min(x,y)$.
  \item Caso $x<y$: $\abs{x-y}=y-x$, luego la primera expresión vale $\frac{2y}{2}=y=\max(x,y)$ y la segunda $x=\min(x,y)$.
  \item Comprobación con $(3,-5)$: $x+y=-2$, $\abs{x-y}=8$; $\max=(-2+8)/2=3$, $\min=(-2-8)/2=-5$.
\end{enumerate}
\resultado \textbf{ambas fórmulas son ciertas}; para $(3,-5)$ dan $3$ y $-5$.
\comprobacion sumando las fórmulas, $\max+\min=x+y$; restándolas, $\max-\min=\abs{x-y}$, ambas evidentes.
\end{solucion}

\begin{solucion}{eje:4-restos-cuadrados}
\datos $n\in\Z$; módulos $3$, $5$ y $8$.
\planteamiento por la \cref{prop:congruencias}, $n^{2}\bmod m$ solo depende del resto de $n$ módulo $m$; se recorren todos los restos.
\pasos
\begin{enumerate}[label=\arabic*.]
  \item Módulo 3, restos $0,1,2$: cuadrados $0,1,4\equiv 1$. Valores posibles: $\{0,1\}$.
  \item Módulo 5, restos $0,\dots,4$: cuadrados $0,1,4,9\equiv 4,16\equiv 1$. Valores posibles: $\{0,1,4\}$.
  \item Módulo 8, restos $0,\dots,7$: cuadrados $0,1,4,9\equiv1,16\equiv0,25\equiv1,36\equiv4,49\equiv1$. Valores: $\{0,1,4\}$; los restos impares $1,3,5,7$ dan siempre~$1$.
  \item (a) $n^{2}+1\equiv 1$ o $2\pmod 3$, nunca $0$. (b) $2$ y $3$ no están en $\{0,1,4\}$. (c) Un impar tiene resto impar módulo~8 y, por el paso~3, $n^{2}\equiv 1$.
\end{enumerate}
\resultado cuadrados: módulo~3 $\{0,1\}$; módulo~5 $\{0,1,4\}$; módulo~8 $\{0,1,4\}$. Las tres afirmaciones son ciertas.
\comprobacion $7^{2}=49=8\cdot 6+1$; $13^{2}=169=5\cdot 33+4$; $5^{2}+1=26$, no múltiplo de~3. El apartado (c) reaparece en el \cref{eje:4-24}.
\end{solucion}

\begin{solucion}{eje:4-triangular-inversa}
\datos $x,y\in\R$.
\planteamiento aplicar el \cref{teo:triangular} dos veces y combinar con $\abs{t}\le c\iff -c\le t\le c$.
\pasos
\begin{enumerate}[label=\arabic*.]
  \item Sean $x,y\in\R$. Como $x=(x-y)+y$, la desigualdad triangular da $\abs{x}\le\abs{x-y}+\abs{y}$, es decir, $\abs{x}-\abs{y}\le\abs{x-y}$.
  \item Como $y=(y-x)+x$, del mismo modo $\abs{y}-\abs{x}\le\abs{y-x}=\abs{x-y}$, es decir, $-\abs{x-y}\le\abs{x}-\abs{y}$.
  \item Juntando, $-\abs{x-y}\le\abs{x}-\abs{y}\le\abs{x-y}$. Si $t=\abs{x}-\abs{y}\ge 0$, $\abs{t}=t\le\abs{x-y}$; si $t<0$, $\abs{t}=-t\le\abs{x-y}$. En ambos casos, $\bigl\lvert\abs{x}-\abs{y}\bigr\rvert\le\abs{x-y}$. $\blacksquare$
\end{enumerate}
\resultado \textbf{desigualdad triangular inversa demostrada.}
\comprobacion $x=3$, $y=-5$: $\bigl\lvert 3-5\bigr\rvert=2\le\abs{8}=8$; $x=3$, $y=5$: $2\le 2$ (igualdad cuando tienen el mismo signo).
\end{solucion}

\begin{solucion}{eje:4-producto-par}
\datos $m,n\in\Z$, $mn$ par, $m$ impar.
\planteamiento la negación de $P\Rightarrow Q$ es $P\wedge\neg Q$: suponer $mn$ par, $m$ impar y $n$ impar.
\pasos
\begin{enumerate}[label=\arabic*.]
  \item Supóngase, por reducción al absurdo, que existen enteros $m$ y $n$ con $mn$ par, $m$ impar y $n$ no par, es decir, impar.
  \item Existen $k,r\in\Z$ con $m=2k+1$ y $n=2r+1$.
  \item Entonces $mn=4kr+2k+2r+1=2(2kr+k+r)+1$ es impar.
  \item Esto contradice que $mn$ sea par, pues ningún entero es par e impar. Luego, si $mn$ es par y $m$ impar, $n$ es par. $\blacksquare$
\end{enumerate}
\resultado \textbf{$n$ es par.}
\comprobacion $m=3$, $mn=12$: $n=4$ es par. La prueba usa la suposición «$n$ impar» en el paso~2, de modo que no es un «sándwich de contradicción».
\end{solucion}

\begin{solucion}{eje:4-cubos}
\datos todo entero $x$; se quieren cubos de enteros.
\planteamiento (a) identidad $(n+1)^{3}+(n-1)^{3}=2n^{3}+6n$; (b) $x^{3}-x=(x-1)x(x+1)$; (c) combinar.
\pasos
\begin{enumerate}[label=\arabic*.]
  \item (a) Sea $N=6n$ con $n\in\Z$. Se desarrolla $(n+1)^{3}+(n-1)^{3}=(n^{3}+3n^{2}+3n+1)+(n^{3}-3n^{2}+3n-1)=2n^{3}+6n$.
  \item Por tanto, $6n=(n+1)^{3}+(n-1)^{3}+(-n)^{3}+(-n)^{3}$, suma de cuatro cubos de enteros.
  \item (b) Sea $x\in\Z$. $x^{3}-x=(x-1)x(x+1)$ es producto de tres enteros consecutivos. Entre $x-1$ y $x$ hay uno par, de modo que $2\divides x^{3}-x$. Por casos según el resto de $x$ módulo~3: si $x\equiv 0$, $3\divides x$; si $x\equiv 1$, $3\divides x-1$; si $x\equiv 2$, $3\divides x+1$. Luego $3\divides x^{3}-x$.
  \item Escríbase $x^{3}-x=2s$ con $s\in\Z$. Como $3\divides 2s$ y $3\divides 3s$, se tiene $3\divides 3s-2s=s$ (\cref{prop:divisibilidad}(ii)); así $s=3t$ y $x^{3}-x=6t$, es decir, $x^{3}\equiv x\pmod 6$.
  \item (c) Sea $x\in\Z$. Por (b), $x^{3}-x=6t$ con $t\in\Z$, luego $x=x^{3}+6(-t)$. Por (a), $6(-t)$ es suma de cuatro cubos, y $x$ es suma de $x^{3}$ y esos cuatro: cinco cubos. $\blacksquare$
  \item Para $x=7$: $7^{3}-7=336=6\cdot 56$, $t=56$, $-t=-56$, y $6(-56)=(-55)^{3}+(-57)^{3}+56^{3}+56^{3}$. Así, $7=7^{3}+(-55)^{3}+(-57)^{3}+56^{3}+56^{3}$.
\end{enumerate}
\resultado \textbf{todo entero es suma de cinco cubos}; $7=7^{3}+(-55)^{3}+(-57)^{3}+2\cdot 56^{3}$.
\comprobacion $343-166\,375-185\,193+175\,616+175\,616=7$ (verificado con Python). No se sabe si basta con cuatro cubos para todo entero~\cite{cs103}.
\end{solucion}

\begin{solucion}{eje:4-euclides-n}
\datos los seis primeros primos $2,3,5,7,11,13$.
\planteamiento calcular los productos y comprobar la primalidad dividiendo por los primos hasta la raíz cuadrada.
\pasos
\begin{enumerate}[label=\arabic*.]
  \item $N_1=3$, $N_2=7$, $N_3=31$, $N_4=211$, $N_5=2311$, $N_6=30\,031$.
  \item $3,7,31$ son primos. $211$: no es divisible por $2,3,5,7,11,13$ y $\sqrt{211}<15$: primo. $2311$: no es divisible por ningún primo hasta $47$ y $\sqrt{2311}<49$: primo.
  \item $30\,031=59\cdot 509$: compuesto.
  \item No hay contradicción: la demostración solo afirma que los divisores primos de $N_k$ no están entre $p_1,\dots,p_k$, y en efecto $59$ y $509$ son mayores que~13.
\end{enumerate}
\resultado $N_1,\dots,N_5$ \textbf{primos}; $N_6=\mathbf{59\cdot 509}$.
\comprobacion $59\cdot 509=30\,031$; factorización confirmada con Python (\texttt{sympy.factorint}).
\end{solucion}

\begin{solucion}{eje:4-4k3}
\datos primos $\equiv 3\pmod 4$; hay al menos uno, el~3.
\planteamiento reducción al absurdo con $N=4p_1\cdots p_m-1$.
\pasos
\begin{enumerate}[label=\arabic*.]
  \item Supóngase que solo hay un número finito de primos de la forma $4k+3$, y sean $p_1,\dots,p_m$ todos ellos ($m\ge 1$, pues $3$ es uno). Sea $N=4p_1\cdots p_m-1$. Entonces $N\ge 4\cdot 3-1=11$ y $N\equiv -1\equiv 3\pmod 4$; en particular, $N$ es impar.
  \item Por el \cref{teo:factorizacion}, $N$ es producto de primos $q_1\cdots q_s$ (con $s=1$ si $N$ es primo). Todos son impares porque $N$ lo es, luego cada $q_j\equiv 1$ o $q_j\equiv 3\pmod 4$.
  \item Si todos fueran $\equiv 1\pmod 4$, su producto sería $\equiv 1\cdot1\cdots1=1\pmod 4$ (\cref{prop:congruencias}(ii)), en contra de $N\equiv 3$. Luego algún $q_j\equiv 3\pmod 4$.
  \item Por la hipótesis, $q_j=p_i$ para algún $i$. Entonces $q_j$ divide a $N$ y a $4p_1\cdots p_m$, luego divide a su diferencia, $1$. Imposible, porque $q_j\ge 3$.
  \item Por tanto, hay infinitos primos de la forma $4k+3$. $\blacksquare$
\end{enumerate}
\resultado \textbf{existen infinitos primos $\equiv 3\pmod 4$.}
\comprobacion con la lista $\{3\}$: $N=11$, primo y $\equiv 3$; con $\{3,7,11\}$: $N=923=13\cdot 71$, y $71\equiv 3\pmod 4$ no está en la lista. (La analogía con $4p_1\cdots p_m+1$ no funciona para los primos $4k+1$: un producto de primos $\equiv 3$ puede ser $\equiv 1$.)
\end{solucion}

\begin{solucion}{eje:4-raiz-primo}
\datos $p$ primo; lema de Euclides.
\planteamiento como en el \cref{teo:raiz2}; el lema de Euclides con $a=b=n$ da $p\divides n^{2}\Rightarrow p\divides n$.
\pasos
\begin{enumerate}[label=\arabic*.]
  \item Supóngase, por reducción al absurdo, que $\sqrt p=a/b$ con $a,b\in\Z$, $b>0$ y fracción irreducible. Entonces $a^{2}=pb^{2}$, luego $p\divides a\cdot a$.
  \item Por el lema de Euclides, $p\divides a$: $a=pc$ con $c\in\Z$. Sustituyendo, $p^{2}c^{2}=pb^{2}$, es decir, $b^{2}=pc^{2}$, y de nuevo $p\divides b$.
  \item Así $p$ divide a $a$ y a $b$, y como $p\ge 2$, contradice la irreducibilidad. Luego $\sqrt p$ es irracional. $\blacksquare$
  \item Con $4$ (no primo) falla el paso~2: $4\divides 2\cdot 2$ pero $4\nodivide 2$. De hecho $\sqrt4=2/1$ es racional.
\end{enumerate}
\resultado \textbf{$\sqrt p$ es irracional para todo primo $p$}; con $\sqrt4$ falla el lema de Euclides.
\comprobacion los casos $p=2$ y $p=3$ son el \cref{teo:raiz2} y el \cref{eje:4-raiz3}, donde el lema de Euclides se sustituyó por un argumento de restos.
\end{solucion}

% ---------------------------------------------------------------------
\section{Tema 5}

\begin{solucion}{eje:5-plantilla}
\datos $\sum_{k=1}^{n}(2k-1)=n^{2}$ para $n\ge 1$.
\planteamiento identificar las piezas del \cref{teo:induccion} con caso base $n_0=1$.
\pasos
\begin{enumerate}[label=\arabic*.]
  \item (a) $P(n)$: «$1+3+\dots+(2n-1)=n^{2}$».
  \item (b) $P(1)$: «$1=1^{2}$». $P(n+1)$: «$1+3+\dots+(2n-1)+(2n+1)=(n+1)^{2}$».
  \item (c) Hay que demostrar: para todo $n\ge 1$, si $P(n)$ es cierta, también lo es $P(n+1)$.
  \item Demostración. \emph{Caso base:} $P(1)$ es $1=1$. \emph{Paso:} sea $n\ge 1$ con $P(n)$. Entonces $1+\dots+(2n-1)+(2n+1)=n^{2}+2n+1=(n+1)^{2}$, que es $P(n+1)$. Por inducción, $P(n)$ vale para todo $n\ge 1$. $\blacksquare$
\end{enumerate}
\resultado \textbf{la suma de los $n$ primeros impares es $n^{2}$.}
\comprobacion $1+3+5+7=16=4^{2}$.
\end{solucion}

\begin{solucion}{eje:5-cuadrados}
\datos $P(n)$: $\sum_{k=1}^{n}k^{2}=\frac{n(n+1)(2n+1)}{6}$.
\planteamiento inducción simple desde $n=1$.
\pasos
\begin{enumerate}[label=\arabic*.]
  \item \emph{Caso base.} $n=1$: $1^{2}=1$ y $\frac{1\cdot 2\cdot 3}{6}=1$.
  \item \emph{Paso.} Sea $n\ge 1$ y supóngase $P(n)$. Entonces
  \[
  \begin{aligned}
  \sum_{k=1}^{n+1}k^{2}&=\frac{n(n+1)(2n+1)}{6}+(n+1)^{2}\\
  &=\frac{(n+1)\bigl[n(2n+1)+6(n+1)\bigr]}{6}=\frac{(n+1)(2n^{2}+7n+6)}{6}.
  \end{aligned}
  \]
  \item Como $2n^{2}+7n+6=(n+2)(2n+3)$, la suma es $\frac{(n+1)(n+2)(2(n+1)+1)}{6}$, que es $P(n+1)$.
  \item Por el principio de inducción, $P(n)$ vale para todo $n\in\N$. $\blacksquare$
\end{enumerate}
\resultado $\sum_{k=1}^{n}k^{2}=\dfrac{n(n+1)(2n+1)}{6}$.
\comprobacion $n=4$: $1+4+9+16=30$ y $\frac{4\cdot5\cdot9}{6}=30$. Factorización: $(n+2)(2n+3)=2n^{2}+3n+4n+6$.
\end{solucion}

\begin{solucion}{eje:5-cubos}
\datos $P(n)$: $\sum_{k=1}^{n}k^{3}=\bigl(\frac{n(n+1)}{2}\bigr)^{2}$.
\planteamiento inducción simple; sacar factor común $(n+1)^{2}/4$.
\pasos
\begin{enumerate}[label=\arabic*.]
  \item \emph{Caso base.} $n=1$: $1=(1\cdot 2/2)^{2}=1$.
  \item \emph{Paso.} Sea $n\ge 1$ con $P(n)$. Entonces
  \[
  \sum_{k=1}^{n+1}k^{3}=\frac{n^{2}(n+1)^{2}}{4}+(n+1)^{3}=\frac{(n+1)^{2}\bigl[n^{2}+4(n+1)\bigr]}{4}=\frac{(n+1)^{2}(n+2)^{2}}{4},
  \]
  que es $P(n+1)$.
  \item Por inducción, $P(n)$ vale para todo $n\in\N$. $\blacksquare$
\end{enumerate}
\resultado $\sum_{k=1}^{n}k^{3}=\Bigl(\dfrac{n(n+1)}{2}\Bigr)^{2}$, el cuadrado de la suma del \cref{ej:gauss}.
\comprobacion $n=3$: $1+8+27=36=6^{2}$.
\end{solucion}

\begin{solucion}{eje:5-telescopica}
\datos $P(n)$: $\sum_{k=1}^{n}\frac{1}{k(k+1)}=\frac{n}{n+1}$.
\planteamiento inducción; después, suma telescópica.
\pasos
\begin{enumerate}[label=\arabic*.]
  \item \emph{Caso base.} $n=1$: $\frac{1}{2}=\frac{1}{2}$.
  \item \emph{Paso.} Si vale $P(n)$, entonces $\sum_{k=1}^{n+1}\frac{1}{k(k+1)}=\frac{n}{n+1}+\frac{1}{(n+1)(n+2)}=\frac{n(n+2)+1}{(n+1)(n+2)}=\frac{(n+1)^{2}}{(n+1)(n+2)}=\frac{n+1}{n+2}$, que es $P(n+1)$. Por inducción, vale para todo $n$. $\blacksquare$
  \item \emph{Segunda prueba.} Como $\frac1k-\frac1{k+1}=\frac{k+1-k}{k(k+1)}=\frac{1}{k(k+1)}$, la suma es $\bigl(1-\tfrac12\bigr)+\bigl(\tfrac12-\tfrac13\bigr)+\dots+\bigl(\tfrac1n-\tfrac1{n+1}\bigr)$; cada término intermedio aparece una vez con signo $+$ y otra con signo $-$, y queda $1-\frac{1}{n+1}=\frac{n}{n+1}$.
\end{enumerate}
\resultado $\sum_{k=1}^{n}\dfrac{1}{k(k+1)}=\dfrac{n}{n+1}$.
\comprobacion $n=3$: $\frac12+\frac16+\frac1{12}=\frac{6+2+1}{12}=\frac{9}{12}=\frac34$. El argumento telescópico, escrito con rigor, es a su vez una inducción.
\end{solucion}

\begin{solucion}{eje:5-potencia}
\datos $P(n)$: $2^{n}>n^{2}$, para $n\ge 5$.
\planteamiento inducción con caso base $n=5$; en el paso hay que probar $2n^{2}\ge(n+1)^{2}$ para $n\ge 5$.
\pasos
\begin{enumerate}[label=\arabic*.]
  \item \emph{Caso base.} $2^{5}=32>25=5^{2}$.
  \item \emph{Desigualdad auxiliar.} Para $n\ge 5$: $2n^{2}-(n+1)^{2}=n^{2}-2n-1=(n-1)^{2}-2\ge 4^{2}-2>0$.
  \item \emph{Paso.} Sea $n\ge 5$ con $2^{n}>n^{2}$. Entonces $2^{n+1}=2\cdot 2^{n}>2n^{2}>(n+1)^{2}$, por el paso~2.
  \item Por inducción, $2^{n}>n^{2}$ para todo $n\ge 5$. $\blacksquare$
  \item Valores pequeños: $n=1$: $2>1$ (cierta); $n=2$: $4>4$ (falsa); $n=3$: $8>9$ (falsa); $n=4$: $16>16$ (falsa).
\end{enumerate}
\resultado \textbf{cierta para $n\ge 5$ y para $n=1$; falsa para $n=2,3,4$.}
\comprobacion $n=6$: $64>36$; $n=10$: $1024>100$. El paso auxiliar falla para $n=2$ ($(n-1)^{2}-2=-1$), lo que explica que no pueda empezarse antes.
\end{solucion}

\begin{solucion}{eje:5-div9}
\datos $P(n)$: $9\divides 4^{n}+6n-1$.
\planteamiento inducción; expresar $4^{n+1}+6(n+1)-1$ en función de $4^{n}+6n-1$.
\pasos
\begin{enumerate}[label=\arabic*.]
  \item \emph{Caso base.} $n=0$: $1+0-1=0=9\cdot 0$.
  \item \emph{Paso.} Sea $n\ge 0$ con $4^{n}+6n-1=9m$, $m\in\Z$. Entonces $4^{n}=9m-6n+1$ y
  \[
  4^{n+1}+6(n+1)-1=4(9m-6n+1)+6n+5=36m-18n+9=9(4m-2n+1).
  \]
  \item Como $4m-2n+1\in\Z$, vale $P(n+1)$. Por inducción, $P(n)$ para todo $n\in\N_0$. $\blacksquare$
\end{enumerate}
\resultado \textbf{$9\divides 4^{n}+6n-1$ para todo $n\ge 0$.}
\comprobacion $n=1$: $4+6-1=9$; $n=2$: $16+12-1=27$; $n=3$: $64+18-1=81$. Verificado con Python hasta $n=49$.
\end{solucion}

\begin{solucion}{eje:5-bernoulli}
\datos $x\ge -1$; $P(n)$: $(1+x)^{n}\ge 1+nx$.
\planteamiento inducción; multiplicar la hipótesis por $1+x\ge 0$.
\pasos
\begin{enumerate}[label=\arabic*.]
  \item \emph{Caso base.} $n=0$: $1\ge 1$.
  \item \emph{Paso.} Sea $n\ge 0$ con $(1+x)^{n}\ge 1+nx$. Como $x\ge -1$, el número $1+x$ es no negativo, y multiplicar por él conserva la desigualdad: $(1+x)^{n+1}\ge(1+nx)(1+x)$.
  \item $(1+nx)(1+x)=1+(n+1)x+nx^{2}\ge 1+(n+1)x$, porque $nx^{2}\ge 0$.
  \item Luego $(1+x)^{n+1}\ge 1+(n+1)x$. Por inducción, la desigualdad vale para todo $n\in\N_0$. $\blacksquare$
\end{enumerate}
\resultado \textbf{$(1+x)^{n}\ge 1+nx$ para $x\ge -1$}; la hipótesis se usa en el paso~2.
\comprobacion sin la hipótesis puede fallar: con $x=-4$ y $n=3$, $(-3)^{3}=-27<1-12=-11$. Con $x=0{,}1$ y $n=10$: $1{,}1^{10}\approx 2{,}594\ge 2$.
\end{solucion}

\begin{solucion}{eje:5-sellos}
\datos $n\ge 12$; $n=4a+5b$ con $a,b\ge 0$.
\planteamiento inducción fuerte; el paso añade un~4 al caso $n-3$, que exige cuatro casos base.
\pasos
\begin{enumerate}[label=\arabic*.]
  \item \emph{Casos base.} $12=4\cdot 3$; $13=4\cdot 2+5$; $14=4+5\cdot 2$; $15=5\cdot 3$.
  \item \emph{Paso.} Sea $n\ge 15$ y supóngase el enunciado cierto para $12,13,\dots,n$. Como $12\le n-3\le n$, existen $a,b\in\N_0$ con $n-3=4a+5b$, y entonces $n+1=4(a+1)+5b$.
  \item Por inducción fuerte, todo $n\ge 12$ admite la forma. $\blacksquare$
  \item Menores que~12: los representables son $4,5,8,9,10$; \textbf{no} lo son $1,2,3,6,7,11$.
\end{enumerate}
\resultado \textbf{todo $n\ge 12$ es $4a+5b$}; fallan $1,2,3,6,7,11$.
\comprobacion $11$ no es $4a+5b$: con $b=0,1,2$ quedan $11,6,1$, ninguno múltiplo de~4. Es el mayor no representable, $4\cdot 5-4-5=11$ (fórmula de Sylvester, ampliación).
\end{solucion}

\begin{solucion}{eje:5-fib-suma}
\datos $F_1=F_2=1$, $F_{n+2}=F_{n+1}+F_n$.
\planteamiento inducción simple usando la recurrencia en el paso.
\pasos
\begin{enumerate}[label=\arabic*.]
  \item \emph{Caso base.} $n=1$: $F_1=1$ y $F_3-1=2-1=1$.
  \item \emph{Paso.} Sea $n\ge 1$ con $\sum_{k=1}^{n}F_k=F_{n+2}-1$. Entonces $\sum_{k=1}^{n+1}F_k=F_{n+2}-1+F_{n+1}=F_{n+3}-1$, por la recurrencia.
  \item Por inducción, la igualdad vale para todo $n\in\N$. $\blacksquare$
\end{enumerate}
\resultado $\sum_{k=1}^{n}F_k=F_{n+2}-1$.
\comprobacion $1+1+2+3+5=12=F_7-1=13-1$.
\end{solucion}

\begin{solucion}{eje:5-fib-cota}
\datos $F_n\le(5/3)^{n}$.
\planteamiento inducción fuerte: $F_{n+1}$ depende de $F_n$ y $F_{n-1}$.
\pasos
\begin{enumerate}[label=\arabic*.]
  \item \emph{Casos base.} $F_1=1\le 5/3$ y $F_2=1\le 25/9$.
  \item \emph{Paso.} Sea $n\ge 2$ y supóngase la cota para $1,\dots,n$. Entonces
  \[
  F_{n+1}=F_n+F_{n-1}\le\Bigl(\tfrac53\Bigr)^{n}+\Bigl(\tfrac53\Bigr)^{n-1}=\Bigl(\tfrac53\Bigr)^{n-1}\Bigl(\tfrac53+1\Bigr)=\Bigl(\tfrac53\Bigr)^{n-1}\cdot\tfrac83.
  \]
  \item Como $\frac83=\frac{24}{9}<\frac{25}{9}=\bigl(\frac53\bigr)^{2}$, se obtiene $F_{n+1}\le\bigl(\frac53\bigr)^{n+1}$.
  \item Por inducción fuerte, $F_n\le(5/3)^{n}$ para todo $n\in\N$. $\blacksquare$
  \item Hacen falta dos casos base porque el paso usa $F_{n-1}$: con $n=1$ el paso hablaría de $F_0$, que no está definido en el enunciado. El paso empieza en $n=2$ y produce $F_3$; $F_1$ y $F_2$ deben comprobarse aparte.
\end{enumerate}
\resultado \textbf{$F_n\le(5/3)^{n}$ para todo $n\ge 1$.}
\comprobacion $F_{10}=55\le(5/3)^{10}\approx 165{,}4$. Verificado con Python hasta $n=39$.
\end{solucion}

\begin{solucion}{eje:5-binario}
\datos $n\in\N$.
\planteamiento inducción fuerte; si $n$ es par, $n=2m$ con $m<n$; si es impar mayor que~1, $n=2m+1$ con $m<n$.
\pasos
\begin{enumerate}[label=\arabic*.]
  \item Sea $P(n)$: «$n$ es suma de potencias de~2 distintas». \emph{Caso base:} $1=2^{0}$.
  \item \emph{Paso.} Sea $n\ge 1$ y supóngase $P(1),\dots,P(n)$. Se prueba $P(n+1)$ distinguiendo dos casos.
  \item \emph{Caso $n+1$ par.} $n+1=2m$ con $1\le m\le n$. Por hipótesis, $m=2^{e_1}+\dots+2^{e_r}$ con exponentes distintos. Entonces $n+1=2^{e_1+1}+\dots+2^{e_r+1}$, y los exponentes siguen siendo distintos.
  \item \emph{Caso $n+1$ impar.} Como $n+1\ge 2$ es impar, $n+1\ge 3$ y $n+1=2m+1$ con $1\le m\le n$. Escribiendo $2m=2^{e_1+1}+\dots+2^{e_r+1}$ como antes, todos los exponentes son al menos~1, luego $n+1=2^{0}+2^{e_1+1}+\dots+2^{e_r+1}$ con exponentes distintos.
  \item Por inducción fuerte, $P(n)$ vale para todo $n\in\N$. $\blacksquare$
\end{enumerate}
\resultado \textbf{todo natural tiene representación binaria} (la unicidad también es cierta, aunque no se pedía).
\comprobacion $13$ es impar: $13=2\cdot 6+1$, $6=2\cdot 3$, $3=2\cdot 1+1$; resulta $13=2^{0}+2^{2}+2^{3}$, es decir, $1101_2$. En el módulo~12 esta representación es la base de la codificación de bytes.
\end{solucion}

\begin{solucion}{eje:5-bo-geometrica}
\datos $P(n)$: $\sum_{k=1}^{n}(2k-1)=n^{2}$.
\planteamiento plantilla del menor contraejemplo~\cite[§2.2]{mcs}.
\pasos
\begin{enumerate}[label=\arabic*.]
  \item Sea $C=\{n\in\N:\ P(n)\text{ es falsa}\}$. Supóngase, por reducción al absurdo, que $C\neq\varnothing$.
  \item Por el principio del buen orden, $C$ tiene un elemento mínimo $m$.
  \item $m\neq 1$, porque $P(1)$ es $1=1^{2}$, cierta. Luego $m\ge 2$ y $m-1\in\N$; por la minimalidad de $m$, $m-1\notin C$, es decir, $P(m-1)$ es cierta: $\sum_{k=1}^{m-1}(2k-1)=(m-1)^{2}$.
  \item Sumando $2m-1$: $\sum_{k=1}^{m}(2k-1)=(m-1)^{2}+2m-1=m^{2}$, es decir, $P(m)$ es cierta, en contra de $m\in C$.
  \item Por tanto $C=\varnothing$ y $P(n)$ vale para todo $n\in\N$. $\blacksquare$
\end{enumerate}
\resultado \textbf{$\sum_{k=1}^{n}(2k-1)=n^{2}$}, demostrado por el buen orden.
\comprobacion los pasos 3 y 4 son exactamente el caso base y el paso de inducción del \cref{eje:5-plantilla}, reorganizados: es la idea de la \cref{prop:bo-induccion}.
\end{solucion}

\begin{solucion}{eje:5-caballos}
\datos la «demostración» propuesta.
\planteamiento comprobar el paso de inducción para cada $n$, en particular el menor.
\pasos
\begin{enumerate}[label=\arabic*.]
  \item El paso afirma que los dos grupos de $n$ caballos (sin el primero y sin el último) comparten algún caballo, y que ese caballo común transmite el color de un grupo al otro.
  \item Para $n=1$, el conjunto tiene $n+1=2$ caballos $\{c_1,c_2\}$. Los grupos son $\{c_2\}$ y $\{c_1\}$, que \textbf{no} comparten ningún caballo. No hay forma de relacionar el color de $c_1$ con el de $c_2$.
  \item Para $n\ge 2$ el argumento sí funciona (los grupos comparten $c_2,\dots,c_n$), pero la cadena se rompe en $P(1)\Rightarrow P(2)$, y sin ese eslabón no se deduce nada para $n\ge 2$.
\end{enumerate}
\resultado \textbf{el paso falla para $n=1$} (paso de 1 a 2 caballos), porque los dos grupos son disjuntos.
\comprobacion el enunciado $P(2)$ es falso (dos caballos de colores distintos), coherente con que el paso $P(1)\Rightarrow P(2)$ sea el defectuoso, ya que $P(1)$ es verdadera.
\end{solucion}

\begin{solucion}{eje:5-factorial}
\datos $S_n=\sum_{k=1}^{n}k\cdot k!$.
\planteamiento calcular, conjeturar y demostrar por inducción.
\pasos
\begin{enumerate}[label=\arabic*.]
  \item $S_1=1$; $S_2=1+4=5$; $S_3=5+18=23$; $S_4=23+96=119$; $S_5=119+600=719$.
  \item Los factoriales $(n+1)!$ son $2,6,24,120,720$: en cada caso $S_n=(n+1)!-1$. Conjetura: $S_n=(n+1)!-1$.
  \item \emph{Caso base:} $S_1=1=2!-1$. \emph{Paso:} si $S_n=(n+1)!-1$, entonces $S_{n+1}=(n+1)!-1+(n+1)(n+1)!=(n+1)!\,(1+n+1)-1=(n+2)!-1$. Por inducción, la conjetura es cierta para todo $n\in\N$. $\blacksquare$
\end{enumerate}
\resultado $S_n=(n+1)!-1$; valores $1,5,23,119,719$.
\comprobacion la identidad $k\cdot k!=(k+1)!-k!$ hace la suma telescópica: $S_n=(n+1)!-1!$. Valores comprobados con Python.
\end{solucion}

\begin{solucion}{eje:5-sumas}
\datos $\sum_{k=1}^{n}(3k-1)$; $\sum_{k=0}^{7}2^{k}$.
\planteamiento linealidad del sumatorio y \cref{ej:gauss}; \cref{ej:geometrica} con $r=2$.
\pasos
\begin{enumerate}[label=\arabic*.]
  \item $\sum_{k=1}^{n}(3k-1)=3\sum_{k=1}^{n}k-\sum_{k=1}^{n}1=\frac{3n(n+1)}{2}-n=\frac{3n^{2}+n}{2}=\frac{n(3n+1)}{2}$.
  \item Inducción. \emph{Caso base:} $n=1$: $2=\frac{1\cdot 4}{2}$. \emph{Paso:} si vale para $n$, $\frac{n(3n+1)}{2}+3(n+1)-1=\frac{3n^{2}+n+6n+4}{2}=\frac{3n^{2}+7n+4}{2}=\frac{(n+1)(3n+4)}{2}$, que es la fórmula para $n+1$. $\blacksquare$
  \item $\sum_{k=0}^{7}2^{k}=\frac{2^{8}-1}{2-1}=255$.
  \item Un byte tiene 8 bits; el mayor número representable es $11111111_2=2^{7}+\dots+2^{0}=255$, y hay $256$ valores, de $0$ a $255$.
\end{enumerate}
\resultado (a) $\dfrac{n(3n+1)}{2}$; (b) $255$.
\comprobacion (a) con $n=4$: $2+5+8+11=26=\frac{4\cdot 13}{2}$. (b) $1+2+4+\dots+128=255$; en hexadecimal, \texttt{0xFF} (módulo~12).
\end{solucion}

\begin{solucion}{eje:5-log}
\datos $p_1,\dots,p_n>0$; $\ln(ab)=\ln a+\ln b$ para $a,b>0$.
\planteamiento inducción sobre $n$, usando la definición recursiva del productorio, $\prod_{t=1}^{n+1}p_t=\bigl(\prod_{t=1}^{n}p_t\bigr)p_{n+1}$.
\pasos
\begin{enumerate}[label=\arabic*.]
  \item Sea $P(n)$: «para cualesquiera reales positivos $p_1,\dots,p_n$, $\ln\prod_{t=1}^{n}p_t=\sum_{t=1}^{n}\ln p_t$». \emph{Caso base:} $n=1$, ambos miembros valen $\ln p_1$.
  \item \emph{Paso.} Sea $n\ge 1$, supóngase $P(n)$ y sean $p_1,\dots,p_{n+1}>0$. El número $a=\prod_{t=1}^{n}p_t$ es positivo (producto de positivos), y $b=p_{n+1}>0$.
  \item Por la propiedad del logaritmo y la hipótesis de inducción, $\ln\prod_{t=1}^{n+1}p_t=\ln(ab)=\ln a+\ln b=\sum_{t=1}^{n}\ln p_t+\ln p_{n+1}=\sum_{t=1}^{n+1}\ln p_t$.
  \item Por inducción, $P(n)$ vale para todo $n\in\N$. $\blacksquare$
  \item En el módulo~1, $P(x_{1:T})=\prod_{t}P(x_t\mid x_{1:t-1})$ es un producto de $T$ factores en $(0,1]$, que para $T$ grande es un número diminuto (por ejemplo, $0{,}1^{1000}=10^{-1000}$, que no cabe en coma flotante). Su logaritmo es una suma de $T$ términos moderados ($1000\cdot\ln 0{,}1\approx -2302{,}6$), con la que se trabaja sin problemas numéricos; además, maximizar una suma es más sencillo que maximizar un producto.
\end{enumerate}
\resultado \textbf{el logaritmo de un producto finito es la suma de los logaritmos.}
\comprobacion $\ln(2\cdot 3\cdot 5)=\ln 30\approx 3{,}4012$ y $\ln 2+\ln 3+\ln 5\approx 0{,}6931+1{,}0986+1{,}6094=3{,}4012$.
\end{solucion}

% ---------------------------------------------------------------------
\section{Tema 6}

\begin{solucion}{eje:6-clasificar}
\datos cuatro enunciados.
\planteamiento identificar la forma lógica (\cref{tema:2}) y aplicar las reglas de demostración y de refutación del \cref{tema:6}.
\pasos
\begin{enumerate}[label=(\alph*)]
  \item Universal, $\forall p\,(p\text{ primo}\wedge p>2\Rightarrow p\text{ impar})$. Se demuestra en general; se refutaría con un primo par mayor que~2.
  \item Existencial. Se demuestra con un testigo ($101$); se refutaría probando que ningún entero entre $100$ y $110$ es primo.
  \item Existencia y unicidad. Se demuestra con un testigo ($x=2$) y un argumento de unicidad; se refutaría mostrando que no hay solución o que hay dos.
  \item Universal negativo, $\forall n\in\Z\ (n^{2}\neq 3)$. Se demuestra en general; se refutaría con un entero $n$ con $n^{2}=3$.
\end{enumerate}
\resultado (a) y (d) universales; (b) existencial; (c) de existencia y unicidad.
\comprobacion los cuatro son verdaderos: en (d), $1^{2}<3<2^{2}$ y $\abs{n}\ge 2\Rightarrow n^{2}\ge 4$.
\end{solucion}

\begin{solucion}{eje:6-contraejemplos}
\datos cuatro enunciados universales falsos.
\planteamiento dar un caso concreto que cumpla la hipótesis y no la tesis.
\pasos
\begin{enumerate}[label=(\alph*)]
  \item $x=\tfrac12$: $x^{2}=\tfrac14<\tfrac12$. (También $x=0$ o $x=1$, donde $x^{2}=x$.)
  \item $a=b=1$: $\sqrt2\neq 2$.
  \item $n=2$: $n^{2}=4$ es múltiplo de~4 y $2$ no lo es.
  \item $n=3$: $2^{3}+1=9=3\cdot 3$.
\end{enumerate}
\resultado contraejemplos $x=\tfrac12$; $a=b=1$; $n=2$; $n=3$.
\comprobacion en (d), $2^{n}+1$ es primo para $n=1,2,4$ ($3,5,17$), lo que podría inducir a error, como en el \cref{ej:euler}.
\end{solucion}

\begin{solucion}{eje:6-demostrar-refutar}
\datos cuatro enunciados.
\planteamiento decidir primero con ejemplos y después demostrar o refutar.
\pasos
\begin{enumerate}[label=(\alph*)]
  \item \textbf{Falso}: $\sqrt2$ y $-\sqrt2$ son irracionales (el segundo por el \cref{eje:4-racional-irracional} con $a=-1$) y su suma es $0$.
  \item \textbf{Verdadero}: si $a=m/n$ y $b=p/q$ con $n,q\neq 0$, entonces $ab=mp/(nq)$ con $nq\neq 0$.
  \item \textbf{Falso}: $n^{2}+n+1=(n^{2}+n)+1$, y $n^{2}+n$ es par (\cref{prop:consecutivos}), luego $n^{2}+n+1$ es impar para todo $n$. Se refuta un existencial con una demostración universal.
  \item \textbf{Falso}: $x=-1$ no tiene raíz cuadrada real, pues $y^{2}\ge 0>-1$.
\end{enumerate}
\resultado (a) falso; (b) verdadero; (c) falso; (d) falso.
\comprobacion en (c) no basta probar $n=1,2,3$ ($3,7,13$, impares): hace falta el argumento general.
\end{solucion}

\begin{solucion}{eje:6-testigos}
\datos dos enunciados existenciales.
\planteamiento dar testigos concretos y verificarlos.
\pasos
\begin{enumerate}[label=\arabic*.]
  \item (a) Sean $a=\tfrac32$ y $b=\tfrac12$. Entonces $\floor{a}=1$, $\ceil{b}=1$ y $\floor{a}\ceil{b}=1$, mientras que $ab=\tfrac34$ y $\floor{ab}=0$. Como $1\neq 0$, el par $(a,b)$ es testigo.
  \item (b) Sean $a=9$, $b=6$, $c=4$, $d=2$. Se cumple $9>6>4>2>0$ y $81+36+16+4=137$. $\blacksquare$
\end{enumerate}
\resultado (a) $a=\tfrac32$, $b=\tfrac12$; (b) $(9,6,4,2)$.
\comprobacion una búsqueda exhaustiva con Python muestra que $(9,6,4,2)$ es la única cuaterna estrictamente decreciente de naturales positivos con suma de cuadrados $137$. En (a), con $a=b=\tfrac12$ no sirve: $0\cdot 1=0=\floor{1/4}$.
\end{solucion}

\begin{solucion}{eje:6-irracional-entre}
\datos $r,s\in\Q$, $r<s$.
\planteamiento construir el testigo $x=r+\dfrac{s-r}{\sqrt2}$.
\pasos
\begin{enumerate}[label=\arabic*.]
  \item Sean $r<s$ racionales y $x=r+\dfrac{s-r}{\sqrt2}$.
  \item Como $s-r>0$ y $\sqrt2>1$, se tiene $0<\dfrac{s-r}{\sqrt2}<s-r$, luego $r<x<s$.
  \item $x$ es irracional: $\dfrac{1}{\sqrt2}$ es irracional (si fuera racional, su inverso $\sqrt2$ también lo sería); por el \cref{eje:4-racional-irracional}, $(s-r)\cdot\dfrac{1}{\sqrt2}$ es irracional por ser $s-r$ racional no nulo, y $r+\dfrac{s-r}{\sqrt2}$ es irracional por ser suma de un racional y un irracional.
  \item Por tanto, $x$ es un irracional entre $r$ y $s$. $\blacksquare$
\end{enumerate}
\resultado \textbf{entre dos racionales distintos hay un irracional}, por ejemplo $r+(s-r)/\sqrt2$.
\comprobacion $r=0$, $s=1$: $x=1/\sqrt2\approx 0{,}7071\in(0,1)$. La prueba es constructiva.
\end{solucion}

\begin{solucion}{eje:6-paridad-unica}
\datos $n\in\Z$; \cref{teo:division} con $b=2$.
\planteamiento existencia y unicidad del cociente; después, reducción al absurdo directa para «par e impar».
\pasos
\begin{enumerate}[label=\arabic*.]
  \item \emph{Existencia.} Por el \cref{teo:division} con $a=n$ y $b=2$, existen $k,r\in\Z$ con $n=2k+r$ y $r\in\{0,1\}$; es decir, $n=2k$ o $n=2k+1$.
  \item \emph{Unicidad.} Si $k$ y $k'$ cumplen la condición, entonces $n=2k+r=2k'+r'$ con $r,r'\in\{0,1\}$, y la unicidad de la división euclídea (\cref{prop:division-unicidad}) da $k=k'$ (y $r=r'$).
  \item \emph{Ningún entero es par e impar} (prueba directa). Supóngase, por reducción al absurdo, que $n=2k=2\ell+1$ con $k,\ell\in\Z$. Entonces $2(k-\ell)=1$, luego $k-\ell=\tfrac12$. Pero $k-\ell$ es entero y no hay enteros estrictamente entre $0$ y $1$. Contradicción. $\blacksquare$
\end{enumerate}
\resultado \textbf{$k$ existe y es único; ningún entero es par e impar.}
\comprobacion $n=-7=2(-4)+1$: $k=-4$ (no $k=-3$, que daría $-6+1=-5$).
\end{solucion}

\begin{solucion}{eje:6-terna}
\datos $p$, $p+2$, $p+4$ primos.
\planteamiento existencia con $p=3$; unicidad con los restos módulo~3.
\pasos
\begin{enumerate}[label=\arabic*.]
  \item \emph{Existencia.} $p=3$: $3$, $5$ y $7$ son primos.
  \item \emph{Unicidad.} Sea $p$ un primo tal que $p,p+2,p+4$ son primos. Por la división euclídea, $p\equiv 0$, $1$ o $2\pmod 3$.
  \item Si $p\equiv 1$, entonces $p+2\equiv 0\pmod 3$, y como $p+2>3$, $p+2$ sería compuesto. Si $p\equiv 2$, entonces $p+4\equiv 6\equiv 0\pmod 3$ y $p+4>3$ sería compuesto. Ambos casos son imposibles.
  \item Luego $p\equiv 0\pmod 3$, es decir, $3\divides p$, y como $p$ es primo, $p=3$. $\blacksquare$
\end{enumerate}
\resultado \textbf{$p=3$ es el único} (terna $3,5,7$).
\comprobacion búsqueda con Python hasta $1000$: solo $p=3$. Entre tres números $p,p+2,p+4$ uno siempre es múltiplo de~3, porque sus restos módulo~3 son $p,p+2,p+1$, distintos.
\end{solucion}

\begin{solucion}{eje:6-mersenne}
\datos (a) $n$ primo $\Rightarrow 2^{n}-1$ primo; (b) $2^{n}-1$ primo $\Rightarrow n$ primo.
\planteamiento (a) buscar contraejemplo; (b) contrarrecíproco con la factorización indicada.
\pasos
\begin{enumerate}[label=\arabic*.]
  \item (a) \textbf{Falso}: $n=11$ es primo y $2^{11}-1=2047=23\cdot 89$.
  \item (b) Se demuestra el contrarrecíproco: si $n\ge 2$ no es primo, $2^{n}-1$ no es primo. Sea $n=ab$ con $1<a,b<n$.
  \item Con $x=2^{a}$ y $m=b$: $2^{n}-1=(2^{a})^{b}-1=(2^{a}-1)\bigl(2^{a(b-1)}+\dots+2^{a}+1\bigr)$.
  \item Como $a\ge 2$, $2^{a}-1\ge 3>1$; como $a<n$, $2^{a}-1<2^{n}-1$. Luego $2^{a}-1$ es un divisor de $2^{n}-1$ estrictamente entre $1$ y $2^{n}-1$, y $2^{n}-1$ es compuesto. $\blacksquare$
\end{enumerate}
\resultado (a) \textbf{falso} ($2^{11}-1=23\cdot 89$); (b) \textbf{verdadero}.
\comprobacion $n=6$: $2^{6}-1=63=(2^{2}-1)(2^{4}+2^{2}+1)=3\cdot 21$. La identidad del paso~3 es la serie geométrica del \cref{ej:geometrica} multiplicada por $x-1$.
\end{solucion}

\begin{solucion}{eje:6-promedio}
\datos $x_1+\dots+x_{10}=100$.
\planteamiento reducción al absurdo.
\pasos
\begin{enumerate}[label=\arabic*.]
  \item Supóngase, por reducción al absurdo, que $x_i<10$ para todo $i$.
  \item Sumando las diez desigualdades, $x_1+\dots+x_{10}<100$, en contra de la hipótesis.
  \item Luego existe $i$ con $x_i\ge 10$. $\blacksquare$
  \item La prueba \textbf{no es constructiva}: no indica qué $x_i$ cumple la propiedad.
  \item No existe necesariamente $i$ con $x_i>10$: con $x_1=\dots=x_{10}=10$ ninguno lo cumple.
\end{enumerate}
\resultado \textbf{existe $x_i\ge 10$} (prueba no constructiva); no necesariamente $x_i>10$.
\comprobacion es el principio del promedio: algún valor es mayor o igual que la media, $100/10=10$.
\end{solucion}

\begin{solucion}{eje:6-critica-1}
\datos la prueba propuesta.
\planteamiento revisar la redacción con las pautas del \cref{tema:6} y la lógica de cada paso.
\pasos
\begin{enumerate}[label=\arabic*.]
  \item \emph{Razonamiento.} Usa la misma $k$ para $x$ e $y$, lo que supone $y=x+1$: solo trata pares consecutivos (\cref{obs:3-errores}(b)).
  \item \emph{Redacción.} No presenta $x$ e $y$ como enteros arbitrarios, y concluye «impar» sin escribir $x+y$ en la forma $2m+1$ con $m$ entero.
  \item El teorema es \textbf{verdadero}. Prueba correcta: «Sean $x$ un entero par e $y$ un entero impar. Existen $k,\ell\in\Z$ con $x=2k$ e $y=2\ell+1$. Entonces $x+y=2(k+\ell)+1$, y como $k+\ell\in\Z$, $x+y$ es impar. $\blacksquare$»
\end{enumerate}
\resultado error de \textbf{variable reutilizada} y conclusión sin justificar; teorema cierto, prueba en el paso~3.
\comprobacion $x=4$, $y=7$: $k=2$, $\ell=3$, $x+y=11=2\cdot 5+1$. La prueba original no cubre este caso, pues $7\neq 4+1$.
\end{solucion}

\begin{solucion}{eje:6-critica-2}
\datos la prueba propuesta.
\planteamiento identificar qué condicional se ha probado.
\pasos
\begin{enumerate}[label=\arabic*.]
  \item El enunciado es $P\Rightarrow Q$ con $P$: $n\equiv 0\pmod 2$ y $Q$: $n\equiv 0\pmod 4$. Su contrarrecíproco es «si $n\not\equiv 0\pmod 4$, entonces $n\not\equiv 0\pmod 2$».
  \item La prueba supone $Q$ y deduce $P$: demuestra el \textbf{recíproco} $Q\Rightarrow P$, que es cierto pero no equivale al enunciado.
  \item El enunciado es \textbf{falso}: $n=2$ cumple $2\equiv 0\pmod 2$ y $2\not\equiv 0\pmod 4$.
\end{enumerate}
\resultado se ha demostrado el recíproco; el enunciado es falso (contraejemplo $n=2$).
\comprobacion es el error de la \cref{obs:3-errores}(d); un contraejemplo al enunciado basta para saber que ninguna prueba correcta puede existir.
\end{solucion}

\begin{solucion}{eje:6-uno-dos}
\datos la cadena de igualdades.
\planteamiento revisar qué operación no es válida.
\pasos
\begin{enumerate}[label=\arabic*.]
  \item Las igualdades $a^{2}=ab$, $a^{2}-b^{2}=ab-b^{2}$ y $(a+b)(a-b)=b(a-b)$ son correctas.
  \item El paso «simplificando» divide ambos miembros por $a-b$, que vale $0$ porque $a=b$. De $X\cdot 0=Y\cdot 0$ no se sigue $X=Y$.
  \item (Si además $b=0$, el último paso divide también por cero.)
\end{enumerate}
\resultado el error es \textbf{dividir por $a-b=0$}.
\comprobacion con $a=b=1$: $(1+1)\cdot 0=1\cdot 0$ es cierto ($0=0$), pero $2\neq 1$.
\end{solucion}

\begin{solucion}{eje:6-redactar}
\datos $a,b\in\N$, $a\divides b$, $b\divides a$.
\planteamiento redacción directa con frases completas y justificación de cada paso.
\pasos
\begin{enumerate}[label=\arabic*.]
  \item Sean $a,b\in\N$ con $a\divides b$ y $b\divides a$. Por la \cref{def:divide}, existen enteros $k$ y $\ell$ con $b=ak$ y $a=b\ell$.
  \item Sustituyendo la primera igualdad en la segunda, $a=ak\ell$. Como $a\ge 1$, en particular $a\neq 0$, y puede dividirse por $a$: $k\ell=1$.
  \item Como $a,b>0$ y $b=ak$, el entero $k$ es positivo, y por la misma razón $\ell>0$. Dos enteros positivos con producto~1 son ambos iguales a~1: si fuera $k\ge 2$, sería $k\ell\ge 2\ell\ge 2$. Así $k=\ell=1$.
  \item Por tanto, $b=a\cdot 1=a$. $\blacksquare$
\end{enumerate}
\resultado \textbf{$a=b$}, con la demostración redactada en los pasos 1–4.
\comprobacion la hipótesis de positividad es necesaria: en $\Z$, $a=2$ y $b=-2$ se dividen mutuamente y son distintos (ahí $k=\ell=-1$). Compárese con el borrador del enunciado: las flechas se han sustituido por frases que explican \emph{por qué} cada paso es válido.
\end{solucion}

\begin{solucion}{eje:6-unica-raiz}
\datos $a\neq 0$; $ax^{2}+bx+c=0$.
\planteamiento completar el cuadrado: $ax^{2}+bx+c=a\Bigl[\bigl(x+\tfrac{b}{2a}\bigr)^{2}-\tfrac{\Delta}{4a^{2}}\Bigr]$ con $\Delta=b^{2}-4ac$.
\pasos
\begin{enumerate}[label=\arabic*.]
  \item Se comprueba la identidad del planteamiento desarrollando: $a\bigl(x^{2}+\tfrac{b}{a}x+\tfrac{b^{2}}{4a^{2}}\bigr)-\tfrac{b^{2}-4ac}{4a}=ax^{2}+bx+c$. Como $a\neq 0$, $x$ es solución si y solo si $\bigl(x+\tfrac{b}{2a}\bigr)^{2}=\tfrac{\Delta}{4a^{2}}$.
  \item ($\Leftarrow$) Si $\Delta=0$, la condición es $\bigl(x+\tfrac{b}{2a}\bigr)^{2}=0$, que equivale a $x=-\tfrac{b}{2a}$. Hay solución y es única.
  \item ($\Rightarrow$) Por contrarrecíproco: si $\Delta\neq 0$, no hay solución única. Si $\Delta<0$, el miembro derecho es negativo y el izquierdo no: no hay solución. Si $\Delta>0$, sea $s=\sqrt{\Delta}/(2\abs{a})>0$; entonces $x=-\tfrac{b}{2a}+s$ y $x=-\tfrac{b}{2a}-s$ son dos soluciones distintas.
  \item En ambos casos falla la existencia o la unicidad. Por tanto, hay solución única si y solo si $\Delta=0$, y entonces es $x=-b/(2a)$. $\blacksquare$
\end{enumerate}
\resultado \textbf{solución única $\iff b^{2}=4ac$}, y es $x=-b/(2a)$.
\comprobacion $x^{2}-6x+9=(x-3)^{2}$: $\Delta=36-36=0$ y $x=6/2=3$. $x^{2}-1$: $\Delta=4>0$, soluciones $\pm1$. $x^{2}+1$: $\Delta=-4$, ninguna.
\end{solucion}


\backmatter
% =====================================================================
\begin{thebibliography}{99}
\addcontentsline{toc}{chapter}{Bibliografía}
\markboth{Bibliografía}{}
\small

\bibitem{hammack}
R.~Hammack.
\newblock \emph{Book of Proof}, 3.ª ed. (edición~3.4).
\newblock Richard Hammack, 2018. Libro de texto de acceso libre.
\newblock \url{https://richardhammack.github.io/BookOfProof/}

\bibitem{mcs}
E.~Lehman, F.~T. Leighton y A.~R. Meyer.
\newblock \emph{Mathematics for Computer Science}, revisión del 18 de mayo de 2015.
\newblock Texto del curso 6.042J del MIT.
\newblock \url{https://ocw.mit.edu/courses/6-042j-mathematics-for-computer-science-spring-2015/}

\bibitem{mit6042}
A.~R. Meyer y A.~Chlipala.
\newblock 6.042J: Mathematics for Computer Science.
\newblock Massachusetts Institute of Technology, MIT OpenCourseWare, primavera de 2015. Unidad~1, «Proofs».
\newblock \url{https://ocw.mit.edu/courses/6-042j-mathematics-for-computer-science-spring-2015/}

\bibitem{polya}
G.~Polya.
\newblock \emph{How to Solve It: A New Aspect of Mathematical Method}.
\newblock Princeton University Press, 1945.

\bibitem{cs103}
Stanford University.
\newblock CS103: Mathematical Foundations of Computing.
\newblock Edición archivada \texttt{cs103.1272}, consultada en octubre de 2026: clases «Mathematical Proofs», «Indirect Proofs», «Propositional Logic» y «First-Order Logic»; guías «Guide to Proofs», «Guide to Negation» y «Proofwriting Checklist»; \emph{Problem Set}~1.
\newblock \url{https://web.stanford.edu/class/cs103/}

\bibitem{velleman}
D.~J. Velleman.
\newblock \emph{How to Prove It: A Structured Approach}, 3.ª ed.
\newblock Cambridge University Press, 2019.

\end{thebibliography}


\end{document}
